\documentclass[11pt,letterpaper]{article}
\usepackage[T1]{fontenc}
\usepackage[utf8]{inputenc}
\usepackage{lmodern}
\usepackage[margin=1in,headheight=15pt]{geometry}
\usepackage{amsmath,amssymb,amsthm,mathtools,mathrsfs}
\usepackage{microtype,booktabs,longtable,tabularx,array,enumitem}
\usepackage[dvipsnames]{xcolor}
\usepackage{fancyhdr}
\usepackage[colorlinks=true,linkcolor=MidnightBlue,citecolor=MidnightBlue,urlcolor=MidnightBlue]{hyperref}
\usepackage[nameinlink,noabbrev]{cleveref}
\hypersetup{%
 pdftitle={Surcomplex Polynomial Algebra: Root Geometry, Newton Profiles, Multiscale Stability, and Residue Duality},
 pdfauthor={Merged research exposition},
 pdfsubject={Surcomplex polynomials; arbitrary-rank Hahn workspaces; root geometry; perturbation; residue and Bezoutian duality}}
\pagestyle{fancy}\fancyhf{}
\fancyhead[L]{\small Surcomplex polynomial algebra}
\fancyhead[R]{\small Order, scales, stability, duality}
\fancyfoot[C]{\thepage}
\setlength{\parskip}{0.25em}
\setlength{\parindent}{1.25em}
\setlist{leftmargin=2em,itemsep=0.15em,topsep=0.3em}
\setcounter{tocdepth}{2}
\numberwithin{equation}{section}
\newtheorem{theorem}{Theorem}[section]
\newtheorem{proposition}[theorem]{Proposition}
\newtheorem{lemma}[theorem]{Lemma}
\newtheorem{corollary}[theorem]{Corollary}
\theoremstyle{definition}
\newtheorem{definition}[theorem]{Definition}
\newtheorem{example}[theorem]{Example}
\theoremstyle{remark}\newtheorem{remark}[theorem]{Remark}
\newcommand{\No}{\mathbf{No}}
\newcommand{\SC}{\mathbf{SC}}
\newcommand{\C}{\mathbb C}
\newcommand{\R}{\mathbb R}
\newcommand{\N}{\mathbb N}
\newcommand{\Z}{\mathbb Z}
\newcommand{\Q}{\mathbb Q}
\newcommand{\OO}{\mathcal O}
\newcommand{\mm}{\mathfrak m}
\newcommand{\HH}{\mathscr H_\Gamma}
\newcommand{\HHp}{\mathscr H_{\Gamma,>0}}
\newcommand{\HHn}{\mathscr H_{\Gamma,\ge0}}
\newcommand{\st}{\operatorname{st}}
\newcommand{\supp}{\operatorname{supp}}
\newcommand{\ord}{\operatorname{ord}}
\newcommand{\Res}{\operatorname{Res}}
\newcommand{\Disc}{\operatorname{Disc}}
\newcommand{\Tr}{\operatorname{Tr}}
\newcommand{\Nm}{\operatorname{Nm}}
\newcommand{\Hom}{\operatorname{Hom}}
\newcommand{\lc}{\operatorname{lc}}
\newcommand{\conv}{\operatorname{conv}}
\newcommand{\id}{\operatorname{id}}
\newcommand{\sig}{\operatorname{sig}}
\newcommand{\sgn}{\operatorname{sgn}}
\newcommand{\rem}{\operatorname{rem}}
\newcommand{\NF}{\operatorname{NF}}
\newcommand{\adj}{\operatorname{adj}}
\newcommand{\Bge}[2]{B_{\ge}(#1,#2)}
\newcommand{\Bgt}[2]{B_{>}(#1,#2)}
\newcommand{\abs}[1]{\left|#1\right|}
\newcommand{\sh}{^{\sharp}}
\newcolumntype{L}{>{\raggedright\arraybackslash}X}
\newcolumntype{P}[1]{>{\raggedright\arraybackslash}p{#1}}
\begin{document}
\hypersetup{pageanchor=false}
\begin{titlepage}
\centering
\vspace*{0.7cm}
{\large\scshape A merged continuation of the supplied surcomplex manuscripts\par}
\vspace{1cm}
{\Huge\bfseries Surcomplex\par Polynomial Algebra\par}
\vspace{0.55cm}
{\LARGE Root Geometry, Newton Profiles,\par Multiscale Stability, and Residue Duality\par}
\vspace{0.85cm}
{\large September 21, 2026\par}
\vspace{0.8cm}
\begin{minipage}{0.94\textwidth}
\begin{abstract}
We develop finite-degree polynomial algebra over the surcomplex class field
$\SC=\No[i]$, keeping two geometries deliberately apart: the surreal-valued
modulus, which gives ordered disks, and the natural Hahn valuation, which
gives valuation balls of arbitrary ordered-group rank. All proofs take place
in set-sized algebraically closed Hahn workspaces $K_\Gamma=\C((t^\Gamma))$,
and every degree is an ordinary finite integer.

The ordered part proves root bounds at arbitrary surreal radii, Gauss--Lucas
with surreal convex weights and its weighted incomplete-polynomial converse,
Jensen's disk theorem, a sharp Schoenberg second-moment inequality with its
collinearity equality case, a precisely delimited real-closed-field transfer
giving a polynomial Rouch\'e theorem on order circles, and Hermite's signature
criterion for counting real roots. It also determines the \emph{exact} root
locus of a polynomial under bounded coefficient uncertainty, in both
geometries, and exhibits the structural contrast between them: a sum of
available magnitudes against a minimum of available valuations.

The scale-sensitive part is organized by one initial-polynomial identity. It
counts roots in every closed ball, open ball, shell and residue direction,
proves Newton's root-valuation rule over an arbitrary ordered exponent group,
gives exact polynomial image balls and fiber degrees, and --- after
differentiation --- shows that a ball containing $k\ge1$ roots contains
exactly $k-r$ roots of the $r$th derivative for $1\le r\le k$. A finite root-cluster tree, an
explicit branch polynomial, two discriminant layer formulas and a nearest
neighbour identity locate the intervening critical directions.

Three perturbation theorems with genuinely different hypotheses are kept
apart: an optimal $\varepsilon/n$ valuation-H\"older matching that needs no
root separation at all; a rootwise threshold $T(P)=\max_i(d_i+\delta_i)$ that
preserves the entire matrix of pairwise distance valuations together with
their leading coefficients, the derivative valuations and the discriminant
valuation, with an exact displacement valuation and a second-order Newton
error term; and the coarser discriminant certificate, shown to
demand strictly more precision on an explicit quartic and to be sharp at equality on an
explicit quadratic.

Support-controlled coprime factor lifting is proved with the certificate
$S^*$, and in a parameter form over $\mathscr H_\Gamma(U)=\OO(U)((t^\Gamma))$
with no shrinking of the ordinary domain. The algebraic package ends with a
universal residue pairing that is perfect over every commutative ring and
through every collision, an explicit dual basis with no discriminant
denominator, the trace identity $\Tr(m_h)=\lambda_P(P'h)$, finite polynomial
maps with branch-value and ramification formulas, a class-safe
Nullstellensatz for finite data, and --- in several variables --- finite free
normal forms, a coefficient-independent B\'ezoutian pairing, the Jacobian
trace formula through collisions, and an identification of the supplied
coherent contour residue in a delimited polynomial subclass.
\end{abstract}
\end{minipage}
\vfill
\begin{minipage}{0.94\textwidth}\small
\textbf{Status.} This article is a merged, theorem-driven research
exposition. It does not claim that any statement below is a first occurrence
in the literature, and it does not claim to resolve a named open conjecture.
Several central statements are classical over real-closed, valued or
arbitrary commutative-ring settings, and are identified as such. The
literature search behind the attributions was targeted rather than
exhaustive. The article has not been independently refereed. The repository's
formalization ledger records partial Lean coverage; no complete formalization
of this article is claimed. The symbolic checks recorded in the source archives validate
displayed finite examples and are not machine verification of the general
proofs. Per-theorem non-claims are kept where they were made and are
collected in \cref{polynomial:sec:scope}.
\end{minipage}
\end{titlepage}
\hypersetup{pageanchor=true}
\pagenumbering{roman}
{\small\tableofcontents}
\clearpage
\pagenumbering{arabic}

\section{Subject, sources, and conventions}\label{polynomial:sec:intro}

\subsection{What changes when the coefficients are surcomplex}
A surcomplex number is $x+iy$ with $x,y\in\No$ and $i^2=-1$. That this class
field is algebraically closed is established surreal theory
\cite{polynomial:Conway,polynomial:Gonshor,polynomial:Poonen}, not a new
fundamental theorem of algebra. For a polynomial of degree $n$, algebraic
closedness already supplies $n$ roots with multiplicity. It answers none of
the geometric questions. How do roots sit relative to a disk of infinite
radius? How many roots occupy a prescribed infinitesimal subcluster? Which
critical points sit \emph{between} two such subclusters? How much coefficient
precision guarantees a one-to-one matching of roots, and when can an
arbitrarily small perturbation make distinct roots collide? What survives a
collision on the residue-theoretic side?

Two different structures answer these questions, and they must not be
identified. The modulus $|x+iy|=\sqrt{x^2+y^2}$ is surreal-valued and
satisfies the ordinary triangle inequality. The natural valuation records a
leading Hahn exponent and is ultrametric. An order disk $|z-a|\le r$ is not
the valuation ball $v(z-a)\ge v(r)$: for instance $\Bge00$ contains every
finite element, not merely the modulus unit disk, while $2t^\rho\in\Bge0\rho$
but $2t^\rho\notin\{|z|\le t^\rho\}$. The whole article is organized by that
separation.

\subsection{The three merged manuscripts}
This article merges three supplied manuscripts that share a title and a
spine. We cite them as
\cite{polynomial:MsFactorization,polynomial:MsNewton,polynomial:MsTrees} and
identify their individual contributions wherever they differ.

\cite{polynomial:MsNewton} (``Order Geometry, Newton Profiles, and
Scale-Resolved Stability'') is the largest and supplies the backbone of the
scale-sensitive theory: the initial-polynomial identity
(\cref{polynomial:thm:initialroots}), Newton's root-valuation rule over an
arbitrary ordered exponent group (\cref{polynomial:thm:newton}), the exact
image balls and fiber degrees (\cref{polynomial:thm:imageballs}), the
$r$th-derivative critical count (\cref{polynomial:thm:criticalballs}), the
optimal $\varepsilon/n$ valuation-H\"older matching
(\cref{polynomial:thm:holder}), the second-order Newton error estimate
(\cref{polynomial:thm:stability}\eqref{polynomial:eq:newtonerror}), the
Jensen disk theorem (\cref{polynomial:thm:jensen}) and the sharp Schoenberg
second-moment inequality with its equality case
(\cref{polynomial:thm:schoenberg}).

\cite{polynomial:MsFactorization} (``Factorization, Root Geometry,
Multiscale Stability, and Residue Duality'') supplies the exact
coefficient-uncertainty root loci in both geometries
(\cref{polynomial:thm:pseudozeros}), the named rootwise threshold $T(P)$ with
its full cluster-preservation conclusion and its two-sided calibration
(\cref{polynomial:thm:stability,polynomial:cor:discprecision,polynomial:ex:sharpdisc,polynomial:ex:quartic}),
the telescoping cluster formula for the discriminant valuation
(\cref{polynomial:prop:clusterdisc}), the universal residue pairing with an
explicit dual basis and no discriminant denominator
(\cref{polynomial:thm:residuepairing}), the universal trace identity
(\cref{polynomial:thm:trace}), the finite flat map, branch-value and
ramification package
(\cref{polynomial:thm:finitemap,polynomial:thm:branchvalues,polynomial:prop:ramification}),
the class-safe Nullstellensatz for finite data
(\cref{polynomial:thm:nullstellensatz}) and the non-Noetherian witness in
\cref{polynomial:rem:nonnoetherian}.

\cite{polynomial:MsTrees} (``Ordered Geometry, Valuative Root Trees, and
Residue Duality'') is the only one of the three to carry the
degree-controlled complete-intersection and B\'ezoutian package beyond one
variable --- all three leave one variable somewhere, the other two for the
Nullstellensatz and finite-algebra material. It
supplies Hermite's signature criterion over surreal real closed fields
(\cref{polynomial:thm:hermite}), the weighted incomplete-polynomial converse
to Gauss--Lucas (\cref{polynomial:prop:weighted}), the nearest-neighbour
critical-point identity (\cref{polynomial:cor:nearest}), the branch polynomial
at a node of the cluster tree and the tree form of the discriminant
(\cref{polynomial:thm:branch,polynomial:prop:disctree}), the escape example
(\cref{polynomial:ex:escape}), and the entire several-variable package:
finite free normal forms over an arbitrary commutative ring
(\cref{polynomial:thm:globalfree}), the coefficient-independent B\'ezoutian
pairing (\cref{polynomial:thm:multiperfect}), the B\'ezoutian kernel and
Jacobian trace identity (\cref{polynomial:thm:jacobian}), Euler--Jacobi
vanishing (\cref{polynomial:cor:eulerjacobi}), global polynomial families on
one ordinary parameter domain (\cref{polynomial:thm:polyparameters}) and the
identification of the supplied coherent contour residue in a delimited
polynomial subclass (\cref{polynomial:thm:analyticcomparison}).

All three are explicitly continuations of earlier supplied manuscripts, cited
here as \cite{polynomial:SourceAnalysis} (\emph{Surcomplex Analysis:
Infinitesimal Calculus, Hahn-Coherent Holomorphy, and Contour Theory}),
\cite{polynomial:SourceResearch} (\emph{Hartogs Coherence, Finite Maps, and
Residue Duality over the Surcomplex Numbers}) and
\cite{polynomial:SourceGlobal} (\emph{Global Divisors and Cohomological
Obstructions in Hahn-Coherent Surcomplex Analysis}). We take from them the
notation $t^\gamma=\omega^{-\gamma}$, set-sized Hahn workspaces, strong
summability, the coherent category $\HH(U)$ and its strong evaluation, and
nothing else. As all three merged manuscripts say in one form or another: we
use that framework as motivation, not as a prerequisite for the polynomial
results, and we do not silently identify their coefficient sheaves with
ordinary convergent power-series rings or assume their proposed novelty
assessments as facts.

\subsection{Conventions: one symbol per concept}\label{polynomial:subsec:conventions}
The three sources use three notations for several of the same objects. We fix
one symbol per concept here and use only that symbol thereafter. Synonyms are
recorded once, in this paragraph, and never used again.

\begin{itemize}
\item \textbf{Workspace.} $\Gamma\subseteq\No$ is a set-sized divisible
ordered additive subgroup; $F=F_\Gamma=\R((t^\Gamma))$ and
$K=K_\Gamma=\C((t^\Gamma))=F_\Gamma[i]$, with $t^\gamma=\omega^{-\gamma}$.
(Synonym in the sources: $R_\Gamma$ for $F_\Gamma$.)
\item \textbf{Valuation and reduction.} $v$, $\lc$, $\OO_K=\{v\ge0\}$,
$\mm_K=\{v>0\}$, and the coefficient-at-zero map
$\st:\OO_K\to\C$. The map $\st$ is called the \emph{standard part};
\emph{reduction} and \emph{shadow} are synonyms used in the sources, and
$\operatorname{red}$ is the symbol one of them uses for it. We write $\st$ throughout and
never $\operatorname{red}$. The topology it induces is called the \emph{standard-part
topology}; the names \emph{residue topology}, \emph{reduction topology} and
\emph{shadow topology} occur in the sources for the same thing and are not
used again.
\item \textbf{Balls and disks.} $\Bge a\rho=\{v(z-a)\ge\rho\}$,
$\Bgt a\rho=\{v(z-a)>\rho\}$ are the closed and open \emph{valuation balls};
$D(a,r)$, $\overline D(a,r)$ are the open and closed \emph{order disks}.
(Synonyms: $B_{\mathrm{cl}},B_{\mathrm{op}}$; ``valuative disk''.) The words
``open'' and ``closed'' describe the defining inequalities only; no
fine-topological connectedness or compactness is used anywhere.
\item \textbf{Scale data.} The scale exponent is always $\rho\in\Gamma$
(synonyms in the sources: $\delta$, $\gamma$). The weighted Gauss valuation is
$w_{a,\rho}(P)$ (synonym: $g_P(a,\rho)$) and the \emph{initial polynomial} is
$I_{a,\rho}(P)\in\C[X]$ (synonym: the \emph{reduction polynomial}
$R_{a,\rho}$). Its variable is always $X$; the polynomial variable over $K$ is
always $z$.
\item \textbf{Roots and condition numbers.} Roots are $\alpha_i$ (synonym:
$r_i$); $d_{ij}=v(\alpha_i-\alpha_j)$; $d_i=v(P'(\alpha_i))=\sum_{j\ne i}d_{ij}$
(synonym: $\sigma_i$); $\delta_i=\max_{j\ne i}d_{ij}$; and the threshold is
\[
 T(P)=\max_i\,(d_i+\delta_i)
\]
(synonym: $\tau(P)$). A perturbation is always $E$ with coefficient precision
$\varepsilon=\min_j v([z^j]E)$ (synonyms: $H$ with $\mu$; $\Delta$).
\item \textbf{Residues.} $\lambda_P(h)=[z^{n-1}]\rem_P h$ in one variable and
$\lambda_F(H)=[x^{\mathrm{top}}]\NF_F(H)$ in several; $\Delta_P$, $\Delta_F$
are B\'ezoutians; $J_F$ is the Jacobian determinant.
\end{itemize}

\subsection{Ring discipline}\label{polynomial:subsec:rings}
Four coefficient rings with nearly identical notation circulate in the wider
set of manuscripts from which this report is drawn: a fixed-domain ring, a
common-domain germ ring, a radius-free ring whose Hahn coefficients are
individually convergent germs with no common polydisk, and a
formal-coefficient ring. Theorems about them are \emph{not}
interchangeable, and the notation does not distinguish them. We therefore
state the rule once and then obey it.

\begin{center}
\begin{tabularx}{\textwidth}{@{}P{0.30\textwidth}L@{}}
\toprule
Ring & Role in this article\\
\midrule
$K=\C((t^\Gamma))$, and $\OO_K$, $\mm_K$ & The coefficient field of every
univariate polynomial theorem below, and its valuation ring. Everything in
\crefrange{polynomial:sec:foundations}{polynomial:sec:multivariate} that
mentions roots, balls, valuations or precision is stated over $K$ or
$\OO_K$.\\[0.5em]
$\HH(U)=\OO(U)((t^\Gamma))$, its nonnegative-support subring $\HHn(U)$,
and the positive-support ideal $\HHp(U)\subseteq\HHn(U)$ &
The \emph{fixed ordinary domain} parameter ring: $U\subseteq\C^d$ is one
ordinary connected domain, fixed once and never shrunk. Exactly three
theorems are stated over it:
\cref{polynomial:thm:parameterhensel,polynomial:thm:polyparameters,polynomial:thm:analyticcomparison},
together with \cref{polynomial:subsec:coherentfamilies}.\\[0.5em]
An arbitrary commutative ring $R$ (one variable) or $A$ (several) with
identity & The coefficient ring of the universal duality theorems
\cref{polynomial:thm:residuepairing,polynomial:thm:trace,polynomial:thm:globalfree,polynomial:thm:multiperfect,polynomial:thm:jacobian}.
No Noetherian, reduced, domain or field hypothesis is used there.\\[0.5em]
$\C[z]((t^\Gamma))$ & Appears only as a \emph{counterexample ring} in
\cref{polynomial:rem:degree,polynomial:subsec:uniformdegree}: it permits
unbounded polynomial degrees among Hahn coefficients and is never a
coefficient ring for a theorem.\\
\bottomrule
\end{tabularx}
\end{center}

\noindent
Every theorem below names its ring in its statement. No theorem is
transferred from one of these rings to another, and in particular:

\begin{itemize}
\item nothing below is stated over a radius-free ring $\C\{z\}((t^\Gamma))$
whose Hahn coefficients are individually convergent germs with no common
polydisk, and nothing below is stated over a formal-coefficient ring
$\C[[z]]((t^\Gamma))$. Neither ring occurs in any of the three merged
manuscripts, and no result here should be read as a statement about them;
\item the parameter theorems are stated for one \emph{fixed} ordinary domain
$U$. They are not germ statements, and their proofs give no common radius and
claim none;
\item the universal-ring duality theorems are stated over an arbitrary
commutative ring. Specializing them to $\HH(U)$ or to $K$ is a ring
homomorphism and is always performed explicitly, never assumed.
\end{itemize}

\subsection{Degree, and what ``polynomial'' means}
Throughout, a polynomial has \emph{ordinary finite degree} in its variables.
Its coefficients may have arbitrary set-sized Hahn supports. We do not use
ordinal degrees, proper-class sums, or a polynomial-like series of unbounded
degree under the name ``polynomial.'' Every theorem about an individual
finite root configuration is proved in a set-sized field containing its data
and then interpreted in $\SC$. The degree hypothesis is a genuine
restriction, not a disguised assertion about all isolated complete
intersections; \cref{polynomial:subsec:uniformdegree} makes the dividing line
exact.

\section{Foundations: workspaces, supports, and two notions of size}
\label{polynomial:sec:foundations}

\subsection{A set-sized field for every finite problem}
We use a class set theory admitting the usual treatment of $\No$, for example
NBG with global choice, in the same way as the supplied manuscripts. Each
individual support, recursion, coefficient list, matrix and algebraic
certificate is a \emph{set}. Conway normal forms give, for a set-sized ordered
abelian subgroup $\Gamma\subseteq\No$,
\begin{equation}\label{polynomial:eq:workspace}
 K_\Gamma=\C((t^\Gamma))\subseteq\SC,\qquad
 F_\Gamma=\R((t^\Gamma)),\qquad t^\gamma=\omega^{-\gamma},
\end{equation}
an element being $\sum_{\gamma\in S}c_\gamma t^\gamma$ with $S\subseteq\Gamma$
well ordered in the increasing order and $c_\gamma\in\C$.

\begin{proposition}[A workspace for every set of data]\label{polynomial:prop:workspace}
Every set of surcomplex constants is contained in a field $K_\Gamma$ with
$\Gamma$ divisible and set-sized. Such a field is algebraically closed; its
fixed field under complex conjugation is the real closed ordered field
$F_\Gamma=\R((t^\Gamma))$, and $K_\Gamma=F_\Gamma[i]$. A polynomial already
split in $K_\Gamma$ acquires no additional roots, and no changed
multiplicities, after passage to a larger surcomplex workspace. More
generally, if $A=K_\Gamma[x_1,\ldots,x_s]/I$ is finite-dimensional, every
surcomplex point annihilating $I$ has all coordinates in $K_\Gamma$, and its
local multiplicity is unchanged by enlargement.
\end{proposition}
\begin{proof}
Take the union of the normal-form supports of the given constants, generate
its additive subgroup in $\No$, and take the divisible hull; these are set
operations. The Hahn-field algebraic-closure theorem applies because $\Gamma$
is divisible and the residue field $\C$ is algebraically closed
\cite[Corollary~4]{polynomial:Poonen}. Real and imaginary parts of a Hahn
series lie in $F_\Gamma$, so $K_\Gamma=F_\Gamma[i]$; an ordered field whose
extension by $i$ is algebraically closed is real closed, which gives the
second assertion. If $P=a_n\prod_{j=1}^n(z-\alpha_j)$ already in $K_\Gamma[z]$,
that identity persists in every extension field, and a product of linear
factors cannot vanish at an element different from all the $\alpha_j$. For the
last assertion, the coordinates of a point of a finite-dimensional quotient
are roots of the characteristic polynomials of the commuting multiplication
matrices, hence lie in $K_\Gamma$; after scalar extension each local factor
still has a nilpotent maximal ideal and the extension field as residue field,
so it is still local of the same vector-space dimension.
\end{proof}

We write $F,K$ for one chosen pair and enlarge workspaces to include centers,
scales and finitely many parameters without further comment. A theorem stated
for all surcomplex centers and scales means that each particular finite
selection is placed in such a workspace. It does not require one set-sized
field to contain every surreal number.

\begin{remark}[{What $\SC[z]$ is not}]\label{polynomial:rem:classsafety}
The notation $\SC[z]$ is convenient for the class of finite polynomials, but
it does not authorize an application of Zorn's lemma to a proper-class family
of ideals. Euclidean calculations concern finitely many polynomials. Maximal
ideals, finite-dimensional algebras and the Nullstellensatz of
\cref{polynomial:sec:multivariate} are first handled in the \emph{set} ring
$K_\Gamma[z]$ or $K_\Gamma[x_1,\ldots,x_d]$, and only then interpreted in
$\SC$.
\end{remark}

\subsection{Modulus, valuation, and the standard part}
For $z=x+iy\in K$ put $\bar z=x-iy$ and $|z|=\sqrt{z\bar z}=\sqrt{x^2+y^2}\in
F_{\ge0}$, the positive square root existing by real closedness. The
algebraic Cauchy--Schwarz identity
$(x^2+y^2)(u^2+v^2)-(xu+yv)^2=(xv-yu)^2$ gives $|zw|=|z||w|$ and
$|z+w|\le|z|+|w|$ after squaring two nonnegative sides. These arguments use
no limits.

For $A\ne0$ set $v(A)=\min\supp(A)$, $\lc(A)=A_{v(A)}$, and $v(0)=+\infty$.
Then
\begin{equation}\label{polynomial:eq:valuation}
 v(AB)=v(A)+v(B),\qquad v(A+B)\ge\min\{v(A),v(B)\},\qquad v(|A|)=v(A),
\end{equation}
the last from the leading positive coefficient of $A\bar A$. Every nonzero
ordinary complex number has valuation zero; in particular
\begin{equation}\label{polynomial:eq:integersunit}
 v(m)=0\qquad\text{for every nonzero integer }m .
\end{equation}
This residue-characteristic-zero fact is responsible for the exact critical
point counts of \cref{polynomial:sec:critical}. Write
$\OO_K=\{v\ge0\}$ and $\mm_K=\{v>0\}$ for the finite elements and the
infinitesimals; the coefficient at exponent zero gives
$\st:\OO_K\to\C$ with $\OO_K/\mm_K\cong\C$, and every finite $A$ is uniquely
$\st(A)+\varepsilon$ with $\varepsilon\in\mm_K$.

\begin{definition}[Order disks and valuation balls]\label{polynomial:def:balls}
For $a\in K$, $r\in F_{>0}$ and $\rho\in\Gamma$,
\[
 D(a,r)=\{z:|z-a|<r\},\qquad \overline D(a,r)=\{z:|z-a|\le r\},
\]
\[
 \Bge a\rho=\{z:v(z-a)\ge\rho\}=a+t^\rho\OO_K,
 \qquad
 \Bgt a\rho=\{z:v(z-a)>\rho\}=a+t^\rho\mm_K.
\]
A \emph{residue-direction ball} at this center and scale is
$a+t^\rho(c+\mm_K)$ for $c\in\C$; it is the open valuation ball
$\Bgt{a+t^\rho c}{\rho}$.
\end{definition}

For example $2t^\rho\in\Bge0\rho$ but $2t^\rho\notin\overline D(0,t^\rho)$, and
$|1+1|=2$, so the modulus is not ultrametric. If $v(x)>v(y)$ with $y\ne0$,
then $x/y$ is infinitesimal, so $|x|<|y|/m$ for every positive integer $m$:
valuation geometry forgets ordinary nonzero complex factors while recording
every surreal scale. This is exactly why its perturbation criteria differ from
modulus criteria, a point made precise in \cref{polynomial:thm:pseudozeros}.

\subsection{The support lemma, and the two things it does not say}
A family of Hahn series is \emph{strongly summable} if the union of its
supports is well ordered and only finitely many members contribute at each
exponent; the sum is then defined by finite coefficientwise addition. For a
well-ordered $S\subseteq\Gamma_{>0}$ put
$S^*=\{e_1+\cdots+e_k:k\ge0,\ e_j\in S\}$, including the empty sum $0$.

\begin{lemma}[Hahn--Neumann support calculus]\label{polynomial:lem:support}
If $S,T$ are well-ordered subsets of an ordered abelian group, then $S+T$ is
well ordered and each element has finitely many representations $s+t$. If
$S$ is well ordered and positive, then $S^*$ is well ordered and each exponent
has finitely many representations by nondecreasing finite words in $S$.
Consequently, if $T$ is well ordered, a family of monomial terms indexed by
$b\in T$ and finite words $(e_1,\ldots,e_k)$ in $S$, each supported at
$b+e_1+\cdots+e_k$, is strongly summable provided each such pair indexes
only finitely many terms.
\end{lemma}
\begin{proof}
The first two assertions are the classical support facts
\cite{polynomial:Neumann}; the same formulation is used in all three merged
manuscripts. For the last, decompose an output exponent as $s+\eta$ with
$\eta\in S^*$: there are finitely many such pairs, finitely many nondecreasing
words for each $\eta$, and finitely many orderings of each finite word. Both
requirements for strong summability have then been checked.
\end{proof}

The lemma justifies $(1+B)^{-1}=\sum_{k\ge0}(-B)^k$ for $B$ of positive
support, as an \emph{exact} Hahn identity. Two negative companions are
load-bearing and are stated here once, since every construction below is
checked coefficient by coefficient rather than by an improving estimate.

\begin{remark}[Strong summability is not a valuation limit]
\label{polynomial:rem:notalimit}
\Cref{polynomial:lem:support} does \emph{not} say that the partial sums
converge in the fine surreal topology. With $\Gamma=\Q+\Q\omega$ the strong
sum of $t^k$ over $k\ge0$ is $(1-t)^{-1}$, yet every remainder after finitely
many terms exceeds the positive tolerance $t^\omega$: the tails never reach
valuation $\omega$. Nor does the lemma say that $k\rho$ eventually exceeds
every positive value-group element; that is false in higher rank. No argument
of the form ``the valuation improves by a fixed positive amount at each
iteration'' justifies any construction in this article, and we never
substitute either false cofinality claim for
\cref{polynomial:lem:support}.
\end{remark}

\begin{remark}[Finitely many dependencies is not finitely many exponents
below]\label{polynomial:rem:belowbound}
``Only finitely many members contribute at a given exponent'' does not mean
``only finitely many exponents lie below it.'' The latter already fails below
$\omega$ in $\Q+\Q\omega$, and it fails in
\cref{polynomial:ex:nondiscrete} within a support of order type $\omega$ all
of whose exponents are below $1$. The distinction is what makes the
finite-parameter-replacement steps in the recursions of
\cref{polynomial:sec:hensel} legitimate. Well-ordering also cannot be weakened
to positivity: $\{1/m:m\ge1\}$ is positive and not well ordered.
\end{remark}

\subsection{The coherent parameter ring on a fixed ordinary domain}
\label{polynomial:subsec:coherentring}
For an ordinary connected domain $U\subseteq\C^d$, the supplied category is
\[
 \HH(U)=\OO(U)((t^\Gamma)),
\]
Hahn families of \emph{ordinary holomorphic} functions on \emph{one common,
fixed} ordinary domain $U$. The nonnegative-support series form the subring
$\HHn(U)$; those of strictly positive support form its ideal $\HHp(U)$,
which does not contain $1$. A datum $F=\sum f_\gamma t^\gamma$ is evaluated
at $c+\varepsilon$ with $c\in U$ and $\varepsilon\in\mm_K^d$ by Taylor
expanding each ordinary coefficient,
\begin{equation}\label{polynomial:eq:coherentevaluation}
 F\sh(c+\varepsilon)=\sum_\gamma\sum_{\alpha\in\N^d}
 \frac{\partial^\alpha f_\gamma(c)}{\alpha!}\,t^\gamma\varepsilon^\alpha ,
\end{equation}
whose support lies in $\supp(F)+(\bigcup_j\supp\varepsilon_j)^*$ and has
finite multiplicity by \cref{polynomial:lem:support}. This convention is
inherited from
\cite{polynomial:SourceAnalysis,polynomial:SourceResearch}; it is not a new
interpretation of a fine limit. The subscripts $\ge0$ and $>0$ are used here
with their stated meaning throughout, to avoid the notational mismatch between
the supplied papers, where a superscript or subscript $+$ carries different
conventions.

\begin{remark}[Uniform degree, and a ring that lacks it]\label{polynomial:rem:degree}
A polynomial with coefficients in $\HH(U)$ has a uniform finite degree bound
in $z$; conversely, if every coefficient polynomial in a Hahn datum has degree
at most one fixed finite $n$, collecting the $n+1$ powers of $z$ gives such a
polynomial. Without that bound this fails. The ring $\C[z]((t^\Gamma))$
permits unbounded polynomial degrees among its Hahn coefficients, and
\begin{equation}\label{polynomial:eq:geometricwarning}
 \sum_{k\ge0}t^kz^k
\end{equation}
belongs to it when $1\in\Gamma$. It is a legitimate coherent datum on the
thickening of every ordinary domain, where it equals $(1-tz)^{-1}$; it is
\emph{not} a finite polynomial, and it is not a global pole-free function of
all surreal scales, since the rational function has a pole at the infinite
point $z=t^{-1}$. This elementary example separates the subject of this
article from unrestricted Hahn-coherent entire data. All global polynomial
evaluations here are finite sums.
\end{remark}

\begin{example}[A global root invisible in a finite monad]\label{polynomial:ex:escape}
Let $t=\omega^{-1}$ and $F(z)=z+tz^2=z(1+tz)$. As a global polynomial it has
the two roots $0$ and $-t^{-1}$, and $K[z]/(F)$ has dimension two. On the
thickening of any fixed ordinary disk about zero, $1+tz$ is a Hahn-coherent
unit with strong inverse $\sum_{k\ge0}(-tz)^k$, so the corresponding analytic
quotient has rank one and sees only the root zero. Neither statement
contradicts the other. In particular, positive Hahn support of a perturbation
does \emph{not} by itself preserve the global polynomial degree or the total
number of roots over all of $\SC$. The same phenomenon in the form
$z^2+tz^3=z^2(1+tz)$ shows a double root in the infinitesimal monad of zero
together with a further root at $-t^{-1}$: the local standard part is $z^2$
while the global degree is three.
\end{example}

\section{Finite polynomial algebra and multiplicity}\label{polynomial:sec:algebra}
Throughout this section $K$ is an algebraically closed workspace as in
\cref{polynomial:prop:workspace}, so all statements apply to surcomplex
polynomials.

\subsection{Factorization, division, and the squarefree part}
\begin{theorem}[Finite polynomial algebra over $K$]\label{polynomial:thm:fta}
Every nonconstant $P\in\SC[z]$ of degree $n$ admits a factorization
\begin{equation}\label{polynomial:eq:factorization}
 P(z)=a_n\prod_{j=1}^n(z-\alpha_j)
 =a_n\prod_{i=1}^s(z-\alpha_i)^{m_i},\qquad a_n\ne0,\ \textstyle\sum_i m_i=n,
\end{equation}
with the multiset of roots unique, all of them lying in one algebraically
closed set-sized Hahn workspace containing the coefficients. For
$A,B\in K[z]$ with $B\ne0$ there are unique $Q,R$ with $A=BQ+R$ and
$\deg R<\deg B$; every pair not identically zero has a monic greatest common
divisor $D=UA+VB$; and every ideal of $K[z]$ is principal. The multiplicity of
$\alpha$ in $P$ is the least $j$ with $P^{(j)}(\alpha)\ne0$, and
\begin{equation}\label{polynomial:eq:logderivative}
 \frac{P'(z)}{P(z)}=\sum_{i}\frac{m_i}{z-\alpha_i},
 \qquad
 \gcd(P,P')=\prod_i(z-\alpha_i)^{m_i-1}
\end{equation}
for monic $P$, so $P$ is squarefree exactly when $\gcd(P,P')=1$.
\end{theorem}
\begin{proof}
Use \cref{polynomial:prop:workspace}, choose one root and divide by its linear
factor; induction on the finite degree gives the factorization, and canceling
a vanishing factor gives multiset uniqueness. Leading-term division lowers a
nonnegative integer degree at each step, so it terminates and gives $Q,R$; the
Euclidean algorithm then produces $D$ and the B\'ezout coefficients. A
nonzero member of least degree generates any nonzero ideal; the zero ideal
is generated by zero. The multiplicity
characterization is the finite Taylor identity
$P(\alpha+Y)=\sum_{j\le n}(P^{(j)}(\alpha)/j!)Y^j$ together with
characteristic zero. The product rule gives the first identity in
\eqref{polynomial:eq:logderivative}; writing $P=(z-\alpha_i)^{m_i}U$ with
$U(\alpha_i)\ne0$ gives $\ord_{\alpha_i}P'=m_i-1$, because $m_i\ne0$ in
characteristic zero, which gives the second.
\end{proof}

This is an application of established algebraic closedness, not a new proof of
the surreal normal-form theorem. These are finite algebraic facts, unaffected
by the sizes of the coefficients.

A monic polynomial in $\OO_K[z]$ has all its roots in $\OO_K$: at a root of
negative valuation the leading term has strictly smaller valuation than every
other term and cannot be cancelled. Conversely a monic polynomial with finite
roots has finite coefficients, by Vieta. Reduction therefore gives the exact
multiset identity
\begin{equation}\label{polynomial:eq:reductionfactor}
 \st(P)(z)=\prod_{j=1}^n\bigl(z-\st(\alpha_j)\bigr)
 \qquad(P\in\OO_K[z]\text{ monic}).
\end{equation}
Multiple roots of the reduction can split at many different infinitesimal
scales; reduction records their total multiplicity, not their separation.

\subsection{Chinese remainders, Hermite jets, and idempotents}
Let $P$ be monic as in \eqref{polynomial:eq:factorization}, and put
$A_i=(z-\alpha_i)^{m_i}$, $P_i=P/A_i$. The inverse of $P_i$ modulo $A_i$ is
the finite Taylor polynomial at $\alpha_i$ of the rational function $1/P_i$,
\begin{equation}\label{polynomial:eq:inversejet}
 C_i(z)=\sum_{j=0}^{m_i-1}\frac1{j!}
 \left(\frac{d^j}{dz^j}\frac1{P_i(z)}\right)_{z=\alpha_i}(z-\alpha_i)^j ,
\end{equation}
whose denominators are nonzero because distinct $\alpha_i$ are separated as
field elements, even when their differences are infinitesimal.

\begin{theorem}[Finite Hermite interpolation and CRT]\label{polynomial:thm:crt}
The map
\begin{equation}\label{polynomial:eq:finiteCRT}
 K[z]/(P)\longrightarrow\prod_{i=1}^sK[u]/(u^{m_i}),
 \qquad [H]\longmapsto\Bigl[\textstyle\sum_{j<m_i}\frac{H^{(j)}(\alpha_i)}{j!}u^j\Bigr]_i
\end{equation}
is an isomorphism: arbitrary jets of orders less than $m_i$ at the $\alpha_i$
are realized by a unique polynomial interpolant of degree less than $n$. The
classes $E_i=[P_iC_i]\in K[z]/(P)$ are orthogonal idempotents with sum $1$ and
give the identities of the factors. When all roots are simple this is the
Lagrange formula
\begin{equation}\label{polynomial:eq:lagrange}
 H(z)=\sum_{i=1}^nH(\alpha_i)\frac{P(z)}{(z-\alpha_i)P'(\alpha_i)},
 \qquad\deg H<n .
\end{equation}
\end{theorem}
\begin{proof}
The finite Taylor identity gives $P_iC_i\equiv1\pmod{A_i}$ and
$P_iC_i\equiv0\pmod{A_j}$ for $j\ne i$, so the remainder modulo $P$ of
$\sum_iP_i C_i J_i(z-\alpha_i)$ realizes prescribed jets $J_i$. A polynomial
in the kernel is divisible by each pairwise coprime $A_i$, hence by their
product $P$; unique division by $P$ supplies the degree bound and uniqueness.
The idempotent identities hold in each factor and therefore in the quotient.
\end{proof}

This is a genuinely finite interpolation theorem: only finitely many
coefficient supports are combined, so it meets none of the infinite
symmetric-support obstruction of \cite{polynomial:SourceGlobal}. Its inverse
can nevertheless introduce inverse powers of small separations --- already
interpolation at $0$ and $t^\rho$ uses $z/t^\rho$ to realize the values $0,1$
--- and the denominators in \eqref{polynomial:eq:lagrange} may be infinitely
large after inversion. No separation lower bound in ordinary real units is
required.

\subsection{Vieta, Newton sums, resultants and discriminants}
Write $P(z)=z^n+c_1z^{n-1}+\cdots+c_n$ and $s_k=\sum_j\alpha_j^k$. Expanding
the finite product gives $c_k=(-1)^ke_k(\alpha_1,\ldots,\alpha_n)$, and the
Newton identities
\begin{align}
 s_k+c_1s_{k-1}+\cdots+c_{k-1}s_1+kc_k&=0 &&(1\le k\le n),
 \label{polynomial:eq:newton1}\\
 s_k+c_1s_{k-1}+\cdots+c_ns_{k-n}&=0 &&(k>n)
 \label{polynomial:eq:newton2}
\end{align}
follow by expanding \eqref{polynomial:eq:logderivative} at infinity as a
formal Laurent series in $z^{-1}$ and comparing finitely many coefficients
after multiplying by $P$. No infinite analytic evaluation is involved.

For $P=a_n\prod_i(z-\alpha_i)$ and $Q=b_m\prod_j(z-\beta_j)$ define
\begin{equation}\label{polynomial:eq:resultantproduct}
 \Res(P,Q)=a_n^m\prod_{i=1}^nQ(\alpha_i)=a_n^mb_m^n\prod_{i,j}(\alpha_i-\beta_j),
\end{equation}
symmetric up to $(-1)^{nm}$, a polynomial in the coefficients, and equal to
the Sylvester determinant. For monic $P$ we will always use the
multiplication-determinant definition $\Res(P,Q)=\det(m_Q:K[z]/(P)\to
K[z]/(P))$, justified in \cref{polynomial:thm:trace}. Set
\begin{equation}\label{polynomial:eq:discdef}
 \Disc(P)=a_n^{2n-2}\prod_{i<j}(\alpha_i-\alpha_j)^2,
 \qquad
 \Disc(P)=(-1)^{n(n-1)/2}\Res(P,P')\ \text{ for monic }P,
\end{equation}
which vanishes exactly at a repeated root. When the relevant products do not
vanish,
\begin{align}
 v(\Res(P,Q))&=m\,v(a_n)+n\,v(b_m)+\sum_{i,j}v(\alpha_i-\beta_j),
 \label{polynomial:eq:resval}\\
 v(\Disc(P))&=(2n-2)v(a_n)+2\sum_{i<j}v(\alpha_i-\alpha_j)
 =\sum_i d_i\quad(P\text{ monic, squarefree}).\label{polynomial:eq:discval}
\end{align}
These quantities record collisions across all scales, not merely ordinary
coincidence of standard parts. The resultant vanishes exactly when $P,Q$ have
a common root, equivalently when the finite B\'ezout map is singular; this is
a coefficient-level test and not a promise of an effective algorithm for
equality of arbitrary infinitely supported surreal coefficients.

\section{Ordered (modulus) root geometry}\label{polynomial:sec:ordergeometry}

\subsection{Root bounds at arbitrary surreal radii}
\begin{proposition}[Cauchy bounds and a radial bound]\label{polynomial:prop:rootbounds}
Let $P(z)=a_nz^n+\cdots+a_0$ with $a_n\ne0$ and $n\ge1$. Every root satisfies
\begin{equation}\label{polynomial:eq:cauchybound}
 |\alpha|<1+\max_{j<n}|a_j/a_n| .
\end{equation}
If $a_0\ne0$ it also satisfies $|\alpha|>\bigl(1+\max_{1\le j\le
n}|a_j/a_0|\bigr)^{-1}$. With $M=\max_{j<n}|a_j/a_n|^{1/(n-j)}$ one has
$|\alpha|\le2M$, and if $M=0$ all roots are zero.
\end{proposition}
\begin{proof}
Divide by $a_n$ and set $A=\max_{j<n}|a_j/a_n|$; the case $A=0$ is immediate.
For $r=|\alpha|\ge1+A>1$ the equation $P(\alpha)=0$ and the finite geometric
identity give $r^n\le A(r^n-1)/(r-1)\le r^n-1$, a contradiction. Apply this to
the reciprocal polynomial $z^nP(1/z)/a_0$ for the lower bound. Finally if
$r>2M>0$ then $\sum_{j<n}|a_j/a_n|r^{j-n}\le\sum_{k=1}^n2^{-k}<1$, so the
leading term cannot be cancelled; the positive roots defining $M$ exist in the
real closed field.
\end{proof}

The bounds hold literally for infinite or infinitesimal coefficients, and no
Archimedean estimate is used --- the geometric sum is finite, and the argument
is valid for an infinite surreal $r$ as well as a finite one. Applying them to
$P(a+rz)$ gives corresponding bounds at any center and scale.

\subsection{Gauss--Lucas with surreal convex weights}
For a finite $S\subseteq K$ the \emph{$F$-convex hull} is
\[
 \conv_F(S)=\Bigl\{\sum_{r\in S}\theta_r r:\theta_r\in F,\ \theta_r\ge0,\
 \sum_r\theta_r=1\Bigr\}.
\]
Allowing surreal, rather than only ordinary real, weights is essential:
restricting them to ordinary real numbers would be a different and generally
too small a hull. This is a finite-dimensional algebraic notion of convexity;
no topological compactness of a proper-class set is involved. Every order disk
and every order half-plane is convex, by the triangle inequality and linearity
of the real part.

\begin{theorem}[Surcomplex Gauss--Lucas]\label{polynomial:thm:gausslucas}
Every root of $P'$ lies in the $F$-convex hull of the roots of a nonconstant
$P\in\SC[z]$. Explicitly, if $w$ is a critical point with $P(w)\ne0$ then
\begin{equation}\label{polynomial:eq:barycentric}
 w=\frac{\displaystyle\sum_i\frac{m_i\alpha_i}{|w-\alpha_i|^2}}
        {\displaystyle\sum_i\frac{m_i}{|w-\alpha_i|^2}},
\end{equation}
all weights being strictly positive elements of $F$. The same convex hull
contains the zeros of every nonzero higher derivative.
\end{theorem}
\begin{proof}
At a nonroot critical point the logarithmic derivative
\eqref{polynomial:eq:logderivative} vanishes. Since
$1/(w-\alpha_i)=(\bar w-\bar\alpha_i)/|w-\alpha_i|^2$, conjugating gives
$\sum_i m_i(w-\alpha_i)/|w-\alpha_i|^2=0$; every weight is positive, their sum
is nonzero and may be divided out, giving
\eqref{polynomial:eq:barycentric}. A critical point that is itself a root is
already in the hull. Repeating for successive derivatives and using convexity
gives the last assertion.
\end{proof}

The displayed proof is finite and algebraic; it does not call on compactness
of a convex hull or a topological separation theorem. Its classical analogue
is Gauss--Lucas, and modern formal developments of the ordinary theorem and
related root geometry are available \cite{polynomial:Eberl}. Neither
conclusion describes how critical points are distributed among infinitesimal
clusters; that more precise question is answered in
\cref{polynomial:sec:critical}.

\begin{proposition}[Weighted incomplete polynomials, and a converse]
\label{polynomial:prop:weighted}
For $\alpha_1,\ldots,\alpha_n\in K$ and weights $\lambda_j\ge0$ in $F$ with
$\sum_j\lambda_j>0$, every root of
$Q_\lambda(z)=\sum_{j=1}^n\lambda_j\prod_{k\ne j}(z-\alpha_k)$ belongs to
$\conv_F\{\alpha_1,\ldots,\alpha_n\}$. Conversely, for $n\ge2$ the union of
these root sets as $\lambda$ varies is \emph{exactly} that convex hull.
\end{proposition}
\begin{proof}
At a root distinct from all $\alpha_j$, divide by $\prod_j(z-\alpha_j)$ and
repeat the proof of \cref{polynomial:thm:gausslucas} with weights
$\lambda_j/|z-\alpha_j|^2$. For the converse let $w=\sum_j\theta_j\alpha_j$
with $\theta_j\ge0$, $\sum_j\theta_j=1$. If $w$ is not a listed root, take
$\lambda_j=\theta_j|w-\alpha_j|^2$; then
$\sum_j\lambda_j/(w-\alpha_j)=\sum_j\theta_j(\bar w-\bar\alpha_j)=0$.
At a listed root $\alpha_i$, choose a nonzero weight at an index $j\ne i$, so
that the incomplete product contains the factor $z-\alpha_i$.
\end{proof}

This is an ordered-field form of a classical extension of Gauss--Lucas
\cite{polynomial:Zhang}; the proof exhibits precisely why it survives for
surreal coefficients. It will be used again in
\cref{polynomial:thm:branch}, where the branch polynomial at a node of the
root-cluster tree is a weighted incomplete polynomial of exactly this kind.

\subsection{Jensen disks for real-surreal coefficients}
\begin{theorem}[Surcomplex Jensen disk theorem]\label{polynomial:thm:jensen}
Suppose $P\in F[z]$ is nonconstant. Every nonreal critical point belongs to at least one
closed order disk having a pair of conjugate nonreal roots as endpoints of a
diameter; equivalently, for some root pair $a\pm ib$ with $b\ne0$ the critical
point $w$ satisfies $|w-a|\le|b|$.
\end{theorem}
\begin{proof}
Conjugation preserves the root multiset. A critical point that is a root lies
on the appropriate boundary, so assume $w=x+iy$ is not a root and $y\ne0$. For
a real root $r$, $y^{-1}\operatorname{Im}(w-r)^{-1}=-|w-r|^{-2}<0$. For a
conjugate pair $a\pm ib$,
\begin{equation}\label{polynomial:eq:jensencalc}
 \frac1y\operatorname{Im}\left(\frac1{w-a-ib}+\frac1{w-a+ib}\right)
 =\frac{-2\bigl((x-a)^2+y^2-b^2\bigr)}{|w-a-ib|^2|w-a+ib|^2}.
\end{equation}
If $w$ were outside every indicated closed disk, every contribution to
$y^{-1}\operatorname{Im}(P'/P)(w)$, with its positive multiplicity, would be
negative, and their finite sum could not vanish. This contradicts $P'(w)=0$.
\end{proof}

The nonconstant hypothesis is necessary: for $P=1$ the derivative is zero,
so every nonreal point is a zero of $P'$ while there is no root pair to define
a disk. In particular a nonconstant real-rooted polynomial has only real critical points. The
assertion concerns $F$-valued disks; it is not a statement about ordinary
pictures obtained by discarding infinitesimals. Neither of the other two
merged manuscripts proves a Jensen disk theorem; one of them proves
\cref{polynomial:thm:hermite} in its place.

\section{A sharp second-moment inequality for critical points}
\label{polynomial:sec:moment}
The ordinary Schoenberg inequality is established mathematics, proved with
matrix differentiators and related spectral methods
\cite{polynomial:Pereira,polynomial:KushelTyaglov}. We give an algebraic proof
over a general real closed field, including the equality case, so that no
Hilbert-space completeness assumption is hidden in a transfer. This section is
the contribution of \cite{polynomial:MsNewton} alone; neither of the other two
merged manuscripts contains any second-moment or matrix-compression material.

\subsection{A matrix whose eigenvalues are the critical points}
Let $P$ be monic of degree $n\ge2$ with roots $\alpha_1,\ldots,\alpha_n$
repeated by multiplicity, and set
$D=\operatorname{diag}(\alpha_1,\ldots,\alpha_n)$,
$e=n^{-1/2}(1,\ldots,1)^T$, $\Pi=I-ee^*$, where $*$ is conjugate transpose.
Algebraic Gram--Schmidt over $F$ gives a real $n\times(n-1)$ matrix $U$ with
$U^*U=I$, $UU^*=\Pi$. Put $B=U^*DU$.

\begin{lemma}[Differentiating compression]\label{polynomial:lem:compression}
The characteristic polynomial of $B$ is $P'(z)/n$.
\end{lemma}
\begin{proof}
The unitary $(e\ U)$ turns $zI-D$ into a block matrix whose lower-right block
is $zI_{n-1}-B$, so its upper-left cofactor is $\det(zI_{n-1}-B)$. In the
original coordinates that cofactor is $e^*\adj(zI-D)e$, which for diagonal $D$
equals $n^{-1}\sum_i\prod_{j\ne i}(z-\alpha_j)=P'(z)/n$. All expressions are
polynomial identities, so no invertibility at a root of $P$ is assumed.
\end{proof}

\begin{lemma}[Finite-dimensional Schur inequality]\label{polynomial:lem:schur}
If $A\in\operatorname{Mat}_m(K)$ has eigenvalues $\lambda_1,\ldots,\lambda_m$
with multiplicity, then $\sum_j|\lambda_j|^2\le\Tr(AA^*)$, with equality
exactly when $A$ is normal.
\end{lemma}
\begin{proof}
Algebraic closedness provides an eigenvector; normalize it by the positive
square root of its squared norm, extend to an orthonormal basis by
Gram--Schmidt and induct, giving a unitary upper triangularization. In that
basis $\Tr(AA^*)$ is the sum of squared moduli of all entries, while the left
side is the sum over the diagonal; the difference is the sum over the strict
upper triangle and vanishes exactly when those entries vanish. A diagonal
matrix is normal and normality is unitarily invariant; conversely an
upper-triangular normal matrix is diagonal, since equality of the $(1,1)$
entries of $TT^*$ and $T^*T$ kills the first row above the diagonal and
induction handles the rest.
\end{proof}

\subsection{The inequality and its equality case}
\begin{theorem}[Sharp surcomplex Schoenberg inequality]\label{polynomial:thm:schoenberg}
Let $w_1,\ldots,w_{n-1}$ be the roots of $P'$ with multiplicity. Then
\begin{equation}\label{polynomial:eq:schoenberg}
 \sum_{j=1}^{n-1}|w_j|^2\le\frac{n-2}{n}\sum_{i=1}^n|\alpha_i|^2
 +\frac1{n^2}\Bigl|\sum_{i=1}^n\alpha_i\Bigr|^2 ,
\end{equation}
equivalently, with the root centroid $c=n^{-1}\sum_i\alpha_i$,
\begin{equation}\label{polynomial:eq:variance}
 \boxed{\displaystyle
 \sum_{j=1}^{n-1}|w_j-c|^2\le\frac{n-2}{n}\sum_{i=1}^n|\alpha_i-c|^2 .}
\end{equation}
Equality holds if and only if all the roots lie on an affine $F$-line in $K$;
the one-point configuration counts as collinear.
\end{theorem}
\begin{proof}
Apply \cref{polynomial:lem:schur} to $B$ and use
\cref{polynomial:lem:compression}. Cyclic invariance of the finite trace gives
\[
 \Tr(BB^*)=\Tr(D\Pi D^*\Pi)
 =\Bigl(1-\frac2n\Bigr)\sum_i|\alpha_i|^2
 +\frac1{n^2}\Bigl|\sum_i\alpha_i\Bigr|^2 ,
\]
which is \eqref{polynomial:eq:schoenberg}; Vieta for $P'$ gives
$\sum_jw_j=(n-1)c$, and replacing every root by its difference from $c$ gives
\eqref{polynomial:eq:variance}.

For equality put $v=\Pi De$ and $u=\Pi D^*e=\bar v$, the last because $\Pi$ is
real. On $e^\perp$, expanding $\Pi=I-ee^*$ and using $DD^*=D^*D$ represents
$BB^*-B^*B$ by $uu^*-vv^*$. If $v=0$ all roots equal $c$. Otherwise equality
of these two rank-one matrices forces $u=\zeta v$ with $\zeta\bar\zeta=1$,
their common nonzero image being one-dimensional. The coordinates of $v$ are
$(\alpha_i-c)/\sqrt n$, so normality is equivalent to
$\overline{\alpha_i-c}=\zeta(\alpha_i-c)$ for all $i$. Choosing a nonzero
$d=\alpha_j-c$ and dividing shows every ratio $(\alpha_i-c)/d$ is fixed by
conjugation, hence lies in $F$, so all roots lie on $c+dF$. Conversely that
collinearity makes $\bar v$ a unit-modulus multiple of $v$, hence $B$ normal.
\end{proof}

The coefficient $(n-2)/n$ cannot be improved, since nonconstant collinear
configurations attain equality. Both sides may be infinite, finite or
infinitesimal surreal quantities; they are compared in the ordered field, not
after taking standard parts. The roots $1,i,-1,-i$ of $z^4-1$ have all
critical points at zero, so the inequality is strict; multiplying every root
by a nonzero surcomplex scale gives an equally valid strict example at that
scale. The matrix proof uses the positive definite form on the
finite-dimensional space $K^n$ and algebraic square roots; it does not assert
completeness of an infinite-dimensional surcomplex Hilbert space.

\section{Real-closed-field transfer: Rouch\'e on order circles, and Hermite
signatures}\label{polynomial:sec:transfer}

\subsection{Exactly what transfer permits}\label{polynomial:subsec:transferscope}
Tarski's theorem gives quantifier elimination for real closed ordered fields, so a
first-order sentence in the ordered-field language true in $\R$ is true in
every real closed field \cite{polynomial:TarskiSurvey,polynomial:Swan}. For
any \emph{fixed} degree bound a complex-coefficient polynomial is finitely
many pairs of real coordinates; polynomial identities, circle equations and
squared-modulus inequalities are then formulas in that language. Root counts
with multiplicity are also encodable: introduce a finite list of roots, impose
the coefficient identities of a complete factorization, and impose by a finite
disjunction that exactly $k$ entries of the list lie in the specified disk.
This counts repeated entries rather than only distinct points; degrees below a
chosen bound are handled by a finite disjunction on the last nonzero
coefficient.

This is a theorem \emph{scheme}, one sentence for each fixed finite degree
bound. It does not transfer unrestricted quantification over functions,
sequences, arbitrary subsets, the standard natural numbers as a distinguished
definable set, or all fine-continuous functions. It also does not make a
proper class into a first-order set-sized model: transfer is applied to
$F_\Gamma$, after all parameters have been placed in that field. Constructive
real-algebraic winding-number and Sturm-chain methods give another route to
root counting over real closed fields \cite{polynomial:Eisermann}; we do not
need that machinery for the implications stated here.

\begin{theorem}[Rouch\'e on an order-modulus circle]\label{polynomial:thm:orderrouche}
Let $P,Q\in\SC[z]$ be nonzero, $a\in\SC$, $r\in\No_{>0}$. Suppose
\begin{equation}\label{polynomial:eq:orderrouchehyp}
 |Q(z)-P(z)|<|P(z)|\qquad\text{for every }z\in\SC\text{ with }|z-a|=r .
\end{equation}
Then $P$ and $Q$ have the same number of roots in $D(a,r)$, counted with
multiplicity, and neither has a root on the circle.
\end{theorem}
\begin{proof}
Fix the two finite degrees. Hypothesis \eqref{polynomial:eq:orderrouchehyp} is
a first-order formula: quantify the two real coordinates of $z$, write the
circle as $(\Re z-\Re a)^2+(\Im z-\Im a)^2=r^2$, and square both nonnegative
sides. The conclusion is first order at those fixed degrees by the encoding
above. The resulting universally quantified implication is true over $\R$ by
ordinary polynomial Rouch\'e, hence in every real closed field. Choose
$F_\Gamma$ containing the real and imaginary parts of the coefficients, $a$
and $r$; the hypothesis holds there because it holds in $\SC$, and by
\cref{polynomial:prop:workspace} the root lists already contain every
surcomplex root, so the conclusion is the desired count in $\SC$. Finally the
strict inequality excludes a boundary root of $P$, and a boundary root of $Q$
would make the two sides equal.
\end{proof}

This is a proof by transfer from a classical theorem, and is identified as
such. It does not define an integral along a fine-continuous real-parameter
surreal circle. It is also not the same statement as the leading-coefficient
Rouch\'e theorem on ordinary thickened domains in
\cite{polynomial:SourceAnalysis}: its hypothesis concerns the \emph{full}
order circle at the chosen surreal scale.

\begin{corollary}[Monomial dominance]\label{polynomial:cor:pellet}
Write $P(a+Z)=\sum_{j=0}^np_jZ^j$. If for some $k$ and $r>0$
\begin{equation}\label{polynomial:eq:pellet}
 |p_k|r^k>\sum_{j\ne k}|p_j|r^j ,
\end{equation}
then $P$ has exactly $k$ roots in $D(a,r)$ counted with multiplicity, and none
on the boundary. The case $k=0$ (no roots) is included.
\end{corollary}
\begin{proof}
On $|z-a|=r$ the difference between $P(z)$ and $p_k(z-a)^k$ has modulus at
most the right side of \eqref{polynomial:eq:pellet}. Apply
\cref{polynomial:thm:orderrouche} with that monomial as reference; it has
exactly $k$ roots in the disk, counting multiplicity at $a$.
\end{proof}

No strict inequality among the valuations of the terms is required. The
criterion can distinguish two terms at the same valuation by their leading
ordinary moduli --- a thing the valuation criterion of
\cref{polynomial:cor:valrouche} cannot do, and the clearest illustration of
why the two are not interchangeable.

\begin{example}[Testing only ordinary circle points is insufficient]
\label{polynomial:ex:circlewarning}
Let $t=\omega^{-1}$ and put
\[
 b=\frac{1+it}{\sqrt{1+t^2}},\qquad \alpha=(1-t^3)b,\qquad \beta=(1+t^3)b,
 \qquad P(z)=z-\alpha,\quad Q(z)=z-\beta .
\]
Then $|b|=1$, $|\alpha|=1-t^3<1$ and $|\beta|=1+t^3>1$, so $P$ and $Q$ have
different root counts in the unit order disk. Nevertheless
$|Q(\zeta)-P(\zeta)|<|P(\zeta)|$ at every \emph{ordinary} complex unit point
$\zeta$: for $\zeta\ne1$ one has $v(P(\zeta))=0$, while at $\zeta=1$ it is
$1$, the leading displacement of $b$ being $it$; the error has valuation $3$
in both cases. At the nonordinary unit point $b$, however, $|P(b)|=t^3$ and
$|Q(b)-P(b)|=2t^3$, so the hypothesis genuinely fails there. The exact
quantifier in \eqref{polynomial:eq:orderrouchehyp} cannot be replaced by
sampling ordinary points.
\end{example}

\subsection{Hermite trace forms count real roots}
Let $R$ be a real closed field, $P\in R[z]$ monic of degree $n$,
$A=R[z]/(P)$, and for $q\in A$
\begin{equation}\label{polynomial:eq:hermiteform}
 T_q(f,g)=\Tr_A(m_{qfg}).
\end{equation}
A symmetric matrix over a real closed field is diagonalizable by congruence;
its \emph{signature} is the number of positive minus the number of negative
diagonal entries, zeros ignored, well defined by Sylvester's law of inertia.
In the monomial basis its entries are $\Tr\bigl(q(C_P)C_P^{i+j}\bigr)$, where
$C_P$ is the companion matrix.

\begin{theorem}[Hermite's signature criterion over surreal real closed fields]
\label{polynomial:thm:hermite}
With the above notation,
\begin{equation}\label{polynomial:eq:hermitesignature}
 \sig T_q=\sum_{\substack{r\in R\\ P(r)=0}}\sgn q(r),
\end{equation}
the sum over \emph{distinct} real roots. In particular $\sig T_1$ is their
number, the rank of $T_1$ is the number of distinct roots over $R[i]$, and
$T_1$ is positive semidefinite if and only if every root is real. For $a<b$
with $P(a)P(b)\ne0$ and $q(z)=(z-a)(b-z)$,
\begin{equation}\label{polynomial:eq:intervalcount}
 \#\{r\in(a,b):P(r)=0\}=\tfrac12\bigl(\sig T_1+\sig T_q\bigr),
\end{equation}
again counting distinct roots; apply it to the squarefree multiplicity factors
for a count with multiplicity.
\end{theorem}
\begin{proof}
Factor $P$ over $R$ into powers of distinct real linear factors and
irreducible quadratics; Chinese remaindering makes the summands orthogonal. On
a local factor at a real root $r$ of multiplicity $m$, multiplication by $h$
has trace $mh(r)$, since its nilpotent part is triangular with zero diagonal;
the form therefore factors through the one-dimensional residue field as
$mq(r)uv$, with signature $\sgn q(r)$. For a nonreal conjugate pair the
residue field is $R[i]$ and, with $q(r)=a_0+ib_0$, the residue form is $m$
times $\begin{pmatrix}2a_0&-2b_0\\-2b_0&-2a_0\end{pmatrix}$, of determinant
$-4(a_0^2+b_0^2)<0$ when $q(r)\ne0$: signature zero, rank two. If $q(r)=0$ it
is zero. Nilpotent directions lie in the kernel in both cases. Adding
signatures gives \eqref{polynomial:eq:hermitesignature}, and $q=1$ gives the
remaining assertions; \eqref{polynomial:eq:intervalcount} follows since $q>0$
exactly on $(a,b)$.
\end{proof}

This classical signature method is related to the residue and B\'ezoutian
forms discussed in \cite{polynomial:BCRS}; it is \emph{not} the nondegenerate
residue pairing of \cref{polynomial:sec:residue}. The two forms are compared
explicitly in \cref{polynomial:thm:discriminant}: the trace pairing degenerates
exactly at a repeated root, the residue pairing never does. Only one of the
three merged manuscripts proves this criterion; the one that proves
\cref{polynomial:thm:jensen} proves no signature criterion, and neither
subsumes the other.

\section{Exact root loci under coefficient uncertainty}
\label{polynomial:sec:pseudozeros}
The next theorem gives a \emph{sharp} description of the reachable root set
under bounded coefficient errors, not merely a sufficient bound, and it
exposes the structural contrast between the two geometries. Weighted
coefficient-uncertainty root regions are called \emph{pseudozero sets} in
ordinary polynomial analysis, and root neighbourhoods of this type have
established applications \cite{polynomial:FaroukiHan}; the argument below is
given in full for the surcomplex settings used here. This section is the
contribution of \cite{polynomial:MsFactorization} alone; neither of the other
two merged manuscripts has an uncertainty-locus section.

Fix a monic degree-$n$ polynomial $P\in\SC[z]$ and a set of perturbable
indices $J\subseteq\{0,\ldots,n-1\}$, and consider $Q(z)=P(z)+\sum_{j\in
J}e_jz^j$, so that all perturbations retain the degree and leading
coefficient.

\begin{theorem}[Exact uncertainty root sets in two geometries]
\label{polynomial:thm:pseudozeros}
Fix $z\in\SC$.

\emph{Order-modulus version.} Given $\eta_j\in\No_{\ge0}$, such a $Q$ with
$Q(z)=0$ and $|e_j|\le\eta_j$ for every $j\in J$ exists \emph{if and only if}
\begin{equation}\label{polynomial:eq:orderpseudozero}
 |P(z)|\le\sum_{j\in J}\eta_j|z^j| .
\end{equation}

\emph{Valuation version.} Given $\nu_j\in\Gamma$, such a $Q$ with $Q(z)=0$ and
$v(e_j)\ge\nu_j$ for every $j\in J$ exists \emph{if and only if}
\begin{equation}\label{polynomial:eq:valpseudozero}
 v(P(z))\ \ge\ \min_{j\in J}\bigl(\nu_j+v(z^j)\bigr).
\end{equation}
Here $z^0=1$, $v(0)=+\infty$, and the minimum of an empty set is $+\infty$. If
the minimum on the right is finite, replacing every bound by $v(e_j)>\nu_j$
replaces $\ge$ by $>$ in \eqref{polynomial:eq:valpseudozero}.
\end{theorem}
\begin{proof}
The triangle inequality gives necessity of
\eqref{polynomial:eq:orderpseudozero}. Let $S$ be its right side. If $S=0$ the
inequality forces $P(z)=0$ and all errors may be taken zero. If $S>0$ put, for
$z^j\ne0$,
\begin{equation}\label{polynomial:eq:phaseerrors}
 e_j=-\frac{P(z)}{S}\,\eta_j\,\frac{\overline{z^j}}{|z^j|},
\end{equation}
and $e_j=0$ otherwise. Then $|e_j|\le\eta_j$ and
$\sum_{j\in J}e_jz^j=-\frac{P(z)}{S}\sum_{j\in J}\eta_j|z^j|=-P(z)$, which
proves sufficiency, including $z=0$ with the stated convention for $z^0$. The
phases in \eqref{polynomial:eq:phaseerrors} are algebraic unit-modulus
factors; no global argument or trigonometric branch is needed.

For the valuation version, the valuation inequality for a finite sum gives
necessity. If the minimum is finite, choose $j_0$ attaining it and take
$e_{j_0}=-P(z)/z^{j_0}$ and $e_j=0$ otherwise; the assumed bound says exactly
that this candidate coefficient has the required valuation. If the minimum is
$+\infty$, every perturbation monomial vanishes at $z$ or there are none, so
$Q(z)=P(z)$ and the condition correctly reduces to $P(z)=0$. For a finite
minimum the same reasoning with strict inequalities gives the last assertion,
discarding monomials vanishing at $z$ when checking necessity: the minimum of
finitely many quantities each strictly above its own bound is strictly above
the minimum of those bounds.
\end{proof}

The contrast is exact, not an analogy based on a heuristic ultrametric norm.
The order formula uses a \emph{sum} of available magnitudes, so aligned phases
let several coefficients cooperate, and the cooperation is realized explicitly
by \eqref{polynomial:eq:phaseerrors}. The valuation formula uses a
\emph{minimum}, and one suitably chosen coefficient always suffices. The proof
permits infinite weights and arbitrary surreal centers because every operation
is finite.

\begin{example}[Constant-term perturbations]\label{polynomial:ex:constantuncertainty}
For $P(z)=z^m-t$ with $m\ge1$, a permitted constant error of modulus at most
$\eta$ gives exactly the region $|z^m-t|\le\eta$, and a constant error of
valuation at least $\nu$ gives exactly $v(z^m-t)\ge\nu$. When $\nu>1$ the root
residues at scale $1/m$ cannot change, although their higher terms can. At
$\nu=1$ the choice $e_0=t$ changes the polynomial to $z^m$ and makes all roots
collide at zero. The comparison anticipates the strictness in
\cref{polynomial:ex:sharpdisc}.
\end{example}

\section{The initial polynomial at a center and a scale}
\label{polynomial:sec:initial}

\subsection{Weighted Gauss valuation}\label{polynomial:subsec:gauss}
Fix a nonzero $P\in K[z]$, a center $a\in K$ and a scale exponent
$\rho\in\Gamma$. Its finite Taylor expansion is
$P(a+Z)=\sum_{j=0}^nb_jZ^j$ with $b_j=P^{(j)}(a)/j!$, and we define
\begin{equation}\label{polynomial:eq:weighted}
 w_{a,\rho}(P)=\min_{b_j\ne0}\{v(b_j)+j\rho\}=:\lambda,
 \qquad
 I_{a,\rho}(P)(X)=\st\bigl(t^{-\lambda}P(a+t^\rho X)\bigr)\in\C[X],
\end{equation}
the standard part taken coefficientwise, with $w_{a,\rho}(0)=+\infty$. All
coefficients inside the parentheses have nonnegative valuation and at least
one is a unit, so the \emph{initial polynomial} $I_{a,\rho}(P)$ is nonzero. It
need not be monic and may have degree smaller than $n$: terms of strictly
larger weighted valuation disappear under reduction. This definition uses no
limiting process and is meaningful for every $\rho\in\Gamma$, including
negative thresholds and non-Archimedean surreal exponents. Put
\begin{equation}\label{polynomial:eq:active}
 J_{a,\rho}(P)=\{j:v(b_j)+j\rho=\lambda\},\quad
 j_-=\min J_{a,\rho}(P),\quad j_+=\max J_{a,\rho}(P),
\end{equation}
so that $\ord_0I_{a,\rho}(P)=j_-$ and $\deg I_{a,\rho}(P)=j_+$.

\begin{lemma}[Multiplicativity]\label{polynomial:lem:gaussvaluation}
For nonzero $P,Q\in K[z]$,
\[
 w_{a,\rho}(PQ)=w_{a,\rho}(P)+w_{a,\rho}(Q),\qquad
 I_{a,\rho}(PQ)=I_{a,\rho}(P)\,I_{a,\rho}(Q),
\]
and $w_{a,\rho}(P+Q)\ge\min\{w_{a,\rho}(P),w_{a,\rho}(Q)\}$.
\end{lemma}
\begin{proof}
Normalize both as in \eqref{polynomial:eq:weighted}. Their reductions are
nonzero and their product is nonzero in the domain $\C[X]$, so the normalized
product has least coefficient valuation exactly zero; this proves both
multiplicative identities. The additive inequality holds coefficient by
coefficient before reduction.
\end{proof}

No cancellation assumption about the original coefficients is needed:
cancellation is already accounted for by computing actual valuations and then
multiplying nonzero reductions. This weighted valuation is the Gauss
valuation; unlike the modulus it is ultrametric, and no supremum over a
topological circle is needed to define it.

\subsection{The exact root formula}
\begin{theorem}[Initial polynomial and all residue-direction root counts]
\label{polynomial:thm:initialroots}
Write $P=a_n\prod_{i=1}^n(z-\alpha_i)$ with multiplicities, over $K$. For
every $a\in K$ and $\rho\in\Gamma$ there is a nonzero constant
$C_{a,\rho}\in\C$ with
\begin{equation}\label{polynomial:eq:initialroots}
 \boxed{\displaystyle
 I_{a,\rho}(P)(X)=C_{a,\rho}\!\!\prod_{v(\alpha_i-a)\ge\rho}\!\!
 \left(X-\st\frac{\alpha_i-a}{t^\rho}\right),}
\end{equation}
and moreover
\begin{equation}\label{polynomial:eq:profileproduct}
 w_{a,\rho}(P)=v(a_n)+\sum_{i=1}^n\min\{\rho,\ v(\alpha_i-a)\},
\end{equation}
with the convention $\min(\rho,+\infty)=\rho$ for roots at $a$. Consequently,
counting multiplicity,
\begin{align}
 \#\{\alpha_i\in\Bge a\rho\}&=\deg I_{a,\rho}(P)=j_+,
 \label{polynomial:eq:closedcount}\\
 \#\{\alpha_i\in\Bgt a\rho\}&=\ord_0I_{a,\rho}(P)=j_-,
 \label{polynomial:eq:opencount}\\
 \#\{i:v(\alpha_i-a)=\rho\}&=j_+-j_- .\label{polynomial:eq:shellcount}
\end{align}
For each $c\in\C$, the number of roots of $P$ in the residue-direction ball
$a+t^\rho(c+\mm_K)$ equals the multiplicity of $c$ as a root of
$I_{a,\rho}(P)$.
\end{theorem}
\begin{proof}
By \cref{polynomial:lem:gaussvaluation} both the weighted valuation and the
initial polynomial are multiplicative, so it suffices to treat one linear
factor. If $v(\alpha-a)\ge\rho$, write $\beta=(\alpha-a)/t^\rho\in\OO_K$;
then $a+t^\rho X-\alpha=t^\rho(X-\beta)$, of weighted valuation $\rho$ and
normalized reduction $X-\st(\beta)$, which is $X$ exactly when
$v(\alpha-a)>\rho$. If $v(\alpha-a)<\rho$, put $\delta=v(\alpha-a)$ and
$\ell=\lc(\alpha-a)$; normalizing $t^\rho X-(\alpha-a)$ by $t^{-\delta}$ gives
reduction the nonzero constant $-\ell$ and weighted valuation $\delta$. The
leading scalar $a_n$ contributes its leading complex coefficient and its
valuation. Multiplying these finitely many results proves
\eqref{polynomial:eq:initialroots} and
\eqref{polynomial:eq:profileproduct}. A finite normalized root reduces to $0$
exactly when its valuation is positive and to $c$ exactly when it lies in
$c+\mm_K$; degree, order at zero and multiplicity at $c$ then give all the
counts, and the nonzero coefficients of $I_{a,\rho}(P)$ occur exactly at the
minimizing indices, which identifies $j_\pm$.
\end{proof}

The theorem packages all roots of the chosen ball into a single ordinary
complex polynomial. Roots outside that ball contribute only a nonzero scalar
to the reduction, but their valuations still contribute to
\eqref{polynomial:eq:profileproduct}. It simultaneously allows $a$ to be
infinite, $t^\rho$ to be infinite or infinitesimal, and $\Gamma$ to have
arbitrary ordered-group rank. It is a finite valued-field argument, not an
appeal to analytic compactness of a ball, and a single reduction carries more
information than an unlabelled Newton polygon, since it also distinguishes
directions \emph{within} a scale.

\begin{corollary}[Valuation Rouch\'e: preservation of the whole initial
polynomial]\label{polynomial:cor:valrouche}
Let $P\ne0$ and $E\in K[z]$ satisfy
\begin{equation}\label{polynomial:eq:valrouche}
 w_{a,\rho}(E)>w_{a,\rho}(P),\qquad w_{a,\rho}(0)=+\infty .
\end{equation}
Then $w_{a,\rho}(P+E)=w_{a,\rho}(P)$ and $I_{a,\rho}(P+E)=I_{a,\rho}(P)$;
consequently $P$ and $P+E$ have equal root counts in $\Bge a\rho$, in
$\Bgt a\rho$, and in every residue-direction ball at that center and scale.
\end{corollary}
\begin{proof}
Normalize by $t^{-w_{a,\rho}(P)}$ after substituting $a+t^\rho X$; every
coefficient of the scaled $E$ then has positive valuation, so its reduction is
zero and cannot cancel the nonzero reduction of $P$. Apply
\cref{polynomial:thm:initialroots}.
\end{proof}

The degrees of $P$ and $P+E$ need not agree, and additional roots of the
higher-degree polynomial may lie outside the selected ball --- this is exactly
the situation of \cref{polynomial:ex:escape}; only the specified ball counts
are asserted. The strict inequality is necessary for the general statement:
$P=z$ and $P+E=z+1$ at scale zero have different residue directions although
the error has the same, rather than a smaller, weighted size. The coefficient
test \eqref{polynomial:eq:valrouche} has the benefit of not quantifying over
boundary points at all. The modulus and valuation versions of Rouch\'e have
different hypotheses and different balls, and neither is identified with the
contour theorem on an ordinary thickening in
\cite{polynomial:SourceAnalysis}; this proof uses no transfer and no version
of contour integration.

\section{Newton profiles over an arbitrary ordered exponent group}
\label{polynomial:sec:newton}

\subsection{The polygon without an Archimedean assumption}
At center zero put $\phi_P(\rho)=\min_j\{v(a_j)+j\rho\}$, a minimum of
finitely many affine expressions in the ordered divisible group $\Gamma$. One
organizes it by its finite breakpoints without embedding $\Gamma$ in $\R$:
candidate breakpoints are $(v(a_j)-v(a_k))/(k-j)$, which lie in $\Gamma$ by
divisibility, and only those at which at least two indices attain the minimum
matter. Equivalently, considering the finitely many coefficient points
$(j,v(a_j))$ with $a_j\ne0$, a \emph{lower face of slope $-\rho$} is specified
\emph{algebraically} by the set of indices minimizing $v(a_j)+j\rho$. This
definition works even when $\Gamma$ cannot be embedded in $\R$; we do not
require a literal Euclidean drawing of its vertical axis.

\begin{theorem}[Newton's root-valuation rule]\label{polynomial:thm:newton}
Over $K$, at a threshold $\rho\in\Gamma$ (rather than the added value
$+\infty$), the number of nonzero roots of
valuation $\rho$ is the difference between the largest and smallest indices
attaining $\phi_P(\rho)$, and every valuation of a nonzero root is such a
breakpoint. Equivalently, a lower face joining indices $j<k$ with slope
$-\rho$ accounts for $k-j$ roots of valuation $\rho$, and the residue
polynomial of that face determines their leading complex coefficients with
multiplicity. Roots equal to the center contribute their exact multiplicity
separately, as roots of valuation $+\infty$.
\end{theorem}
\begin{proof}
The count at a fixed scale is \eqref{polynomial:eq:shellcount}. If a nonzero
root has valuation $\rho$, that count is positive, so at least two indices
attain the minimum; conversely a positive index difference produces exactly
that many roots by the same formula. A lower face is precisely a set of
coefficient points on the line $y+\rho x=\phi_P(\rho)$ with all others above
it, so its horizontal span is the difference of the extreme active indices.
The leading complex coefficients are read from
\eqref{polynomial:eq:initialroots}. Alternatively,
\eqref{polynomial:eq:profileproduct} expresses
$\phi_P(\rho)=v(a_n)+\sum_i\min\{\rho,v(\alpha_i)\}$, whose slope changes at
precisely the root valuations, with the same multiplicities. Both arguments
use only order and rational scalar multiplication in $\Gamma$.
\end{proof}

The familiar Newton polygon is therefore a finite encoding of the
root-valuation multiset. The standard rank-one analytic version and valuation
polygons are discussed in \cite[\S3]{polynomial:Benedetto} and in
\cite{polynomial:ChoiLee}; the proof above requires neither a real-valued norm
nor completeness. The one-variable tropical condition that the minimum is
attained at least twice is thus sufficient as well as necessary for a nonzero
root at that valuation --- a consequence of algebraic closedness together with
the exact count, \emph{not} an additional convergence assertion about a Newton
algorithm.

\begin{example}[A cubic with two root scales]\label{polynomial:ex:newtoncubic}
Let $t=\omega^{-1}$ and $P(z)=z^3-t^2z-t^5$. The nonzero coefficient points
are $(0,5),(1,2),(3,0)$, so $\phi_P(\rho)=\min\{5,\,2+\rho,\,3\rho\}$: at
$\rho=1$ the active indices are $1,3$, giving two roots of valuation one, and
at $\rho=3$ they are $0,1$, giving one root of valuation three. Explicitly
\[
\begin{aligned}
 t^{-3}P(tX)&=X^3-X-t^2, & I_{0,1}(P)&=X^3-X=X(X-1)(X+1),\\
 t^{-5}P(t^3X)&=t^4X^3-X-1, & I_{0,3}(P)&=-X-1 .
\end{aligned}
\]
The three residue directions at scale one are $0,1,-1$; the zero direction has
higher valuation and is resolved at the next relevant scale. The simple-root
recursion of \cref{polynomial:cor:simpleroot} gives
\[
 \alpha_+=t+\tfrac12t^3-\tfrac38t^5+O(t^7),\quad
 \alpha_-=-t+\tfrac12t^3+\tfrac38t^5+O(t^7),\quad
 \alpha_0=-t^3-t^7+O(t^{11}).
\]
Here $O(t^k)$ means an integer-power formal Hahn series supported in exponents
at least $k$, not a fine-topological limiting estimate and not a numerical
asymptotic assertion obtained by varying $t$. The critical points are exactly
$\pm t/\sqrt3$, since $P'=3z^2-t^2$; they lie in none of the three root
submonads at scale one, although both lie in their common valuation ball.
\end{example}

\begin{example}[A scale beyond all integer multiples]\label{polynomial:ex:highrank}
In a divisible workspace containing $1$ and $\omega$, take
$P(z)=(z-t)(z-t^\omega)$. Its two roots have valuations $1$ and $\omega$; the
Newton profile has both breakpoints and \cref{polynomial:thm:newton} applies
unchanged. There is no integer $N$ with $N>\omega$, and an iteration that only
raises a valuation through the sequence $1,2,3,\ldots$ cannot be advertised as
reaching the second scale. The finite initial-polynomial construction has no
such restriction. Compare
\cref{polynomial:rem:notalimit}.
\end{example}

\begin{example}[A coefficient with nondiscrete support]\label{polynomial:ex:nondiscrete}
Let $\gamma_k=1-1/(k+2)$ and $A=\sum_{k\ge0}t^{\gamma_k}$. Its support is
increasing and well ordered of order type $\omega$, although all its exponents
lie below $1$ and there are infinitely many of them below $1$. Since
$v(A)=1/2$, both roots of $P(z)=z^2-A$ have valuation $1/4$; at that scale the
reduction is $X^2-1$, so the residues are $1$ and $-1$. Writing
$A=t^{1/2}(1+B)$ with $B=\sum_{k\ge1}t^{\gamma_k-1/2}$ of positive well-ordered
support, the two exact roots are
\begin{equation}\label{polynomial:eq:binomialhahn}
 \alpha_\pm=\pm t^{1/4}\sum_{j\ge0}\binom{1/2}{j}B^j .
\end{equation}
The binomial series is strongly summable by \cref{polynomial:lem:support}
for the positive support $S=\{\gamma_k-1/2:k\ge1\}$. Thus each root has
support in $1/4+S^*$; the shift comes from the displayed factor $t^{1/4}$.
Its use does not require that the original coefficient
series converge after replacing $t$ by a positive real number $q\in(0,1)$; in
fact $q^{\gamma_k}\to q$, not $0$, so the summands would fail to tend to zero.
No such substitution is used in the Hahn construction, and the support is not
an ordinary integer-power series in one monomial.
\end{example}

\section{Polynomial image balls, finite maps, branch values, and ramification}
\label{polynomial:sec:maps}
The initial-polynomial identity determines a polynomial map on an entire
valuation ball, not only its zero fiber. The \emph{scale-local} statement of
\cref{polynomial:thm:imageballs} and the \emph{global algebraic} statement of
\cref{polynomial:thm:finitemap} answer different questions and neither
contains the other: the first is sharp on one prescribed ball and says nothing
about the rest of $K$; the second is exact on every fiber of the whole map and
says nothing about scale. Both are kept.

\subsection{Exact image balls and two mapping degrees}
\begin{theorem}[Exact image and fiber degree]\label{polynomial:thm:imageballs}
Let $P\in K[z]$ be nonconstant, write $P(a+Y)=b_0+\sum_{j\ge1}b_jY^j$ and put
\[
 \lambda=\min_{j\ge1}\{v(b_j)+j\rho\},\quad
 d_-=\min\{j\ge1:v(b_j)+j\rho=\lambda\},\quad
 d_+=\max\{j\ge1:v(b_j)+j\rho=\lambda\}.
\]
Then
\begin{align}
 P\bigl(\Bge a\rho\bigr)&=\Bge{P(a)}{\lambda},
 \label{polynomial:eq:imageclosed}\\
 P\bigl(\Bgt a\rho\bigr)&=\Bgt{P(a)}{\lambda}.
 \label{polynomial:eq:imageopen}
\end{align}
Every target in \eqref{polynomial:eq:imageclosed} has exactly $d_+$ preimages
in the closed source ball, counted with polynomial multiplicity, and every
target in \eqref{polynomial:eq:imageopen} has exactly $d_-$ preimages in the
open source ball.
\end{theorem}
\begin{proof}
For $v(h)\ge\rho$ every term $b_jh^j$ with $j\ge1$ has valuation at least
$\lambda$, and for $v(h)>\rho$ strictly greater; this gives both inclusions.
Let $v(y-P(a))\ge\lambda$. The polynomial $t^{-\lambda}(P(a+t^\rho X)-y)$ has
finite coefficients and reduction of degree $d_+$, since the changed constant
term cannot affect the nonzero top coefficient; by
\cref{polynomial:thm:initialroots} it has exactly $d_+$ roots in $\OO_K$,
proving surjectivity and the closed count. If $v(y-P(a))>\lambda$, its reduced
constant term vanishes and its order at zero is exactly $d_-$, and the
open-ball count of the same theorem finishes the proof.
\end{proof}

This is an arbitrary-rank polynomial counterpart of the usual
non-Archimedean disk-mapping theorem; compare
\cite[Corollary~3.11]{polynomial:Benedetto}. The finite proof shows why
neither a complete real-valued norm nor a restricted-power-series convergence
hypothesis is needed. The two mapping degrees genuinely differ: for
$P(z)=z^2-z$ at $a=0$, $\rho=0$ one has $\lambda=0$, $d_-=1$, $d_+=2$, so the
finite elements map onto the finite elements with degree two while the
infinitesimal ball at zero maps bijectively to itself, the second preimage of
an infinitesimal target lying in the monad of $1$. ``Finite'' here means exact
fiber cardinality with multiplicity on a ball, not topological properness.

\begin{corollary}[Affine normalization of a finite polynomial problem]
\label{polynomial:cor:normalization}
For finitely many nonzero polynomials of bounded finite degree one can choose one
affine change of coordinates over $K$ so that all their roots are finite in
the rescaled coordinate; each nonzero polynomial can then be made monic with
finite coefficients by a scalar normalization. Concretely, for
$P=z^n+\sum_{j<n}a_jz^j$, $n\ge1$, choose $\rho\in\Gamma$ with
$\rho\le\min_{a_j\ne0}v(a_j)/(n-j)$; then $t^{-n\rho}P(t^\rho X)$ is monic with
integral coefficients, and every root of a monic integral polynomial is
integral.
\end{corollary}
\begin{proof}
Choose a workspace containing all the roots and a $\rho$ no larger than the
finitely many valuations of the nonzero root displacements from the selected
center; then every $(\alpha-a)/t^\rho$ is finite. Divide the rescaled
polynomial by its leading coefficient and use Vieta. For the explicit
criterion, the substitution multiplies the coefficient of $z^j$ by
$t^{(j-n)\rho}$, so integrality is exactly $v(a_j)+(j-n)\rho\ge0$. Integrality
of the roots is the argument after
\cref{polynomial:thm:fta}. Only a minimum among finitely many exponents is
used. Constant nonzero polynomials have no roots to constrain the scale;
if all relevant coefficients vanish, any scale works. The zero polynomial
is excluded: its zero set is the whole field, which has elements of unbounded
magnitude when $\Gamma\ne\{0\}$.
\end{proof}

The normalization changes the numerical valuation precision of the
coefficients, so precision must always be measured \emph{after} the stated
change of coordinates. This is what makes the integral-coefficient hypotheses
of \cref{polynomial:sec:stability} harmless, and
\cref{polynomial:subsec:rescaled} records the exact bookkeeping.

\subsection{Every fiber has the same algebraic length}
\begin{theorem}[A polynomial is a finite flat map of its degree]
\label{polynomial:thm:finitemap}
Let $P\in K[z]$ be monic of degree $n\ge1$ and regard $K[z]$ as a
$K[Y]$-algebra by $Y\mapsto P(z)$. It is free of rank $n$ with basis
$1,z,\ldots,z^{n-1}$. For every $y\in K$ the fiber algebra $K[z]/(P(z)-y)$ has
dimension $n$ and decomposes as
\begin{equation}\label{polynomial:eq:fiberdecomp}
 \prod_{P(\alpha)=y}K[u]/(u^{e_\alpha}),\qquad
 e_\alpha=\ord_\alpha(P-y),\qquad \sum_{P(\alpha)=y}e_\alpha=n .
\end{equation}
\end{theorem}
\begin{proof}
The presentation $K[Y,z]/(P(z)-Y)\cong K[z]$ by eliminating $Y$; the defining
polynomial is monic in $z$, so division makes the algebra free over $K[Y]$
with the stated basis, and specializing $Y$ to $y$ leaves the same basis and
rank. \Cref{polynomial:thm:fta,polynomial:thm:crt} give
\eqref{polynomial:eq:fiberdecomp}.
\end{proof}

The words ``finite flat'' describe this free algebra and the associated map of
\emph{set-sized} spectra. They do not assert topological properness or path
lifting in the full fine surreal topology; the ordinary connected
real-parameter paths excluded in \cite{polynomial:SourceAnalysis} play no role.

\begin{theorem}[Branch-value polynomial]\label{polynomial:thm:branchvalues}
Let $P\in K[z]$ be monic of degree $n\ge2$ and list the roots of $P'$ as
$c_1,\ldots,c_{n-1}$ with multiplicity. Then
\begin{equation}\label{polynomial:eq:branchvalues}
 \Disc_z\bigl(P(z)-Y\bigr)=(-1)^{n(n-1)/2}n^n\prod_{j=1}^{n-1}\bigl(P(c_j)-Y\bigr).
\end{equation}
It has degree $n-1$ in $Y$, and its zeros are exactly the values with a
nonreduced fiber, that is, the critical values of $P$, with the multiplicities
indicated by the product.
\end{theorem}
\begin{proof}
Put $A(z)=P(z)-Y$. The monic discriminant identity, the resultant
symmetry sign $(-1)^{n(n-1)}=1$, and the leading coefficient $n$ of $P'$
give
\[
 \Disc_zA=(-1)^{n(n-1)/2}\Res_z(A,P'),\qquad
 \Res_z(A,P')=\Res_z(P',A)=n^n\prod_{j=1}^{n-1}A(c_j),
\]
which together are \eqref{polynomial:eq:branchvalues}. A fiber has a multiple root exactly when
$A$ and $A'=P'$ have a common root, so the zero set is as claimed; the leading
coefficient in $Y$ is nonzero in characteristic zero.
\end{proof}

Coincident critical values combine their multiplicities in this product;
there need not be $n-1$ distinct branch values. The formula continues to make
sense over coherent coefficient rings as a universal polynomial identity, even
when the values cannot be labelled globally.

\begin{proposition}[Finite and infinite ramification]\label{polynomial:prop:ramification}
For a nonconstant degree-$n$ polynomial over $K$, the local index at a finite
point $\alpha$ is $e_\alpha=\ord_\alpha(P-P(\alpha))$ and
\begin{equation}\label{polynomial:eq:finiteRH}
 \sum_{\alpha\in K}(e_\alpha-1)=n-1,
\end{equation}
only finitely many summands being nonzero. At infinity the local index is $n$,
so the total sum on the algebraic projective line, including infinity, is
$2n-2$.
\end{proposition}
\begin{proof}
The finite Taylor expansion at $\alpha$ begins with a nonzero term of degree
$e_\alpha$, so the multiplicity of $\alpha$ in $P'$ is $e_\alpha-1$; summing
the multiplicities of the degree-$(n-1)$ polynomial $P'$ gives
\eqref{polynomial:eq:finiteRH}. At infinity use $u=1/z$, $w=1/Y$: with leading
coefficient $a_n$, $1/P(1/u)=a_n^{-1}u^n(1+\cdots)^{-1}$ has a unit factor with
nonzero constant term, so the order is $n$; adding $n-1$ finishes. This is the
polynomial instance of the ramification formula, derived directly rather than
from a topological covering or surface argument.
\end{proof}

\begin{corollary}[Surjectivity and polynomial injections]
\label{polynomial:cor:polynomialmaps}
Every nonconstant polynomial $\SC\to\SC$ is surjective, and the injective ones
are exactly the affine polynomials with nonzero slope.
\end{corollary}
\begin{proof}
Algebraic closure supplies a root of $P-y$ for every target. If $n\ge2$,
choose $y$ outside the finite set of critical values of
\cref{polynomial:thm:branchvalues}; its fiber has $n$ distinct roots, so the
map is not injective. The affine case is immediate. This proof is finite and
does not require an analytic inverse theorem.
\end{proof}

\section{Critical points at every scale}\label{polynomial:sec:critical}
The residue field is $\C$, so every nonzero integer is a unit by
\eqref{polynomial:eq:integersunit}. That seemingly small fact makes the
following counts exact, and they would fail in positive residue
characteristic.

\subsection{Differentiation commutes with nonconstant reduction}
\begin{lemma}[The derivative of an initial polynomial]
\label{polynomial:lem:initialderivative}
Over $K$, if $I=I_{a,\rho}(P)$ has degree $k\ge1$ then
\begin{equation}\label{polynomial:eq:initialderivative}
 w_{a,\rho}(P')=w_{a,\rho}(P)-\rho,\qquad I_{a,\rho}(P')=I' ,
\end{equation}
and more generally $w_{a,\rho}(P^{(r)})=w_{a,\rho}(P)-r\rho$ and
$I_{a,\rho}(P^{(r)})=I^{(r)}$ for $1\le r\le k$.
\end{lemma}
\begin{proof}
Differentiating $t^{-\lambda}P(a+t^\rho X)$ with respect to $X$ gives
$t^{\rho-\lambda}P'(a+t^\rho X)$, whose coefficients remain finite and whose
reduction is $I'$. Since $\C$ has characteristic zero and $I$ is nonconstant,
$I'\ne0$, so the displayed normalization has weighted valuation exactly zero,
which is \eqref{polynomial:eq:initialderivative}. For higher derivatives
$I^{(r)}\ne0$ whenever $r\le k$ and the same argument applies.
\end{proof}

The hypothesis cannot simply be removed: a constant initial polynomial has
zero derivative and does not determine the next nonzero scale of $P'$.
Likewise the argument can fail in positive residue characteristic, where a
nonconstant reduction may have zero derivative.

\begin{theorem}[Exact critical-point conservation in occupied balls]
\label{polynomial:thm:criticalballs}
Let $P\in K[z]$ be nonconstant and let $B$ be either $\Bge a\rho$ or
$\Bgt a\rho$. If $B$ contains exactly $k\ge1$ roots of $P$ counted with
multiplicity, then for every $1\le r\le k$ it contains exactly $k-r$ roots of
$P^{(r)}$, counted with multiplicity; in particular exactly $k-1$ critical
points of $P$. Moreover, if $I=I_{a,\rho}(P)$ has positive degree, the residue
directions and critical multiplicities inside the closed ball are read off
from $I'$.
\end{theorem}
\begin{proof}
For the closed ball, $\deg I=k$ by
\eqref{polynomial:eq:closedcount}; \cref{polynomial:lem:initialderivative}
gives an initial polynomial $I^{(r)}$ of degree $k-r$ for $P^{(r)}$, and
\cref{polynomial:thm:initialroots} converts that into the closed-ball count.
For the open ball, $\ord_0I=k$ by \eqref{polynomial:eq:opencount}, so
$\deg I\ge k$ and the lemma still applies; in characteristic zero
$\ord_0I^{(r)}=k-r$ in that range, which is the open-ball count. For the last
assertion apply \cref{polynomial:thm:initialroots} to $P'$, whose initial
polynomial is exactly $I'$ by \cref{polynomial:lem:initialderivative}.
\end{proof}

\begin{remark}[A root-free ball is a different case]\label{polynomial:rem:rollehyp}
One must not substitute $k=0$. The open infinitesimal ball at zero contains no
root of $z^2-1$ but does contain the critical point zero. No formula with the
meaningless count $k-1=-1$ is intended. Such critical points occur in residue
directions \emph{between} distinct root branches, which is exactly what
\cref{polynomial:thm:branch} makes explicit. In positive residue
characteristic even the nonempty-ball argument can fail, because $I'$ may lose
its expected degree or vanish identically.
\end{remark}

\begin{corollary}[Ramification count on an arbitrary mapping ball]
\label{polynomial:cor:ramificationball}
With the notation of \cref{polynomial:thm:imageballs}, the closed source ball
contains exactly $d_+-1$ critical multiplicities of $P$ and the open source
ball exactly $d_--1$. In either case, mapping degree one is equivalent to
the absence of a critical point and implies bijectivity onto the image ball.
The converse from bijectivity holds for closed balls, and for open balls
provided $\Gamma\ne\{0\}$.
\end{corollary}
\begin{proof}
Apply \cref{polynomial:thm:criticalballs} to $P-P(a)$, whose closed and open
root counts are $d_+$ and $d_-$ by \cref{polynomial:thm:imageballs}, and whose
derivative is $P'$. Degree one supplies one simple preimage for every target.
Conversely, if the degree exceeds one, choose a target that is not one of the
finitely many critical values. The closed image ball contains
$P(a)+c t^\lambda$ for all $c\in\C$, so is infinite. For an open image ball
and $0<\eta\in\Gamma$, the points $P(a)+t^{\lambda+n\eta}$, $n\ge1$, are
distinct points of that ball. In the asserted cases a noncritical target
therefore exists; its degree-many preimages are simple and distinct, so the
map is not injective.
\end{proof}

The nontriviality condition cannot be dropped from the open-ball converse.
For $\Gamma=\{0\}$, one has $K=\C$ and $\Bgt00=\{0\}$. The restriction of
$P(z)=z^2$ to this singleton is bijective onto its singleton image, although
$d_-=2$ and zero is a critical point. All the multiplicity counts above remain
valid. Bijectivity of one exceptional fiber does not determine mapping degree.

\begin{corollary}[Valuative Gauss--Lucas and nearest neighbours]
\label{polynomial:cor:nearest}
Every valuation ball containing all roots of $P$ contains all its critical
points. If $P\in K[z]$ is squarefree of degree at least two, then
\begin{equation}\label{polynomial:eq:nearest}
 \max_{P'(w)=0}v(w-\alpha_i)=\delta_i=\max_{j\ne i}v(\alpha_i-\alpha_j).
\end{equation}
Thus each root has a critical point at exactly its nearest-neighbour
valuation, and no critical point is closer in the valuation sense.
\end{corollary}
\begin{proof}
The first assertion is \cref{polynomial:thm:criticalballs} with $k=\deg P$.
For the second, $\Bgt{\alpha_i}{\delta_i}$ contains only $\alpha_i$, whereas
$\Bge{\alpha_i}{\delta_i}$ contains at least two roots; so the open ball
contains no critical point and the closed ball contains at least one, and
their difference is the shell $v(w-\alpha_i)=\delta_i$.
\end{proof}

This belongs to the non-Archimedean root-location picture, which has
precedents such as \cite{polynomial:ChoiLee,polynomial:Faber}. It is
\emph{not} a proof of a Euclidean root--critical-point distance conjecture:
valuation distance ignores the ordinary leading coefficient of a modulus, and
its balls are different objects from order disks. Our proof is a direct
arbitrary-rank polynomial argument and uses no Berkovich space.

\subsection{The branch polynomial between child clusters}
\begin{theorem}[Critical residue directions inside a cluster]
\label{polynomial:thm:branch}
Suppose the closed ball $\Bge a\rho$ over $K$ contains $k\ge1$ roots and its
nonempty root residue directions are the distinct $c_1,\ldots,c_s\in\C$ with
multiplicities $m_1,\ldots,m_s$, $\sum_jm_j=k$, so that
\begin{equation}\label{polynomial:eq:clusterinitial}
 I_{a,\rho}(P)=C\prod_{j=1}^s(X-c_j)^{m_j},\qquad C\ne0 .
\end{equation}
Then each occupied child ball $a+t^\rho(c_j+\mm_K)$ contains exactly $m_j-1$
critical multiplicities, and the remaining critical points of the parent ball
have total multiplicity $s-1$, with residue directions given by the roots,
with multiplicity, of the \emph{branch polynomial}
\begin{equation}\label{polynomial:eq:branchpoly}
 W(X)=\sum_{j=1}^sm_j\prod_{\ell\ne j}(X-c_\ell).
\end{equation}
It has degree $s-1$, leading coefficient $k\ne0$, and no root equal to any
$c_j$.
\end{theorem}
\begin{proof}
Differentiating \eqref{polynomial:eq:clusterinitial} gives
$I'=C\prod_j(X-c_j)^{m_j-1}W(X)$. The leading coefficient of $W$ is
$\sum_jm_j=k\ne0$, so $\deg W=s-1$, and
$W(c_j)=m_j\prod_{\ell\ne j}(c_j-c_\ell)\ne0$. Now read every critical
direction and multiplicity from $I'$ by
\cref{polynomial:lem:initialderivative,polynomial:thm:criticalballs}. Counting,
$(k-1)-\sum_j(m_j-1)=s-1$.
\end{proof}

For two child directions the intervening critical direction is
\[
 \frac{m_1c_2+m_2c_1}{m_1+m_2},
\]
with the \emph{opposite} cluster multiplicities as weights; it is generally
not the centroid of the roots. For $X^m(X-1)$ the nonzero critical direction
is $m/(m+1)$ whereas the root centroid is $1/(m+1)$. More generally $W$ is a
weighted incomplete polynomial in the sense of
\cref{polynomial:prop:weighted}, so its roots lie in the ordinary complex
convex hull of the child directions. This is a useful bridge between the two
geometries: the valuation locates the scale, and an ordinary complex
polynomial locates the leading directions at that scale. The roots of $W$ may
themselves repeat and are counted with multiplicity, and $W$ may be applied
again to any smaller root-containing ball.

\subsection{The finite root-cluster tree}\label{polynomial:subsec:tree}
For a finite set of distinct roots with positive multiplicities build a tree
as follows. If more than one distinct root is present, take the full root set
as the root node. For a node $C$ containing at least two distinct roots set
$\delta_C=\min_{\alpha\ne\beta\in C}v(\alpha-\beta)$; by the ultrametric
inequality $v(\alpha-\beta)>\delta_C$ is an equivalence relation on $C$, and it
has at least two classes since the minimum is attained. These classes are the
occupied residue-direction children of $C$. Recurse until the children are
singletons; a leaf is a single root location retaining its multiplicity.
Recursion strictly reduces the number of distinct roots in a child, so the
tree is finite even when the valuations have arbitrarily high ordered-group
rank. Its depths belong to $\Gamma$, not necessarily to $\R$.

Every internal node $C$ is exactly the set of roots in $\Bge a{\delta_C}$ for
$a\in C$: a child lies in one residue subball of its parent and has strictly
larger internal depth, while all roots outside it are separated at the
parent's depth or earlier. Each internal node therefore has an actual scale
and a branch polynomial \eqref{polynomial:eq:branchpoly}.

\begin{corollary}[Critical allocation by the root tree]\label{polynomial:cor:tree}
For a degree-$n$ polynomial over $K$ with $d$ distinct roots, an internal node
with $s$ occupied children accounts for exactly $s-1$ critical multiplicities
in the part of its ball outside those children, and a leaf of root
multiplicity $m$ accounts for $m-1$ at that root. These allocations account
for all $n-1$ critical multiplicities, since
\[
 \sum_{\text{internal nodes}}(s-1)=d-1,\qquad
 \sum_{\text{leaves}}(m-1)=n-d .
\]
\end{corollary}
\begin{proof}
At an internal node \cref{polynomial:thm:branch} gives $s-1$ outside its
occupied children. A child carrying $m$ total root multiplicities contains
$m-1$ critical multiplicities, and its smallest closed root ball contains that
same number by \cref{polynomial:thm:criticalballs}, so none are lost in the
intermediate part of the child ball. At a singleton leaf the exact root
multiplicity gives $m-1$ at the root. The tree identity follows by counting
edges: the sum of all child counts equals the number of nodes minus one. If
$d=1$ the direct multiplicity calculation already gives the assertion. The top
ball contains all $n-1$ critical multiplicities, so the allocation is
exhaustive.
\end{proof}

This is an allocation theorem, not a claim that pairwise root valuations alone
determine every critical point. The residue directions $c_j$ carry additional
data, and further subclustering among the roots of a branch polynomial can
require its own rescaling analysis.

\begin{example}[Critical points between nested roots]\label{polynomial:ex:nested}
Let $t=\omega^{-1}$ and $P(z)=z(z-t^3)(z-t)(z-1)$, whose roots form the nested
clusters $\{0,t^3\}\subset\{0,t^3,t\}\subset\{0,t^3,t,1\}$.
\begin{center}
\begin{tabular}{@{}ccll@{}}
\toprule
$\rho$ & $w_{0,\rho}(P)$ & $I_{0,\rho}(P)$ & $I_{0,\rho}(P')$\\
\midrule
$0$ & $0$ & $X^3(X-1)$ & $X^2(4X-3)$\\
$1$ & $3$ & $-X^2(X-1)$ & $-X(3X-2)$\\
$3$ & $7$ & $X(X-1)$ & $2X-1$\\
\bottomrule
\end{tabular}
\end{center}
The three critical points are simple, with leading terms
\begin{equation}\label{polynomial:eq:nestedcritical}
 w_{\mathrm{outer}}=\tfrac34+\text{infinitesimal},\quad
 w_{\mathrm{middle}}=t\bigl(\tfrac23+\text{infinitesimal}\bigr),\quad
 w_{\mathrm{inner}}=t^3\bigl(\tfrac12+\text{infinitesimal}\bigr).
\end{equation}
At scale zero two critical multiplicities remain in the zero monad and one has
direction $3/4$; at scale one, one remains in the smaller zero child and one
has direction $2/3$; at scale three the final critical point has direction
$1/2$. The finite tree has three binary splits, each supplying one intervening
critical point. Directly,
$P'(z)=4z^3-3(1+t+t^3)z^2+2(t+t^3+t^4)z-t^4$, whose coefficient profile
verifies the critical valuations $0,1,3$; the branch calculation additionally
supplies their leading complex coefficients and shows why all three are
simple.
\end{example}

\subsection{Two layer formulas for the discriminant valuation}
\label{polynomial:subsec:disclayers}
The discriminant valuation of a monic polynomial with distinct finite roots
can be decomposed either by \emph{levels} or by \emph{tree nodes}. The two
decompositions are equivalent; each is convenient for a different purpose, and
both are recorded.

\begin{proposition}[Discriminant as a finite level sum]\label{polynomial:prop:clusterdisc}
Let $P\in K[z]$ be monic with distinct finite roots. List the distinct
positive pairwise distance valuations as $0<\lambda_1<\cdots<\lambda_r$, put
$\lambda_0=0$, and for each $k$ let $\mathcal C_k$ be the classes of the
equivalence $\alpha_i\sim_k\alpha_j\iff v(\alpha_i-\alpha_j)\ge\lambda_k$.
Then
\begin{equation}\label{polynomial:eq:clusterdisc}
 v(\Disc P)=2\sum_{k=1}^r(\lambda_k-\lambda_{k-1})
 \sum_{C\in\mathcal C_k}\binom{|C|}{2}.
\end{equation}
If there are no positive pairwise valuations the sum is zero.
\end{proposition}
\begin{proof}
The ultrametric inequality makes $\sim_k$ an equivalence relation, and all
pairwise distance valuations are nonnegative because the roots are finite. A
pair of distance valuation $\lambda_j$ lies in a common class at levels
$k\le j$ and not above, so its contribution before the factor $2$ is
$\sum_{k\le j}(\lambda_k-\lambda_{k-1})=\lambda_j$; a pair of valuation zero
contributes nothing. Sum over the finitely many pairs and use
\eqref{polynomial:eq:discval}.
\end{proof}

\begin{proposition}[Discriminant as a sum over the root tree]
\label{polynomial:prop:disctree}
Let $P\in K[z]$ be monic and squarefree of degree $n\ge2$. With the tree of
\cref{polynomial:subsec:tree},
\begin{equation}\label{polynomial:eq:disctree}
 v(\Disc P)=n(n-1)\delta_{\mathrm{root}}
 +\sum_{\substack{C\ \text{internal}\\ C\ne\mathrm{root}}}
 |C|\bigl(|C|-1\bigr)\bigl(\delta_C-\delta_{\mathrm{parent}(C)}\bigr).
\end{equation}
The depths $\delta_C$ can be negative, positive, or non-Archimedean elements
of $\Gamma$. Formulas \eqref{polynomial:eq:clusterdisc} and
\eqref{polynomial:eq:disctree} agree whenever both apply.
\end{proposition}
\begin{proof}
The separation valuation of a pair of distinct roots is the depth of its least
common ancestor: start at the root depth and add depth increments for every
proper internal node containing both roots, a telescoping sum. Summing over
unordered pairs counts $\binom{|C|}{2}$ pairs at each node; multiplying by two
and using \eqref{polynomial:eq:discval} gives
\eqref{polynomial:eq:disctree}. For the comparison, both expressions compute
$2\sum_{i<j}v(\alpha_i-\alpha_j)$, the first by grouping the pairs according
to which levels $\lambda_k$ they survive and the second by grouping them
according to their least common ancestor. When the roots are finite,
$\delta_{\mathrm{root}}\ge0$ and the classes $\mathcal C_k$ are exactly the
node sets of the tree cut at height $\lambda_k$.
\end{proof}

Both formulas are meaningful in any ordered divisible value group. They are
finite telescoping identities, not integrals with respect to a real-valued
radius. They explain why discriminant precision can protect a whole nested
root configuration, which is the content of
\cref{polynomial:cor:discprecision}.

\section{Support-controlled factor lifting}\label{polynomial:sec:hensel}
The preceding sections use finite algebra after the coefficients are given.
Constructing factors from a perturbation may require infinitely many Hahn
coefficients, and there the support convention becomes substantive. The
argument is a coprime-factor version of the preparation recursions in
\cite{polynomial:SourceAnalysis,polynomial:SourceGlobal}.

\subsection{The ordinary finite linear splitting}
Let $p_1,\ldots,p_s\in\C[z]$ be pairwise coprime monic polynomials of degrees
$d_i\ge1$, $n=\sum_id_i$, $P_0=\prod_ip_i$. The map
\begin{equation}\label{polynomial:eq:factorlinear}
 L:\bigoplus_{i=1}^s\C[z]_{<d_i}\longrightarrow\C[z]_{<n},
 \qquad (h_i)\longmapsto\sum_ih_i\prod_{j\ne i}p_j
\end{equation}
is an isomorphism: reducing a zero image modulo $p_i$ forces $h_i=0$, because
the other factors are invertible modulo $p_i$, and source and target have the
same finite dimension. This is the linearization of multiplication of monic
factors, and the only linear inverse used below.

\begin{theorem}[Support-controlled coprime Hensel factorization]
\label{polynomial:thm:hensel}
Let $P=P_0+E\in\OO_K[z]$ be monic of degree $n$ with $\deg E<n$ and all
coefficients of $E$ infinitesimal, and let $S\subseteq\Gamma_{>0}$ be the
union of their nonzero supports. There are unique monic $P_i=p_i+A_i$ with
$\deg A_i<d_i$ and corrections of positive support such that
\begin{equation}\label{polynomial:eq:henselfactor}
 P=\prod_{i=1}^sP_i .
\end{equation}
All coefficients of all $A_i$ have support in $S^*\setminus\{0\}$, and the
$P_i$ remain pairwise coprime over $K$. Divisibility of $\Gamma$ is not
required.
\end{theorem}
\begin{proof}
Expanding the product, the equation reads
$E=L(A_1,\ldots,A_s)+N(A_1,\ldots,A_s)$, where every monomial of $N$ contains
at least two correction factors. Proceed by well-founded recursion on
$S^*\setminus\{0\}$. At exponent $\nu$ the nonlinear coefficient $N_\nu$
depends only on correction coefficients at exponents strictly smaller than
$\nu$, since the exponents of the two or more corrections sum to $\nu$ and
each is positive; by \cref{polynomial:lem:support} there are finitely many
contributions. Define
\begin{equation}\label{polynomial:eq:henselrecursion}
 (A_{1,\nu},\ldots,A_{s,\nu})=L^{-1}\bigl(E_\nu-N_\nu\bigr).
\end{equation}
The degree bounds persist because both the linear splitting and every product
difference in the monic product have degree less than $n$, and all factors
remain monic of the prescribed degrees. The union of the constructed supports
lies in the well-ordered $S^*$, and
\eqref{polynomial:eq:henselrecursion} verifies the product at every exponent.

For uniqueness, compare two positive-support factorizations and take the least
exponent at which any correction differs, using the finite union of their
well-ordered supports. At that exponent every product term involving an
additional positive correction uses only smaller differences and vanishes, so
the remaining equation is $L$ applied to the vector of first differences
equals zero; injectivity of $L$ contradicts the choice of that exponent. This
proves uniqueness among all positive-support candidates, not only those
preassigned to $S^*$.

Finally the coefficient matrix of the finite B\'ezout map of a pair $P_i,P_j$
reduces to the corresponding invertible matrix for $p_i,p_j$, so its
determinant has nonzero standard part and is a unit of $\OO_K$; the adjugate
gives an inverse and hence a B\'ezout identity, so $P_i,P_j$ are coprime.
\end{proof}

No iteration through only the natural numbers is assumed cofinal in the output
support: \eqref{polynomial:eq:henselrecursion} is genuinely a recursion on a
well-ordered support monoid, and
\cref{polynomial:rem:notalimit,polynomial:rem:belowbound} are exactly what
must be avoided. The certificate $S^*$ is the universal support bound of this
article.

\begin{corollary}[Simple-root lifting and the first correction]
\label{polynomial:cor:simpleroot}
Every simple root $c$ of $P_0$ lifts to a unique root $c+\eta$ of $P$ with
$\eta\in\mm_K$ and $\supp\eta\subseteq S^*\setminus\{0\}$. If $S$ is
nonempty, at the least positive error exponent $e=\min S$ the coefficient is
\begin{equation}\label{polynomial:eq:firstrootcorrection}
 \eta_e=-\frac{E_e(c)}{P_0'(c)} .
\end{equation}
There may be cancellation making this coefficient zero; no nonvanishing of the
first correction is assumed. If $E=0$, then $S=\varnothing$ and $\eta=0$;
there is no first error exponent.
\end{corollary}
\begin{proof}
Apply \cref{polynomial:thm:hensel} to the coprime pair $z-c$ and
$P_0/(z-c)$, omitting the latter if it is $1$; the lifted first factor is
$z-c-\eta$, and for any root reducing
to $c$, uniqueness of this linear factor identifies it with $c+\eta$.
When $S\ne\varnothing$, expanding $P(c+\eta)=0$ at its least exponent gives
$P_0'(c)\eta_e+E_e(c)=0$.
\end{proof}

\begin{corollary}[Canonical standard-part cluster factors]
\label{polynomial:cor:clusterfactor}
For monic $P\in\OO_K[z]$ write $\st(P)=\prod_{c\in A}(z-c)^{m_c}$ with
$A\subset\C$ finite. Then $P$ has unique monic factors $P_c$ with
$\st(P_c)=(z-c)^{m_c}$, and the roots of $P_c$ are exactly the roots of $P$ in
$c+\mm_K$.
\end{corollary}
\begin{proof}
Apply \cref{polynomial:thm:hensel} to the pairwise coprime ordinary factors.
Every root of $P_c$ is finite and its standard part is a root of $\st(P_c)$ by
\eqref{polynomial:eq:reductionfactor}, so all of them lie over $c$; the
product identity and degree counts show none are omitted.
\end{proof}

Coprimality is essential for the support conclusion. Individual roots of a
multiple cluster need not have supports in the original exponent group: for
$P=z^m-t\in K_\Z[z]$ the reduction $z^m$ has a repeated root, and the
individual roots $\zeta t^{1/m}$ require a divisible workspace and lie outside
the monoid generated by $1$. The theorem controls the \emph{coprime cluster
factors}, not a nonexistent simple-root splitting of their common multiple
ordinary reduction. Symmetric cluster coefficients can have substantially
simpler supports than their separately labelled roots.

\begin{example}[An explicit two-factor recursion]\label{polynomial:ex:hensel}
For $P=z^2+tz-1+t^2$ the ordinary factors are $z-1$ and $z+1$, and the
normalized roots have expansions
\[
 \alpha_+=1-\tfrac12t-\tfrac38t^2-\tfrac9{128}t^4+O(t^6),\qquad
 \alpha_-=-1-\tfrac12t+\tfrac38t^2+\tfrac9{128}t^4+O(t^6),
\]
obtainable either from \eqref{polynomial:eq:henselrecursion} or from
$\alpha_\pm=(-t\pm\sqrt{4-3t^2})/2$ with the strongly summable binomial
series. The two factors are $z-\alpha_+$ and $z-\alpha_-$; their sum and
product have the exact prescribed coefficients, and all displayed support lies
in the positive integer monoid generated by the input.
\end{example}

\subsection{Holomorphic parameters on one fixed ordinary domain}
\begin{theorem}[Parameter factor lifting without repeated shrinking]
\label{polynomial:thm:parameterhensel}
Let $U\subseteq\C^d$ be one ordinary domain, fixed throughout, and let
$A_1(x,z),\ldots,A_s(x,z)\in\OO(U)[z]$ be monic of fixed degrees $d_i$ and
pairwise coprime at \emph{every} $x\in U$. Let
$P=\prod_jA_j+E\in\HHn(U)[z]$ be monic of the same total degree $n$, with $E$
of $z$-degree less than $n$ and of strictly positive Hahn support with
ordinary holomorphic coefficients on $U$, and let $S$ be its total positive
coefficient support. Then \eqref{polynomial:eq:henselfactor} holds with unique
monic factors $A_j+B_j$, $B_j\in\HHp(U)[z]$, $\deg B_j<d_j$, whose corrections
have support in $S^*\setminus\{0\}$, and every ordinary coefficient function
is holomorphic on the \emph{original} domain $U$. No shrinking of $U$ occurs
at any step.
\end{theorem}
\begin{proof}
In the finite power bases, \eqref{polynomial:eq:factorlinear} is a square
matrix $L(x)$ of ordinary holomorphic functions on $U$. Its value is
invertible at every $x\in U$ by pairwise coprimality, so
$L(x)^{-1}=\adj(L(x))/\det L(x)$ is holomorphic on all of $U$; there is no
requirement that its ordinary size be uniformly bounded. Apply
\eqref{polynomial:eq:henselrecursion} coefficientwise with this fixed
holomorphic matrix inverse. Finite products and applications of $L^{-1}$
preserve holomorphy on the same $U$ and introduce no Hahn exponent, so the
support and uniqueness arguments are unchanged.
\end{proof}

This theorem is stated over the fixed-domain ring $\HHn(U)$ and is proved only
there. It is a finite-dimensional parameter version of the matrix inversion
principle in \cite{polynomial:SourceResearch}. It requires an ordinary coprime
factorization on the given base; it does not manufacture globally labelled
roots through ordinary monodromy or across a discriminant locus. When the
leading ordinary polynomial has only simple roots, the same argument on its
ordinary root cover gives a unique coherent lift of every sheet, and local
lifts agree on overlaps by uniqueness, so ordinary analytic continuation
permutes the lifted sheets exactly as it permutes the original ones. That does
\emph{not} produce a globally labelled root on a base with nontrivial ordinary
monodromy, and it asserts nothing about holomorphic individual roots through a
collision.

\section{Perturbation: three theorems with three hypotheses}
\label{polynomial:sec:stability}
Three genuinely different perturbation statements are proved below and are
deliberately \emph{not} merged into one. They differ in hypothesis and in what
they deliver.

\begin{center}
\begin{tabularx}{\textwidth}{@{}P{0.20\textwidth}P{0.26\textwidth}L@{}}
\toprule
Statement & Hypothesis on $P$ & What it gives\\
\midrule
\cref{polynomial:thm:holder} & none beyond monic integral; repeated roots
allowed & a matching of root \emph{multisets} with
$v(\alpha_i-\beta_i)\ge\varepsilon/n$, and $\varepsilon/n$ is optimal
uniformly\\[0.5em]
\cref{polynomial:thm:stability} & distinct finite roots, and
$\varepsilon>T(P)=\max_i(d_i+\delta_i)$ & a \emph{unique} labelling, the whole
matrix $v(\beta_i-\beta_j)=d_{ij}$ with leading coefficients, exact
displacement valuations, a second-order Newton error, $v(Q'(\beta_i))=d_i$ and
$v(\Disc Q)=v(\Disc P)$\\[0.5em]
\cref{polynomial:cor:discprecision} & distinct finite roots and
$\varepsilon>v(\Disc P)$ & everything in \cref{polynomial:thm:stability}, from
a single discriminant valuation; it can require strictly more precision\\
\bottomrule
\end{tabularx}
\end{center}

\noindent
It must be emphasized that \cref{polynomial:thm:holder} is \emph{not} a weaker
version of \cref{polynomial:thm:stability} and does not follow from it. It
assumes no root separation whatsoever --- in particular it applies to
polynomials with repeated roots, where $T(P)$ is undefined --- and its
conclusion is correspondingly much weaker. The two have different hypotheses
and both belong.

\subsection{Optimal valuation-H\"older matching of root multisets}
If $P,Q\in\OO_K[z]$ are monic of degree $n$ with coefficients agreeing to
valuation at least $\varepsilon$, evaluating $Q-P$ at a root of $P$ shows that
\emph{some} root of $Q$ is close. That alone does not produce a bijection of
root multisets: several roots could be sent to the same nearby root. The next
proof balances multiplicities in every relevant ultrametric cluster before
matching.

\begin{theorem}[Optimal valuation-H\"older root matching]
\label{polynomial:thm:holder}
Let $P=z^n+\sum_{j<n}a_jz^j$ and $Q=z^n+\sum_{j<n}b_jz^j$ with
$a_j,b_j\in\OO_K$, $n\ge1$, and suppose $v(a_j-b_j)\ge\varepsilon>0$ for every
$j<n$. Then there are listings of their roots, repeated with multiplicity,
such that
\begin{equation}\label{polynomial:eq:holder}
 \boxed{\displaystyle v(\alpha_i-\beta_i)\ \ge\ \frac{\varepsilon}{n}
 \qquad(1\le i\le n).}
\end{equation}
The exponent $\varepsilon/n$ is optimal uniformly over all such polynomials.
\end{theorem}
\begin{proof}
All roots of both polynomials are finite. Put $\rho=\varepsilon/n$. In the
finite combined root multiset, inspect the valuations of the nonzero pairwise
differences that are strictly less than $\rho$, and choose $\mu$ with
$0<\mu<\rho$ greater than all of these finitely many valuations; such a $\mu$
exists by divisibility --- take the midpoint between $\rho$ and the maximum of
that finite set together with zero.

On the finite collection of root locations the equivalence relations
$v(x-y)\ge\mu$ and $v(x-y)\ge\rho$ now coincide, and their classes are
intersections with closed $\mu$-balls. Take any class and choose a root $a$ in
it, from either polynomial. Translation by the finite center $a$ preserves the
lower bound $\varepsilon$ for the coefficients of $E=Q-P$, since each
translated coefficient is a finite sum of an original coefficient times an
integer and a nonnegative power of $a$, and
\eqref{polynomial:eq:integersunit} applies. Dilation by $t^\mu$ with $\mu>0$
gives $w_{a,\mu}(E)\ge\varepsilon$. On the other hand
\eqref{polynomial:eq:profileproduct} yields
\[
 w_{a,\mu}(P)=\sum_{i=1}^n\min\{\mu,\,v(\alpha_i-a)\}\le n\mu<\varepsilon .
\]
So \cref{polynomial:cor:valrouche} says that this class contains exactly the
same number of $P$-roots as $Q$-roots, with multiplicity; the same argument
rules out a class containing roots of only one polynomial. Match the two
finite multisets arbitrarily within each class. Its members have pairwise
differences of valuation at least $\rho$, so every such pairing satisfies
\eqref{polynomial:eq:holder}.

For optimality take $P=z^n$ and $Q=z^n-t^\varepsilon$: every root of $Q$ has
valuation exactly $\varepsilon/n$ while every root of $P$ is zero, so no
larger uniform bound is possible.
\end{proof}

The proof uses finitely many roots and their finite separation data. There is
no passage to a limit of roots as an auxiliary radius increases, and it works
when $\Gamma$ has arbitrary ordered-group rank. The correspondence is not
generally unique: repeated or unresolved clusters naturally allow several
labellings.

\begin{corollary}[Simultaneous stability for derivatives]
\label{polynomial:cor:derivativeholder}
Under the hypotheses of \cref{polynomial:thm:holder}, for every integer
$0\le r<n$ the root multisets of $P^{(r)}$ and $Q^{(r)}$ admit a matching with
$v(\alpha_i^{(r)}-\beta_i^{(r)})\ge\varepsilon/(n-r)$.
\end{corollary}
\begin{proof}
Divide each $r$th derivative by the nonzero integer $n(n-1)\cdots(n-r+1)$,
making it monic of degree $n-r$. Its coefficients remain finite and the
coefficient differences retain valuation at least $\varepsilon$, because every
nonzero rational factor has valuation zero by
\eqref{polynomial:eq:integersunit}. Apply \cref{polynomial:thm:holder}.
\end{proof}

For first derivatives the denominator is $n-1$, not $n$, and it is sharp:
$P=z^n$ and $Q=z^n-t^\varepsilon z$ change the normalized derivative from
$z^{n-1}$ to $z^{n-1}-t^\varepsilon/n$. The integrality hypothesis is used
twice --- roots are finite, and translating an error by a root does not reduce
its coefficient valuation bound. For arbitrary coefficients, first apply
\cref{polynomial:cor:normalization} to a common root scale and then measure
the rescaled coefficient error; an unnormalized claim with the same
$\varepsilon/n$ at all infinite root locations would discard the conditioning
factors introduced by that change of variables.

\subsection{The rootwise threshold, and full cluster stability}
Let $P\in\OO_K[z]$ be monic of degree $n\ge2$ with distinct finite roots
$\alpha_1,\ldots,\alpha_n$, and recall from
\cref{polynomial:subsec:conventions}
\begin{equation}\label{polynomial:eq:stabilitydata}
 d_{ij}=v(\alpha_i-\alpha_j),\qquad
 d_i=v(P'(\alpha_i))=\sum_{j\ne i}d_{ij},\qquad
 \delta_i=\max_{j\ne i}d_{ij},\qquad
 T(P)=\max_i(d_i+\delta_i).
\end{equation}
All these are nonnegative. Here $\delta_i$ records the closest other root in
valuation geometry --- larger valuation means smaller distance --- while $d_i$
measures the size of the inverse derivative at the root. Distinguishing them
gives a sharper threshold than any single coarse discriminant bound. The open
balls $\Bgt{\alpha_i}{\delta_i}$ are pairwise disjoint, since a point in two of
them would force $v(\alpha_i-\alpha_j)>\min(\delta_i,\delta_j)\ge d_{ij}$.

The following theorem is proved, in the same form and with the same threshold,
by all three merged manuscripts; two of them state it with the letters $T(P)$
and $\tau(P)$ for the same number $\max_i(d_i+\delta_i)$, and with $\sigma_i$
and $d_i$ for the same $v(P'(\alpha_i))$. It is stated once. Its conclusion
list is the union of theirs: the unique labelling and the preserved distance
matrix, the leading-coefficient refinement, the exact displacement valuation,
and the second-order Newton estimate
\eqref{polynomial:eq:newtonerror}, which only one of the three proves.

\begin{theorem}[Exact matching, full cluster stability, and a second-order
Newton error]\label{polynomial:thm:stability}
Let $P\in\OO_K[z]$ be monic of degree $n\ge2$ with distinct finite roots and
data \eqref{polynomial:eq:stabilitydata}, and let $Q=P+E$ be monic of degree
$n$ with $\deg E<n$ and coefficient precision
$\varepsilon=\min_{j<n}v([z^j]E)$, with $\varepsilon=+\infty$ when $E=0$.
Suppose
\begin{equation}\label{polynomial:eq:stabilitythreshold}
 \varepsilon>T(P)=\max_i(d_i+\delta_i).
\end{equation}
Then $Q$ is squarefree with $n$ distinct finite roots, and there is a
\emph{unique} labelling $\beta_1,\ldots,\beta_n$ with
\begin{equation}\label{polynomial:eq:matchingballs}
 v(\beta_i-\alpha_i)>\delta_i\qquad(1\le i\le n),
\end{equation}
that is, exactly one root of $Q$ in each disjoint ball
$\Bgt{\alpha_i}{\delta_i}$, and these are all of them. Moreover
\begin{align}
 v(\beta_i-\beta_j)&=d_{ij}\quad\text{and}\quad
 \lc(\beta_i-\beta_j)=\lc(\alpha_i-\alpha_j) &&(i\ne j),
 \label{polynomial:eq:pairwisepreserved}\\
 v(\beta_i-\alpha_i)&\ge\varepsilon-d_i,\label{polynomial:eq:shiftbound}\\
 v(Q'(\beta_i))&=d_i,\qquad v(\Disc Q)=v(\Disc P).
 \label{polynomial:eq:discpreserved}
\end{align}
If $E(\alpha_i)=0$ then $\beta_i=\alpha_i$. If $E(\alpha_i)\ne0$ then the
displacement has the exact valuation and leading term
\begin{equation}\label{polynomial:eq:leadingdisplacement}
 v(\beta_i-\alpha_i)=v(E(\alpha_i))-d_i,
 \qquad
 \beta_i-\alpha_i=-\frac{E(\alpha_i)}{P'(\alpha_i)}\,(1+\epsilon_i),
 \quad \epsilon_i\in\mm_K,
\end{equation}
and the first Newton correction is a genuine second-order improvement:
\begin{equation}\label{polynomial:eq:newtonerror}
 v\left(\beta_i-\alpha_i+\frac{E(\alpha_i)}{P'(\alpha_i)}\right)
 \ \ge\ 2(\varepsilon-d_i)-\delta_i .
\end{equation}
\end{theorem}
\begin{proof}
Fix $i$ and work at center $\alpha_i$ and scale $\delta_i$. Formula
\eqref{polynomial:eq:profileproduct} gives
\begin{equation}\label{polynomial:eq:rootweight}
 w_{\alpha_i,\delta_i}(P)=\delta_i+\sum_{j\ne i}d_{ij}=\delta_i+d_i,
\end{equation}
since every other separation has $d_{ij}\le\delta_i$; the initial polynomial
therefore has a simple zero at the origin, so exactly one root of $P$ lies in
$\Bgt{\alpha_i}{\delta_i}$. Every coefficient of the translated expansion
$E(\alpha_i+Z)$ has valuation at least $\varepsilon$: use the finite binomial
formula, integrality of $\alpha_i$, and
\eqref{polynomial:eq:integersunit}; since $\delta_i\ge0$, scaling by
$t^{\delta_i}$ cannot lower these valuations. Hence
\[
 w_{\alpha_i,\delta_i}(E)\ \ge\ \varepsilon\ >\ d_i+\delta_i
 =w_{\alpha_i,\delta_i}(P),
\]
and \cref{polynomial:cor:valrouche} gives exactly one root of $Q$, counted
with multiplicity, in $\Bgt{\alpha_i}{\delta_i}$; it is therefore simple.
These $n$ balls are disjoint, their $n$ simple roots exhaust the degree of
$Q$, and finiteness of the roots is immediate from finiteness of the centers
and nonnegativity of the scales. This proves the matching and its uniqueness.

Put $u_i=\beta_i-\alpha_i$. For $i\ne j$ both $u_i,u_j$ have valuation
strictly greater than $d_{ij}$, so $\alpha_i-\alpha_j$ determines both the
valuation and the leading coefficient of $\beta_i-\beta_j$; this is
\eqref{polynomial:eq:pairwisepreserved}. Taking valuations of the finite root
products gives \eqref{polynomial:eq:discpreserved}; alternatively
$Q'(\alpha_i+u_i)=U(\alpha_i+u_i)+u_iU'(\alpha_i+u_i)+E'(\alpha_i+u_i)$ with
$U=\prod_{j\ne i}(z-\alpha_j)$, where $u_iU'/U=\sum_{j\ne i}u_i/(\beta_i-\alpha_j)$
is infinitesimal and the last term has valuation at least
$\varepsilon>d_i$.

For the displacement, all factors $\beta_i-\alpha_j$ with $j\ne i$ have
valuation $d_{ij}$, so $v(P(\beta_i))=v(u_i)+d_i$ when $u_i\ne0$; since
$P(\beta_i)=-E(\beta_i)$ and $\beta_i$ is integral, $v(E(\beta_i))\ge\varepsilon$,
which is \eqref{polynomial:eq:shiftbound} (a zero displacement satisfies it
under the $+\infty$ convention). If $E(\alpha_i)=0$ then $\alpha_i$ itself is a
root of $Q$ in its assigned ball and uniqueness gives $u_i=0$. Otherwise the
exact factorization
\begin{equation}\label{polynomial:eq:productnearroot}
 P(\alpha_i+u_i)=P'(\alpha_i)\,u_i\prod_{j\ne i}
 \left(1+\frac{u_i}{\alpha_i-\alpha_j}\right)
 =P'(\alpha_i)u_i(1+\epsilon_i),\qquad \epsilon_i\in\mm_K,
\end{equation}
has every parenthesized factor a unit with standard part one, because
$v(u_i)>\delta_i\ge d_{ij}$. Also $E(\alpha_i+u_i)-E(\alpha_i)=u_iV_i$ with
$v(V_i)\ge\varepsilon$, by the finite difference-of-powers identity and
integrality of both arguments; since $\varepsilon>d_i$ this difference has
valuation strictly larger than $d_i+v(u_i)$ and cannot affect the leading term
of $P(\alpha_i+u_i)+E(\alpha_i+u_i)=0$. This proves
\eqref{polynomial:eq:leadingdisplacement}.

For \eqref{polynomial:eq:newtonerror} set
$R_P=P(\alpha_i+u_i)-P'(\alpha_i)u_i$ and $R_E=E(\alpha_i+u_i)-E(\alpha_i)$.
Expanding the finite product in \eqref{polynomial:eq:productnearroot} and the
finite Taylor polynomial of $E$ gives
\begin{equation}\label{polynomial:eq:remainderbounds}
 v(R_P)\ge d_i+2v(u_i)-\delta_i,\qquad v(R_E)\ge\varepsilon+v(u_i),
\end{equation}
the first because every term of the product minus one contains at least one
factor $u_i/(\alpha_i-\alpha_j)$ of valuation at least $v(u_i)-\delta_i$, the
second because every translated error coefficient has valuation at least
$\varepsilon$ and $v(u_i)>0$. The root equation is
$P'(\alpha_i)u_i+E(\alpha_i)+R_P+R_E=0$; dividing by $P'(\alpha_i)$ and
inserting \eqref{polynomial:eq:remainderbounds} and then
\eqref{polynomial:eq:shiftbound} gives
\[
 v\left(u_i+\frac{E(\alpha_i)}{P'(\alpha_i)}\right)
 \ \ge\ \min\{2v(u_i)-\delta_i,\ \varepsilon+v(u_i)-d_i\}
 \ \ge\ 2(\varepsilon-d_i)-\delta_i . \qedhere
\]
\end{proof}

The theorem preserves the entire matrix of pairwise valuations, so every root
partition defined by a valuation threshold, and hence the whole cluster tree
of \cref{polynomial:subsec:tree}, is preserved. That is a precise meaning of
preservation of the complete finite cluster geometry. It does \emph{not}
assert equality of the roots, of their ordinary moduli, or of every Hahn
coefficient of the pairwise differences.

When $E\ne0$, the right side of \eqref{polynomial:eq:newtonerror} is strictly greater than
$\varepsilon-d_i$ under \eqref{polynomial:eq:stabilitythreshold}, so
$-E(\alpha_i)/P'(\alpha_i)$ is a genuine first-order correction with a
strictly smaller error, not merely a formal suggestion to apply Newton's
method. No infinite Newton iteration is required for the proof. For degree one
the exact root displacement is simply the negative constant-coefficient error.

\begin{corollary}[Discriminant precision protects all clusters]
\label{polynomial:cor:discprecision}
Under the hypotheses on $P$, put $D=v(\Disc P)$. Then $T(P)\le D$, so a monic
same-degree perturbation with
\begin{equation}\label{polynomial:eq:discriminantprecision}
 \varepsilon>D
\end{equation}
has all the matching and preservation properties of
\cref{polynomial:thm:stability}. One also has $T(P)\le2\max_id_i$.
\end{corollary}
\begin{proof}
By \eqref{polynomial:eq:discval}, $D=2\sum_{i<j}d_{ij}=\sum_i d_i$. Fix $i$
and choose $j_0$ with $d_{ij_0}=\delta_i$; in $d_i+\delta_i$ the pair
$(i,j_0)$ is counted twice, every other pair incident to $i$ once, and the
remaining pairs not at all, and all $d_{ij}$ are nonnegative, so
$d_i+\delta_i\le D$. Take the maximum over $i$. For the last bound,
$\delta_i\le d_i$ since all separations are nonnegative.
\end{proof}

\begin{example}[Why the inequality must be strict]\label{polynomial:ex:sharpdisc}
Let $h>0$ be a value-group element and put
\[
 P(z)=z(z-t^h),\qquad Q(z)=P(z)+\tfrac14t^{2h}=\bigl(z-\tfrac12t^h\bigr)^2 .
\]
The two roots of $P$ are distinct with $d_i=\delta_i=h$, so
$T(P)=D=2h$. The error has valuation exactly $2h$, yet $Q$ has a double root,
which lies in neither separated open ball $\Bgt0h$ nor $\Bgt{t^h}h$. Thus
replacing $>$ by $\ge$ in \eqref{polynomial:eq:stabilitythreshold} or in
\eqref{polynomial:eq:discriminantprecision} is false. The example proves
sharpness of the strict boundary in this family; it does not claim that $D$ is
the optimal threshold for every polynomial, nor optimality of the bound for
every individual perturbation. It also realizes equality in the global
H\"older bound \eqref{polynomial:eq:holder} for $n=2$.
\end{example}

\begin{example}[A complete four-root hierarchy, with $T(P)=7<D=10$]
\label{polynomial:ex:quartic}
Let $t=\omega^{-1}$ and
\begin{equation}\label{polynomial:eq:quartic}
 P(z)=(z-t)(z-t-t^3)(z+t)(z-1),
\end{equation}
with roots labelled in that order. The positive distance valuations are
$d_{12}=3$, $d_{13}=d_{23}=1$, $d_{14}=d_{24}=d_{34}=0$, so
\begin{equation}\label{polynomial:eq:quarticdata}
 (d_1,d_2,d_3,d_4)=(4,4,2,0),\qquad
 (\delta_1,\delta_2,\delta_3,\delta_4)=(3,3,1,0),
\end{equation}
giving $T(P)=7$ whereas $D=10$; the discriminant has leading term $16t^{10}$,
and \eqref{polynomial:eq:clusterdisc} returns the same valuation as
$2\bigl(1\cdot\binom32+(3-1)\binom22\bigr)=10$. The discriminant threshold is
therefore \emph{not} optimal.

The critical points can be placed without solving a cubic. At scale zero about
zero the initial polynomial is $X^3(X-1)$, whose derivative is $X^2(4X-3)$, so
two critical points lie in the infinitesimal monad of zero and one has
standard part $3/4$. At scale one about zero the initial polynomial is, up to
a nonzero scalar, $(X-1)^2(X+1)$, whose derivative is $(X-1)(3X+1)$: one
critical point has $w/t\in1+\mm_K$ and another has $w/t\in-\tfrac13+\mm_K$. In
the ball of scale three about $t$ the initial polynomial is a nonzero scalar
times $X(X-1)$, so the first of these has $(w-t)/t^3\in\tfrac12+\mm_K$. These
three exhaust $\deg P'=3$.

Now take $\varepsilon=8$, between $T(P)$ and $D$. For $Q=P+t^8$,
\cref{polynomial:thm:stability} applies and preserves every pairwise root
distance valuation, with
\begin{align*}
 \beta_1-\alpha_1&=-\tfrac12t^4(1+\epsilon_1), &
 \beta_2-\alpha_2&=\tfrac12t^4(1+\epsilon_2),\\
 \beta_3-\alpha_3&=\tfrac14t^6(1+\epsilon_3), &
 \beta_4-\alpha_4&=-t^8(1+\epsilon_4),
\end{align*}
all $\epsilon_i\in\mm_K$, directly from $-t^8/P'(\alpha_i)$. For
$Q=P+t^8(1+z)$ the four matched displacement valuations are exactly
$4,4,6,8$ by \eqref{polynomial:eq:leadingdisplacement}. The original cluster
scales $0,1,3$ and this quartic's residual critical directions
$3/4,-1/3,1/2$ survive such a perturbation. The analogous statement for
the different quartic of \cref{polynomial:ex:nested} preserves its own
directions $3/4,2/3,1/2$.
\end{example}

\begin{corollary}[Preservation of the tree and of its critical directions]
\label{polynomial:cor:treestability}
Under \cref{polynomial:thm:stability}, at every internal node of the original
root-cluster tree, using its original center and scale, $P$ and $Q$ have the
same initial polynomial. Their critical multiplicities in every residue
direction of that node therefore agree as well.
\end{corollary}
\begin{proof}
A node may be centered at one of its roots $\alpha_i$, and its depth $\rho$ is
at most $\delta_i$; all roots are finite so $\rho\ge0$. By
\eqref{polynomial:eq:profileproduct},
\[
 w_{\alpha_i,\rho}(P)=\rho+\sum_{j\ne i}\min\{\rho,d_{ij}\}
 \le\delta_i+d_i<\varepsilon\le w_{\alpha_i,\rho}(E),
\]
so \cref{polynomial:cor:valrouche} gives equality of the initial polynomials.
At an internal node this polynomial is nonconstant, so
\cref{polynomial:lem:initialderivative} gives equality for the derivatives
too, and the direction counts follow from
\cref{polynomial:thm:initialroots,polynomial:thm:branch}.
\end{proof}

This preserves the root-scale hierarchy and the \emph{allocated} critical
directions, not the individual multiplicities of all critical points. For
instance $z^3-1$ is squarefree but its derivative has a double root at zero;
adding $t^Nz$ with $N>0$ leaves the ordinary root clusters unchanged and
splits that critical point into two at a smaller scale, while the direction
count at the original scale is still two.

\subsection{Using the theorem at infinite scales}\label{polynomial:subsec:rescaled}
Finiteness of the roots in \cref{polynomial:thm:stability} is a normalization
hypothesis, not a restriction to an Archimedean problem. Let $P$ be any monic
squarefree degree-$n$ polynomial over $K$. Choose $a\in\SC$ and a positive
surreal $\rho_0$ large enough that all $(\alpha_i-a)/\rho_0$ are finite, which
\cref{polynomial:prop:rootbounds} supplies from the coefficients, and put
\[
 P_*(Y)=\rho_0^{-n}P(a+\rho_0Y),\qquad Q_*(Y)=\rho_0^{-n}Q(a+\rho_0Y),
\]
both monic of degree $n$. If $E(a+Z)=\sum_{j<n}\widetilde e_jZ^j$, the
normalized precision is
\begin{equation}\label{polynomial:eq:rescaledmu}
 \varepsilon_*=\min_{j<n}\bigl(v(\widetilde e_j)+(j-n)v(\rho_0)\bigr),
\end{equation}
and the discriminant transforms exactly as
\begin{equation}\label{polynomial:eq:rescaleddisc}
 \Disc(P_*)=\rho_0^{-n(n-1)}\Disc(P),\qquad
 D_*=v(\Disc P)-n(n-1)v(\rho_0).
\end{equation}
Thus $\varepsilon_*>D_*$ is a coefficient-level sufficient condition for the
normalized matching, and pulling the roots back by $a+\rho_0Y$ preserves all
pairwise valuation equalities in the original coordinates. Applying the
unnormalized discriminant criterion without this adjustment would be
unjustified when some original roots are infinite.

\section{Universal residue duality and collision-stable traces}
\label{polynomial:sec:residue}
This section deliberately works over a more general coefficient ring than $K$.
Let $R$ be \emph{any} commutative ring with identity and
\[
 P(z)=z^n+p_{n-1}z^{n-1}+\cdots+p_0\in R[z],\qquad n\ge1 .
\]
Monic division works over $R$ without dividing by any coefficient, so
$A=R[z]/(P)$ is free with basis $1,z,\ldots,z^{n-1}$, and we define
\begin{equation}\label{polynomial:eq:lambdadef}
 \lambda_P:A\longrightarrow R,\qquad
 \lambda_P(h)=[z^{n-1}]\rem_P h ,
\end{equation}
the coefficient of $z^{n-1}$ in the remainder of $h$ on division by $P$. For
$h\in A$ write $m_h$ for its multiplication endomorphism. No contour is
chosen and no moving root is labelled. In particular, if $P,h\in\OO_K[z]$ then
$\rem_Ph$ and $\lambda_P(h)$ are integral as well, since monic division uses
no inverse coefficients. The functional detects nilpotent directions that the
trace does not, and it is defined by finite polynomial division, even at
repeated roots and at infinite surcomplex roots.

\subsection{A dual basis with no discriminant denominator}
\begin{theorem}[Universal residue pairing and explicit dual basis]
\label{polynomial:thm:residuepairing}
Over any commutative ring $R$ with identity, the bilinear form
$(f,g)\mapsto\lambda_P(fg)$ on $A=R[z]/(P)$ is perfect, including when $P$ has
repeated roots. With $p_n=1$, the residue-dual basis to $1,z,\ldots,z^{n-1}$
is
\begin{equation}\label{polynomial:eq:dualbasis}
 b^i(z)=\sum_{k=i+1}^np_kz^{k-1-i},\qquad 0\le i<n,
\end{equation}
so that $\lambda_P(z^jb^i)=\delta_{ij}$. Its Gram matrix
$G_{ij}=\lambda_P(z^{i+j})$ has
\begin{equation}\label{polynomial:eq:gramdet}
 \det G=(-1)^{n(n-1)/2},
\end{equation}
and the coefficient matrix of the B\'ezoutian
\begin{equation}\label{polynomial:eq:bezoutkernel}
 \Delta_P(z,Y)=\frac{P(z)-P(Y)}{z-Y}=\sum_{i=0}^{n-1}b^i(z)Y^i
\end{equation}
is $G^{-1}$ in these power bases. The same pairing is perfect over $\OO_K$ on
$\OO_K[z]/(P)$ when $P$ is integral and monic.
\end{theorem}
\begin{proof}
There is a formal Laurent expansion at infinity,
$P(z)^{-1}=z^{-n}(1+p_{n-1}z^{-1}+\cdots+p_0z^{-n})^{-1}$, the inverse being a
formal power series in the independent symbol $z^{-1}$, valid over any $R$.
Division by $P$ gives
\begin{equation}\label{polynomial:eq:lambdaatinfinity}
 \lambda_P(h)=[z^{-1}]\,\frac{h(z)}{P(z)} ,
\end{equation}
since the quotient is polynomial and has no $z^{-1}$ term, and for a remainder
of degree less than $n$ that coefficient is exactly its $z^{n-1}$ coefficient.
The definition of $b^i$ yields
\[
 \frac{b^i(z)}{P(z)}=z^{-i-1}-\frac{\sum_{k=0}^ip_kz^{k-i-1}}{P(z)},
\]
whose second term has highest possible Laurent degree $-n-1$; after
multiplication by $z^j$ with $0\le j<n$ it cannot contribute to degree $-1$.
So \eqref{polynomial:eq:lambdaatinfinity} gives
$\lambda_P(z^jb^i)=\delta_{ij}$, constructing an inverse to the map
$A\to\Hom_R(A,R)$ defined by the pairing and hence proving perfection over $R$
itself.

For the determinant, $G_{ij}=0$ when $i+j<n-1$ and $G_{ij}=1$ when
$i+j=n-1$; reversing the column order gives a lower triangular matrix with
diagonal ones, and the column-reversal sign is $(-1)^{n(n-1)/2}$. Finally the
finite identity $(z^k-Y^k)/(z-Y)=\sum_{i<k}z^{k-1-i}Y^i$ proves
\eqref{polynomial:eq:bezoutkernel}, whose $i$th coefficient column is the dual
vector $b^i$, so its coefficient matrix is $G^{-1}$.
\end{proof}

No inverse root separation, resultant or discriminant appears. This is
stronger than nondegeneracy over a fraction field: it is an isomorphism over
the given ring $R$, even when $R$ has zero divisors or is non-Noetherian. Its
strength for the present problem is that no discriminant denominator ever
appears, which is what makes the pairing survive every collision. This is the
polynomial core of the one-variable residue duality in
\cite{polynomial:SourceResearch}; it is not claimed as a new general duality
theorem.

\subsection{Trace, norm, and the derivative as the trace element}
\begin{theorem}[Universal trace identity]\label{polynomial:thm:trace}
Over any commutative ring $R$ with identity and monic $P$ of degree $n\ge1$,
\begin{equation}\label{polynomial:eq:traceidentity}
 \Tr(m_h)=\lambda_P(P'h)\qquad(h\in A).
\end{equation}
If $G_{\mathrm{tr},ij}=\Tr(m_{z^{i+j}})$ then
\begin{equation}\label{polynomial:eq:tracegram}
 G_{\mathrm{tr}}=G\,M_{P'},\qquad \det G_{\mathrm{tr}}=\Disc(P),
\end{equation}
where $M_{P'}$ is the multiplication matrix of $P'$ and $\Disc$ is the usual
universal coefficient polynomial for monic $P$. For any polynomial $Q$,
$\det(m_Q)=\Res(P,Q)$.
\end{theorem}
\begin{proof}
The coefficient of $z^i$ in the power-basis expansion of $hz^i$ is
$\lambda_P(hz^ib^i)$, so summing the diagonal coefficients gives
$\Tr(m_h)=\lambda_P\bigl(h\sum_{i<n}z^ib^i\bigr)$. The explicit dual basis
satisfies the polynomial identity
\begin{equation}\label{polynomial:eq:eulerderivative}
 \sum_{i=0}^{n-1}z^ib^i=\sum_{k=1}^n\sum_{i=0}^{k-1}p_kz^{k-1}
 =\sum_{k=1}^nkp_kz^{k-1}=P' ,
\end{equation}
proving \eqref{polynomial:eq:traceidentity} over every $R$. The
trace-pairing entry is $\lambda_P(z^iP'z^j)$, the $(i,j)$ entry of $GM_{P'}$.
For $\det(m_Q)=\Res(P,Q)$, both sides are polynomials with integer
coefficients in the finitely many coefficients of $P,Q$; over the fraction
field of the universal integer coefficient ring $P$ has distinct roots and the
multiplication determinant is $\prod_iQ(\alpha_i)$, which is
\eqref{polynomial:eq:resultantproduct}. The universal coefficient ring injects
into that field, so the identity holds there and after every specialization.
Taking determinants in \eqref{polynomial:eq:tracegram} gives
$(-1)^{n(n-1)/2}\Res(P,P')=\Disc(P)$ for monic $P$, and the same
universal-polynomial argument identifies it with
\eqref{polynomial:eq:discdef} over splitting fields.
\end{proof}

The element representing the trace under residue duality is exactly $P'$. This
should not be mistaken for a proof of the general multivariable Jacobian
comparison; the argument here concerns the monogenic algebra $R[z]/(P)$. The
several-variable statement is \cref{polynomial:thm:jacobian}, proved
separately.

\begin{remark}[The B\'ezoutian is the duality kernel]
\label{polynomial:rem:kernelmap}
There is a second, basis-free proof of
\eqref{polynomial:eq:traceidentity} that will be reused verbatim in $n$
variables. The element $\Delta_P\in A\otimes_RA$ of
\eqref{polynomial:eq:bezoutkernel} satisfies $(z-Y)\Delta_P=0$ and
$(\lambda_P\otimes\id)(\Delta_P)=1$, the latter because its coefficient of
$z^{n-1}$ is one and its $z$-degree is $n-1$. A perfect pairing identifies
$A\otimes_RA$ with $\operatorname{End}_R(A)$ by
\begin{equation}\label{polynomial:eq:kernelmap}
 \sum_su_s\otimes v_s\ \longmapsto\ \Bigl(h\mapsto\sum_s\lambda_P(hu_s)v_s\Bigr).
\end{equation}
The first identity says the operator attached to $\Delta_P$ commutes with
multiplication by $z$, hence with multiplication by every polynomial; the
second says it sends $1$ to $1$. It is therefore the identity operator, so for
residue-dual bases $(b_j),(b^j)$ one gets
$\Delta_P=\sum_jb_j(z)b^j(Y)$ and $P'=\sum_jb_jb^j$ in $A$, and taking
diagonal coefficients of multiplication by $h$ recovers
$\Tr(m_h)=\lambda_P(P'h)$. The derivative is the diagonal value of the duality
kernel.
\end{remark}

\subsection{Residues at distinct or colliding roots}
For a rational function, the residue at $\xi$ means the coefficient of
$(z-\xi)^{-1}$ in its formal Laurent expansion at $\xi$; the residue at
infinity is taken in the coordinate $u=1/z$ for the differential, hence is the
negative of the coefficient of $z^{-1}$ in the expansion at infinity. These
are formal Laurent coefficients, never integrals.

\begin{proposition}[Finite residue formula at multiple roots]
\label{polynomial:prop:localresidues}
Let $P$ be monic over a characteristic-zero algebraically closed field, in
particular over $K$, with $P=\prod_i(z-\alpha_i)^{m_i}$ and
$P_i=P/(z-\alpha_i)^{m_i}$. Then for every polynomial $h$,
\begin{equation}\label{polynomial:eq:localresidues}
 \lambda_P(h)=\sum_i\Res_{z=\alpha_i}\frac{h(z)}{P(z)}\,dz
 =\sum_i\frac1{(m_i-1)!}\left(\frac{d^{m_i-1}}{dz^{m_i-1}}
 \frac{h(z)}{P_i(z)}\right)_{z=\alpha_i} ,
\end{equation}
and in the squarefree case
\begin{equation}\label{polynomial:eq:simpleresidues}
 \lambda_P(h)=\sum_{i=1}^n\frac{h(\alpha_i)}{P'(\alpha_i)} .
\end{equation}
\end{proposition}
\begin{proof}
Finite partial fractions express $h/P$ as a polynomial plus a sum of terms
$c_{i,k}(z-\alpha_i)^{-k}$. In the Laurent expansion at infinity only the
terms with $k=1$ contribute to $z^{-1}$, with contributions $c_{i,1}$; by
\eqref{polynomial:eq:lambdaatinfinity} this proves the first equality. Near
$\alpha_i$, with $u=z-\alpha_i$, the coefficient $c_{i,1}$ is the coefficient
of $u^{m_i-1}$ in the Taylor expansion of $h(\alpha_i+u)/P_i(\alpha_i+u)$,
whose denominator has nonzero constant term; that coefficient is exactly the
displayed derivative. For a simple root $P_i(\alpha_i)=P'(\alpha_i)$.
\end{proof}

Both formulas are finite algebraic identities for surcomplex inputs. The local
terms in \eqref{polynomial:eq:simpleresidues} may individually be infinite
even when $\lambda_P(h)$ is finite or zero: this is not divergence but
cancellation in a finite sum. For $P=z(z-t)$ one has
$\lambda_P(h)=(h(t)-h(0))/t$, with infinite individual weights but with
$\lambda_P(h)$ integral for $h\in\OO_K[z]$, by
\eqref{polynomial:eq:lambdadef}. At the repeated polynomial $P=(z-t)^2$ one
has $\lambda_P(h)=h'(t)$ --- an exact formula involving a jet, not a limit
taken as two roots approach each other. At a collision the separate point
weights become singular while the universal functional remains defined and its
local formula retains derivative information. The same partial-fraction
argument gives the algebraic residue theorem for every rational differential
on the projective line: the sum of its finitely many finite residues and its
residue at infinity is zero. Again, no topological loop in $\SC$ is required.

\begin{proposition}[Field-level traces and norms, including collisions]
\label{polynomial:prop:tracenormroots}
Over an algebraically closed characteristic-zero field, with
$P=\prod_i(z-\alpha_i)^{m_i}$,
\begin{equation}\label{polynomial:eq:tracenormroots}
 \Tr(m_h)=\sum_im_ih(\alpha_i),\qquad
 \Nm_A(h):=\det(m_h)=\prod_ih(\alpha_i)^{m_i} .
\end{equation}
The radical of the trace pairing is the nilradical of $A$, whereas the residue
pairing of \cref{polynomial:thm:residuepairing} has zero radical, whether or
not roots repeat.
\end{proposition}
\begin{proof}
Use the CRT factors $K[u]/(u^{m_i})$ of \cref{polynomial:thm:crt}. On the
$i$th factor multiplication by $h(\alpha_i+u)$ is triangular in the powers of
$u$ with every diagonal entry $h(\alpha_i)$, so its trace is $m_ih(\alpha_i)$
and its determinant $h(\alpha_i)^{m_i}$; summing and multiplying gives
\eqref{polynomial:eq:tracenormroots}. On a pair of elements the trace pairing
is $\sum_im_if(\alpha_i)g(\alpha_i)$; if all $f(\alpha_i)=0$ this vanishes for
every $g$, and conversely the CRT idempotent of the $i$th factor gives
$m_if(\alpha_i)=0$, whence $f(\alpha_i)=0$ in characteristic zero. These are
precisely the elements whose image in each local factor is nilpotent. The
residue pairing is perfect by its constructed dual basis.
\end{proof}

\begin{theorem}[Discriminant and trace-pairing degeneration]
\label{polynomial:thm:discriminant}
For monic $P$ of degree $n$ over $K$ with $T_{ij}=\Tr(m_{z^{i+j}})$,
\begin{equation}\label{polynomial:eq:tracedisc}
 \det T=\Disc(P)=(-1)^{n(n-1)/2}\Res(P,P') .
\end{equation}
The trace pairing is nondegenerate exactly when $P$ is squarefree, whereas the
residue pairing of \cref{polynomial:thm:residuepairing} is always
nondegenerate, with the constant Gram determinant
\eqref{polynomial:eq:gramdet}.
\end{theorem}
\begin{proof}
List the $n$ roots with repetition and let $V_{ki}=\alpha_k^i$. The trace
formula gives $T=V^TV$ (transpose, not conjugate transpose), so
$\det T=(\det V)^2=\prod_{i<j}(\alpha_i-\alpha_j)^2$, an identity that remains
valid with repeated rows and so handles repeated roots directly. Also
$\prod_iP'(\alpha_i)=(-1)^{n(n-1)/2}\prod_{i<j}(\alpha_i-\alpha_j)^2$: in the
squarefree case multiply $P'(\alpha_i)=\prod_{j\ne i}(\alpha_i-\alpha_j)$, and
with repeated roots both sides vanish. Apply
\eqref{polynomial:eq:resultantproduct}. The product of squared differences is
nonzero precisely for distinct roots, while the residue determinant is the
fixed unit \eqref{polynomial:eq:gramdet}.
\end{proof}

This determinant formula shows that the product definition
\eqref{polynomial:eq:discdef} is a coefficient polynomial, and links its
valuation to the root-separation precision of
\cref{polynomial:sec:stability}. The trace cannot detect nilpotent directions
in a local factor, whereas the residue functional detects the higher jets
needed for a perfect pairing; these are different bilinear forms and must not
be interchanged. Compare also \cref{polynomial:thm:hermite}, which is a third
form again, built from the trace and used for real-root counting.

\begin{example}[A cubic collision]\label{polynomial:ex:cubicduality}
This calculation appears in the supplied several-variable manuscript
\cite{polynomial:SourceResearch}; it is included as a normalization check and
a bridge to that work, not as a new example. Let $P=z^3-az-b$. In the basis
$(1,z,z^2)$,
\[
 G=\begin{pmatrix}0&0&1\\0&1&0\\1&0&a\end{pmatrix},\qquad
 G^{-1}=\begin{pmatrix}-a&0&1\\0&1&0\\1&0&0\end{pmatrix},\qquad
 G_{\mathrm{tr}}=\begin{pmatrix}3&0&2a\\0&2a&3b\\2a&3b&2a^2\end{pmatrix},
\]
with $\det G=-1$ and $\det G_{\mathrm{tr}}=4a^3-27b^2$. The residue-dual basis
is $(z^2-a,\,z,\,1)$, whose trace element is $(z^2-a)+z^2+z^2=3z^2-a=P'$. At
$a=3t^2$, $b=2t^3$ the polynomial is $(z-2t)(z+t)^2$ and
\begin{equation}\label{polynomial:eq:cubiccollision}
 \lambda_P(h)=\frac{h(2t)-h(-t)}{9t^2}-\frac{h'(-t)}{3t},
 \qquad \Tr(m_h)=h(2t)+2h(-t).
\end{equation}
The trace forgets the nilpotent direction at the double point; the residue
functional remembers its first derivative.
\end{example}

\subsection{Coherent coefficient families on a fixed ordinary domain}
\label{polynomial:subsec:coherentfamilies}
Take $R=\HH(U)$ or $R=\HHn(U)$ for one fixed ordinary domain $U$. Every monic
bounded-degree family $P(x,z)$ gives a free rank-$n$ algebra over $R$.
Polynomial division, multiplication matrices, $\lambda_P$, the dual basis
\eqref{polynomial:eq:dualbasis}, traces, norms and discriminants are computed
by finitely many coefficient-ring operations, hence are legitimate
Hahn-coherent coefficient functions on the \emph{same} $U$. No inverse
discriminant and no repeated domain shrink is needed for the residue pairing.
For an integral family all entries of the residue Gram matrix and of its
explicit inverse are integral, and if the input coefficient supports lie in a
nonnegative support monoid $M$, so do these entries. At a finite surcomplex
parameter, evaluation \eqref{polynomial:eq:coherentevaluation} is a ring
homomorphism, so it commutes with these finite identities and the specialized
pairing is exactly the pairing of the specialized polynomial. This explains
collision stability algebraically, independently of any simultaneous choice of
root labels.

For more general coherent analytic numerators of unbounded Taylor degree one
must first use the fixed-domain division theorems of the supplied manuscripts
before taking a remainder. The purely polynomial argument of this section does
not silently replace that extra analytic step, and nothing here is asserted
over a radius-free or formal-coefficient ring.

\section{Several variables: certificates, finite algebras, and local
multiplicity}\label{polynomial:sec:multivariate}
The one-variable results fit into ordinary algebraic geometry over $K$. We
make the legitimate transfer explicit, because the proper-class notation
$\SC$ by itself is not a substitute for a set-sized ring in a scheme
construction; recall \cref{polynomial:rem:classsafety}.

\subsection{The finite-type field lemma}
\begin{lemma}[Zariski's lemma in the needed form]\label{polynomial:lem:zariski}
If an extension field $L$ of an infinite field $k$ is finitely generated as a
$k$-algebra, then it is finite algebraic over $k$. In particular, if $k$ is
algebraically closed then $L=k$.
\end{lemma}
\begin{proof}
Choose a maximal algebraically independent subcollection $t_1,\ldots,t_r$
among a finite algebra generating set; the other finitely many generators are
algebraic over $k(t_1,\ldots,t_r)$, so $L$ is a finite field extension of that
rational-function field. Clearing the denominators of the coefficients in
their monic algebraic equations supplies a nonzero
$f\in k[t_1,\ldots,t_r]$ such that every remaining generator is integral over
$A=k[t_1,\ldots,t_r,1/f]$. Because $L$ is generated as an algebra by the
original generators and is a field containing $1/f$, it equals the algebra
obtained by adjoining those integral generators to $A$; thus $L$ is integral
over $A$.

If a field is integral over a subring $A$, that subring is a field: for
$0\ne a\in A$, a monic equation for $a^{-1}\in L$, multiplied by a suitable
power of $a$, expresses $a^{-1}$ as an element of $A$. If $r\ge1$, choose
$c\in k$ such that $t_1-c$ does not divide $f$; only finitely many linear
factors divide a fixed nonzero polynomial, and $k$ is infinite. Then
$1/(t_1-c)\notin A$, as one checks by multiplying a putative representation by
a power of $f$ and using divisibility in $k[t_1,\ldots,t_r]$. This contradicts
that $A$ is a field. Hence $r=0$, and the original generators are algebraic.
\end{proof}

\subsection{A class-safe Nullstellensatz for finite data}
\begin{theorem}[Surcomplex Nullstellensatz]\label{polynomial:thm:nullstellensatz}
Let $K$ be a set-sized algebraically closed Hahn workspace. A proper ideal
$I\subset K[x_1,\ldots,x_d]$ has a common zero in $K^d$, and $H\in K[x]$
vanishes at all common zeros of $I$ if and only if $H^N\in I$ for some
positive integer $N$.

In class-safe form for finite data: let
$F_1,\ldots,F_s,G\in\SC[x_1,\ldots,x_d]$ be finitely many finite-degree
polynomials. Then $G$ vanishes at every common surcomplex zero of the $F_i$ if
and only if, for some ordinary integer $N\ge1$ and finite polynomials
$H_i\in\SC[x_1,\ldots,x_d]$,
\begin{equation}\label{polynomial:eq:nullcertificate}
 G^N=\sum_{i=1}^sH_iF_i .
\end{equation}
In particular the $F_i$ have no common surcomplex zero if and only if $1$
belongs to their polynomial ideal, and the certificate may be chosen inside
any algebraically closed Hahn workspace containing all the coefficients.
\end{theorem}
\begin{proof}
Fix such a workspace $K$. A maximal ideal $\mathfrak n$ of the \emph{set} ring
$K[x]$ has field quotient finitely generated as a $K$-algebra, hence equal to
$K$ by \cref{polynomial:lem:zariski}; the images of the variables give
coordinates $a_1,\ldots,a_d\in K$ and $\mathfrak n=(x_1-a_1,\ldots,x_d-a_d)$.
A proper ideal lies in such a maximal ideal and so has a common zero.
Conversely, evaluation at a common zero is a homomorphism annihilating $I$, so
$1\notin I$.

If $H$ vanishes at all common zeros, the ideal
$J=I K[x,y]+(1-yH)$ has no common zero, so the weak form gives $1\in J$.
Membership in this generated ideal uses only finitely many $F_i\in I$ and
therefore gives an identity
$1=\sum_iA_i(x,y)F_i(x)+B(x,y)(1-yH(x))$; no finite generating list for $I$
has been assumed. If $H=0$ the certificate is
immediate; otherwise substitute $y=H^{-1}$ in the localization and multiply by
a sufficiently high power $H^N$ to clear denominators. This is the usual
Rabinowitsch argument, the underlying classical theorem being recorded for
instance in \cite[Tag 00FV]{polynomial:Stacks}.

For the class-safe form, the hypothesis about surcomplex zeros implies the
hypothesis about $K$-zeros, so it supplies the certificate in $K$; conversely
that same finite identity persists in every larger workspace and hence at
every surcomplex tuple. No proper-class maximal-ideal argument is used.
\end{proof}

No priority claim attaches to this classical theorem. The point is that the
field and ideal to which it is applied are specified and set-sized. A
positive-dimensional solution set can acquire new points in a larger
workspace, even though solvability does not change. The statement says nothing
about an effective coefficient oracle for an arbitrary surreal constant: once
coefficients belong to a computably presented field, ordinary finite
algebraic algorithms may be available, but an arbitrary Hahn support need not
admit such a presentation.

\begin{remark}[A non-Noetherian witness]\label{polynomial:rem:nonnoetherian}
We have not asserted Noetherianity of any of the analytic coefficient rings in
the supplied manuscripts, and already over a point it fails. For nonzero
divisible $\Gamma$, the valuation ring $\OO_K$ has the non-finitely generated
ideal $\mm_K$: a finite set of positive generator valuations has a least
member $\delta>0$, but $t^{\delta/2}$ is not in the ideal they generate over
$\OO_K$. The positive part is not an ideal of the full Hahn field, where
positive monomials are units. This distinction is essential, and it is exactly
why \cref{polynomial:thm:residuepairing,polynomial:thm:globalfree} are proved
over an arbitrary commutative ring with no finiteness hypothesis.
\end{remark}

\begin{proposition}[Finite fibers and local lengths]\label{polynomial:prop:finitealgebra}
If $B=K[x_1,\ldots,x_d]/I$ is finite-dimensional of dimension $N>0$, then its
common-zero set $Z\subset K^d$ is finite and
\begin{equation}\label{polynomial:eq:finitealgebraproduct}
 B\cong\prod_{\xi\in Z}B_\xi,\qquad \sum_{\xi\in Z}\dim_KB_\xi=N,
\end{equation}
each $B_\xi$ being local with residue field $K$ and nilpotent maximal ideal.
Thus $|Z|\le N$, with equality precisely when all local factors have dimension
one. After extension to a larger algebraically closed workspace there are no
new point locations and the local dimensions are unchanged. Moreover
\begin{equation}\label{polynomial:eq:multivartrace}
 \Tr(m_H)=\sum_{\xi\in Z}(\dim_KB_\xi)H(\xi),\qquad
 \det(m_H)=\prod_{\xi\in Z}H(\xi)^{\dim_KB_\xi}.
\end{equation}
\end{proposition}
\begin{proof}
Distinct maximal ideals are comaximal, so by the finite Chinese remainder
theorem any $s$ of them give a surjection $B\to K^s$, using
\cref{polynomial:lem:zariski} for their residue fields; hence there are at
most $N$ of them, and they correspond to $Z$ by
\cref{polynomial:thm:nullstellensatz}. Let $J$ be their intersection. The
chain $J^m$ stabilizes by finite dimensionality; for a stable
$M=J^m=JM$, choose finitely many generators, write each as a $J$-combination
of them, and multiply the resulting matrix equation by its adjugate: the
determinant is $1+j$ with $j\in J$, hence a unit annihilating all generators,
so $M=0$. The powers $\mathfrak n_\xi^m$ are pairwise comaximal with
intersection $J^m=0$, and the Chinese remainder theorem gives
\eqref{polynomial:eq:finitealgebraproduct}; each factor has one maximal ideal,
nilpotent, with residue field $K$. After scalar extension each factor still
has a nilpotent ideal and the extension field as quotient, so it is still
local of the same dimension, and the coordinates modulo that ideal are the
original tuple $\xi$. Finally, multiplication by $H-H(\xi)$ on each local
factor is nilpotent, which gives \eqref{polynomial:eq:multivartrace}.
\end{proof}

The integer $\dim_KB_\xi$ is the algebraic local multiplicity, a
zero-dimensional intersection multiplicity. For an arbitrary analytic system,
constructing such a finite algebra is additional work and is not supplied by
this proposition. A monic Gr\"obner basis computed over $K$ remains one after
extension to another field, since its division and critical-pair identities
persist, so its standard monomials and quotient dimensions are unchanged; this
is how elimination may be carried out in a fixed workspace, and it claims no
algorithm for arbitrary surreal coefficients lacking finite descriptions.

\begin{example}[A coupled infinitesimal four-point fiber]\label{polynomial:ex:coupled}
Let $t=\omega^{-1}$ and $f_1=x^2-ty-t^2$, $f_2=y^2-tx$. Since $t$ is a nonzero
field element, set $Y=y/t$; the second equation gives $x=tY^2$ and the first
becomes $t^2(Y^4-Y-1)=0$, so
\begin{equation}\label{polynomial:eq:couplediso}
 K[x,y]/(f_1,f_2)\cong K[Y]/(Y^4-Y-1),
\end{equation}
of dimension four, with $(1,x,y,xy)\mapsto(1,tY^2,tY,t^2Y^3)$. The ordinary
polynomial $Y^4-Y-1$ has four distinct roots: a common root with $4Y^3-1$
would satisfy $Y^3=1/4$ and $Y/4-Y-1=0$, hence $Y=-4/3$, whose cube is not
$1/4$. Its roots $c_1,\ldots,c_4\in\C$ are nonzero, so the actual solutions
are $(x,y)=(tc_j^2,tc_j)$, all infinitesimal and simple, the Jacobian
$4xy-t^2=t^2(4c_j^3-1)$ being nonzero at each. In the stated basis
\[
 W_x=\begin{pmatrix}0&t^2&0&0\\1&0&0&t^2\\0&t&0&t^2\\0&0&1&0\end{pmatrix},
 \qquad
 W_y=\begin{pmatrix}0&0&0&t^3\\0&0&t&0\\1&0&0&t^2\\0&1&0&0\end{pmatrix},
\]
which commute and satisfy $W_x^2=tW_y+t^2I$, $W_y^2=tW_x$; eliminating $x$
gives the resultant $y^4-t^3y-t^4$. A coupled system whose ordinary reduction
is $x^2=y^2=0$ thus splits into four actual surcomplex points while its
finite-algebra length stays four.
\end{example}

\section{Degree-controlled complete intersections over arbitrary rings}
\label{polynomial:sec:globalci}
Only one of the three merged manuscripts leaves one variable, and this section
and the next are its contribution. Let $A$ be a nonzero commutative ring with
identity and
\begin{equation}\label{polynomial:eq:polysystem}
 F_i(x)=x_i^{d_i}+E_i(x)\in A[x_1,\ldots,x_n],\qquad d_i\ge1,\qquad
 \deg E_i<d_i\quad(1\le i\le n),
\end{equation}
where throughout this section and the next, degree means \emph{total} degree.
Put
\[
 \mathcal B=\{\alpha\in\N^n:0\le\alpha_i<d_i\},\quad
 D=\prod_id_i,\quad \rho_\sharp=\sum_i(d_i-1),\quad
 x^{\mathrm{top}}=\prod_ix_i^{d_i-1}.
\]
The errors may be fully coupled and their coefficients may have arbitrary size
or valuation: the hypothesis is on polynomial degree, not on smallness. The
degree hypothesis is a genuine restriction and not a disguised assertion about
all isolated complete intersections.

\begin{theorem}[Global finite-free normal forms]\label{polynomial:thm:globalfree}
For \eqref{polynomial:eq:polysystem}, the quotient
$B=A[x]/(F_1,\ldots,F_n)$ is free over $A$ with basis
$(x^\alpha)_{\alpha\in\mathcal B}$ and rank $D$. Every polynomial $H$ has a
unique normal form in their span and a finite division identity
\begin{equation}\label{polynomial:eq:globalnf}
 H=\sum_iF_iQ_i+\NF_F(H).
\end{equation}
Reduction never increases total degree, and every actual use of a relation
strictly lowers the degree of the replaced monomial. All output coefficients
are polynomial expressions with integer coefficients in the input
coefficients, so all of this commutes with every ring homomorphism on
coefficients. The sequence $F_1,\ldots,F_n$ is regular in any order.
\end{theorem}
\begin{proof}
Use the rules $x_i^{d_i}\mapsto-E_i$; every replacement has smaller total
degree, so all branches of reduction terminate after finitely many steps on a
monomial, and by linearity on every polynomial. The monomials on which no rule
applies are exactly $x^\alpha$, $\alpha\in\mathcal B$.

For compatibility, without assuming coefficients form a field: an ambiguity
involving the rules for $i\ne j$ occurs on a multiple of
$x_i^{d_i}x_j^{d_j}$; reducing the two factors in one order gives a multiple
of $E_iE_j$ and in the other a multiple of $E_jE_i$, the same polynomial, and
any remaining reductions take place in smaller total degrees. Induction on
total degree resolves these ambiguities, and disjoint monomial reductions
commute by linearity. Equivalently these are monic leading terms with
relatively prime leading monomials whose critical-pair differences reduce to
zero. No division by a coefficient occurs, so the completely reduced
polynomial is independent of the reduction choices.

Each rule subtracts a multiple of one $F_i$, giving
\eqref{polynomial:eq:globalnf}; conversely $x_i^{d_i}Q$ and $-E_iQ$ have the
same normal form, so the normal-form map annihilates $F_iQ$ and hence the
ideal, and an already reduced polynomial in the ideal is zero. This proves
both spanning and independence. The coefficients are obtained by finitely many
additions, subtractions and multiplications in $A$, hence commute with all
coefficient homomorphisms.

For regularity, apply the same reduction to any initial subset of the
equations; its quotient has standard monomials with the corresponding
exponents bounded and the remaining exponents unrestricted, and because
reduction strictly lowers degree its total-degree associated graded ring is
$A[x]/(x_i^{d_i}:i\text{ in that subset})$. Multiplication by the highest part
$x_j^{d_j}$ of the next equation is injective there, merely shifting the still
unrestricted exponent of $x_j$; so the highest nonzero graded term of a
nonzero element survives multiplication by $F_j$, and $F_j$ is a
non-zero-divisor in the preceding quotient. The final quotient is nonzero,
being free of positive rank over nonzero $A$. The argument works in every
ordering.
\end{proof}

\begin{corollary}[A global B\'ezout count, including infinite-scale roots]
\label{polynomial:cor:globalbezout}
Over an algebraically closed Hahn field $K$, the system
\eqref{polynomial:eq:polysystem} has total local algebraic multiplicity $D$ in
all of $\SC^n$. If its Jacobian determinant is nonzero at every solution,
there are exactly $D$ distinct solutions. The homogenized system has no
additional projective solutions at infinity.
\end{corollary}
\begin{proof}
Use the rank-$D$ quotient and
\eqref{polynomial:eq:finitealgebraproduct}; workspace independence
(\cref{polynomial:prop:workspace}) includes every surcomplex solution. At a
solution with invertible Jacobian the images of the equations generate the
cotangent space of the coordinate maximal ideal, so Nakayama's lemma in the
finite local algebra makes its maximal ideal zero and that factor has
dimension one. At the projective hyperplane $x_0=0$ the homogenized equations
read $x_i^{d_i}=0$ for every $i$, forcing all projective coordinates to
vanish, which is impossible.
\end{proof}

An infinite surcomplex coordinate is still an affine coordinate in this
statement; it is not the projective point at infinity, and the last sentence
does not assert that all coordinates have nonnegative valuation.

\subsection{Finite support propagation without positivity}
Let $A$ be a Hahn coefficient ring, $S$ contain the coefficient support of
$H$, and $T$ be the union of the coefficient supports of the $E_i$. Neither
support need be positive. With $L=\deg H\ge0$,
\begin{equation}\label{polynomial:eq:finitepolysupport}
 \supp\NF_F(H)\ \cup\ \bigcup_i\supp Q_i
 \ \subseteq\ S+\bigcup_{k=0}^LkT,\qquad 0T=\{0\},
\end{equation}
the empty polynomial having zero outputs. Indeed, following one monomial
through the reduction, each use of an error adds one input support letter and
decreases total degree by at least one, so at most $L$ such uses occur on a
branch; the witness polynomials $Q_i$ accumulate only the preceding
multipliers and obey the same bound. Finite unions and finite Minkowski sums
of well-ordered sets are well ordered, and the coefficientwise products have
finite multiplicity by \cref{polynomial:lem:support}.

This finite-degree mechanism differs from the infinite positive-monoid inverse
$S^*$ of \cref{polynomial:thm:hensel}. It permits \emph{negative} valuations
without sacrificing support control, but only because total degree rules out
an infinite reduction chain. Neither mechanism subsumes the other: analytic
perturbations need not have bounded degree, while polynomial perturbations
need not have positive support.

\section{Multivariate B\'ezoutians and the Jacobian trace formula}
\label{polynomial:sec:bezoutian}
Retain \eqref{polynomial:eq:polysystem} over the arbitrary commutative ring
$A$ and define
\begin{equation}\label{polynomial:eq:multilambda}
 \lambda_F(H)=[x^{\mathrm{top}}]\,\NF_F(H).
\end{equation}
B\'ezoutian and residue duality are classical subjects
\cite{polynomial:BCRS,polynomial:BMP}; what is proved here is an explicit
statement for this family over arbitrary coefficient rings, with the
comparison to the supplied Hahn-coherent framework.

\begin{theorem}[A coefficient-independent perfect pairing]
\label{polynomial:thm:multiperfect}
The pairing $(G,H)\mapsto\lambda_F(GH)$ on $B$ is perfect over $A$. Its Gram
determinant in the rectangular basis is the sign of the permutation
$\alpha\mapsto\boldsymbol d-\boldsymbol1-\alpha$, hence a constant $+1$ or
$-1$ independent of every coefficient of the errors. For $\deg H<\rho_\sharp$
one has $\lambda_F(H)=0$, whereas $\lambda_F(x^{\mathrm{top}})=1$.
\end{theorem}
\begin{proof}
Reduction cannot increase degree, so $\lambda_F(x^\alpha x^\beta)=0$ when
$|\alpha|+|\beta|<\rho_\sharp$. If the degrees sum to $\rho_\sharp$, reduction
of any nonstandard monomial strictly lowers degree, so
\begin{equation}\label{polynomial:eq:gramdegree}
 \lambda_F(x^{\alpha+\beta})=
 \begin{cases}1,&\alpha+\beta=\boldsymbol d-\boldsymbol1,\\
 0,&\text{otherwise},\end{cases}
 \qquad |\alpha|+|\beta|=\rho_\sharp .
\end{equation}
Order the row indices $\alpha$ by nondecreasing total degree and the columns
by the complementary indices $\boldsymbol d-\boldsymbol1-\alpha$. The Gram
matrix is then block lower triangular with identity diagonal blocks: degrees
below the row's complementary level give zero, equal degrees give
\eqref{polynomial:eq:gramdegree}. Its determinant in this reordering is one;
undoing the complement permutation gives the stated determinant and therefore
perfection over $A$. The last two assertions are the same degree argument.
\end{proof}

\subsection{The divided-difference determinant}
For two coordinate tuples $X,Y$ define
\begin{equation}\label{polynomial:eq:divideddiffmatrix}
 D_{ij}(X,Y)=\frac{F_i(Y_1,\ldots,Y_{j-1},X_j,\ldots,X_n)
 -F_i(Y_1,\ldots,Y_j,X_{j+1},\ldots,X_n)}{X_j-Y_j},
\end{equation}
each numerator being divisible by its displayed denominator as a polynomial
identity over $A$, and put
\begin{equation}\label{polynomial:eq:multibezout}
 \Delta_F(X,Y)=\det\bigl(D_{ij}(X,Y)\bigr),\qquad
 J_F(x)=\det\left(\frac{\partial F_i}{\partial x_j}(x)\right),
\end{equation}
viewing $\Delta_F$ in $B\otimes_AB=A[X,Y]/(F(X),F(Y))$.

\begin{theorem}[B\'ezoutian kernel and Jacobian duality]
\label{polynomial:thm:jacobian}
Let $(b_\alpha)=(x^\alpha)_{\alpha\in\mathcal B}$ and let $(b^\alpha)$ be its
residue-dual basis for $\lambda_F$. Then
\begin{align}
 \Delta_F(X,Y)&=\sum_{\alpha\in\mathcal B}b_\alpha(X)b^\alpha(Y)
 \quad\text{in }B\otimes_AB,\label{polynomial:eq:casimir}\\
 J_F&=\sum_{\alpha\in\mathcal B}b_\alpha b^\alpha\quad\text{in }B,
 \label{polynomial:eq:jacobianeuler}\\
 \Tr_B(m_H)&=\lambda_F(J_FH)\qquad(H\in B).
 \label{polynomial:eq:jacobiantrace}
\end{align}
These are identities over $A$, including nonreduced coefficient rings and
parameters at which roots collide.
\end{theorem}
\begin{proof}
The columnwise telescoping identity is $D(X,Y)(X-Y)=F(X)-F(Y)$; multiplying by
the adjugate of $D$ and passing to $B\otimes_AB$ gives
\begin{equation}\label{polynomial:eq:diagonalannihilator}
 (X_j-Y_j)\Delta_F=0\qquad(1\le j\le n),
\end{equation}
with no matrix inverse or determinant denominator anywhere in the argument.

The total degree of row $i$ of $D$ is at most $d_i-1$, and the
highest-degree contribution of that row is diagonal, coming from
$x_i^{d_i}$. So $\Delta_F$ has total degree at most $\rho_\sharp$, with
homogeneous part of degree $\rho_\sharp$ equal to
\begin{equation}\label{polynomial:eq:bezouttop}
 \prod_{i=1}^n\left(\sum_{k=0}^{d_i-1}X_i^{d_i-1-k}Y_i^k\right).
\end{equation}
Apply $\lambda_F$ in the $X$ variables: terms of $X$-degree below
$\rho_\sharp$ vanish by \cref{polynomial:thm:multiperfect}, and a nonstandard
$X$-monomial of degree $\rho_\sharp$ reduces to smaller degree and also
vanishes. The remaining monomial is $X^{\mathrm{top}}$, whose coefficient in
\eqref{polynomial:eq:bezouttop} is one with no $Y$ dependence, so
\begin{equation}\label{polynomial:eq:bezoutnormalization}
 (\lambda_F\otimes\id)(\Delta_F)=1 .
\end{equation}

By \cref{polynomial:thm:multiperfect} the kernel map
\eqref{polynomial:eq:kernelmap}, with $B$ and $\lambda_F$ in place of $A_P$
and $\lambda_P$, is an isomorphism $B\otimes_AB\cong\operatorname{End}_A(B)$.
Equation \eqref{polynomial:eq:diagonalannihilator} says the operator attached
to $\Delta_F$ commutes with multiplication by each coordinate, hence, the
coordinates generating $B$ as an $A$-algebra, that it is $B$-linear;
\eqref{polynomial:eq:bezoutnormalization} says it sends $1$ to $1$. It is
therefore the identity, whose unique representing tensor is the right side of
\eqref{polynomial:eq:casimir}. The multiplication map $B\otimes_AB\to B$
substitutes $X=Y=x$, turning each divided difference into the corresponding
partial derivative, so the diagonal of $\Delta_F$ is $J_F$, which is
\eqref{polynomial:eq:jacobianeuler}. Finally the coefficient of $b_\alpha$ in
$Hb_\alpha$ is $\lambda_F(Hb_\alpha b^\alpha)$; summing these diagonal matrix
coefficients gives \eqref{polynomial:eq:jacobiantrace}.
\end{proof}

\begin{corollary}[Simple-root weights and Euler--Jacobi vanishing]
\label{polynomial:cor:eulerjacobi}
Over an algebraically closed Hahn field $K$, suppose every solution of
\eqref{polynomial:eq:polysystem} is simple. Then
\begin{equation}\label{polynomial:eq:multisimpleresidue}
 \lambda_F(H)=\sum_{\xi\in Z}\frac{H(\xi)}{J_F(\xi)},
\end{equation}
and consequently
\begin{equation}\label{polynomial:eq:eulerjacobi}
 \sum_{\xi\in Z}\frac{H(\xi)}{J_F(\xi)}=0\quad\text{whenever }
 \deg H<\rho_\sharp,
 \qquad
 \sum_{\xi\in Z}\frac{\xi^{\mathrm{top}}}{J_F(\xi)}=1 .
\end{equation}
For arbitrary multiplicities, with $e_\xi$ the local-factor idempotent of
\eqref{polynomial:eq:finitealgebraproduct}, the functional
$\lambda_{F,\xi}(H):=\lambda_F(e_\xi H)$ is a perfect local residue functional
and
\begin{equation}\label{polynomial:eq:localjacobianlength}
 \lambda_{F,\xi}(J_FH)=(\dim_KB_\xi)\,H(\xi).
\end{equation}
\end{corollary}
\begin{proof}
Different local factors are orthogonal for the residue pairing, their products
being zero, so its restriction to each factor is perfect. For a simple factor
$Ke_\xi$ write $\lambda_F(e_\xi H)=c_\xi H(\xi)$; the trace identity
\eqref{polynomial:eq:jacobiantrace} applied to $e_\xi$ gives $1=c_\xi
J_F(\xi)$, which is \eqref{polynomial:eq:multisimpleresidue}. The degree
bounds of \cref{polynomial:thm:multiperfect} then give
\eqref{polynomial:eq:eulerjacobi}. In a factor of length $\mu$ multiplication
by $H$ has trace $\mu H(\xi)$, and the global identity applied to $e_\xi H$
gives \eqref{polynomial:eq:localjacobianlength}.
\end{proof}

The simple-root formula is not used to define the functional. Defining it that
way would obscure the multiple-root case and would falsely suggest that
division by a discriminant is necessary. The top-normal-form definition
\eqref{polynomial:eq:multilambda}, the perfect Gram matrix, and the
B\'ezoutian all remain defined at every collision --- exactly as in one
variable, where \cref{polynomial:thm:residuepairing} has no discriminant
denominator either.

\subsection{Parameters on one fixed ordinary domain, and the analytic
comparison}
\begin{theorem}[Global polynomial families on one ordinary parameter domain]
\label{polynomial:thm:polyparameters}
Let $U\subseteq\C^s$ be one fixed ordinary domain and suppose the finitely
many coefficients in \eqref{polynomial:eq:polysystem} lie in $A=\HH(U)$ or in
$A=\HHn(U)$. Then the global polynomial quotient is finite free of rank $D$
over $A$; normal forms, multiplication matrices, the residue Gram matrix
\emph{and its inverse}, and the B\'ezoutian trace identity all have
coefficients holomorphic on the same $U$. They commute with every allowed
finite coherent parameter substitution and with specialization to a finite
surcomplex parameter. No positivity of the errors is required for the version
over $\HH(U)$, and no ordinary neighbourhood is shrunk.
\end{theorem}
\begin{proof}
Apply
\cref{polynomial:thm:globalfree,polynomial:thm:multiperfect,polynomial:thm:jacobian}
to the indicated coefficient ring. Normal-form computations are finite, obey
\eqref{polynomial:eq:finitepolysupport}, and only add and multiply functions
already holomorphic on $U$. The inverse Gram matrix is its adjugate divided by
the constant determinant $\pm1$, so it too involves only finite polynomial
operations, and all other constructions are finite determinants or matrix
operations. Since evaluation and allowed coherent parameter substitutions are
ring homomorphisms, the base-change assertion of
\cref{polynomial:thm:globalfree} gives compatibility.
\end{proof}

Now suppose in addition that every coefficient of every $E_i$ has
\emph{positive} Hahn support, so that the parameter ring is integral. On an
ordinary product of circles in the fiber variables, the supplied research
manuscript defines its residue by the strong coefficient formula
\begin{equation}\label{polynomial:eq:sourceLambda}
 \Lambda_F(H)=\sum_{k\in\N^n}(-1)^{|k|}
 \left[x_1^{d_1(k_1+1)-1}\cdots x_n^{d_n(k_n+1)-1}\right]
 \Bigl(H\prod_iE_i^{k_i}\Bigr),
\end{equation}
obtained by expanding $1/F_i$ about $x_i^{d_i}$ and applying the ordinary
Cauchy coefficient formula at each Hahn exponent; positive supports make this
family strongly summable by \cref{polynomial:lem:support}. Its values are
coherent parameter functions and $\Lambda_F(F_iH)=0$, since after cancellation
of $F_i$ the coefficientwise integrand is holomorphic in $x_i$ on its disk.
These are the functional and annihilation property proved in
\cite{polynomial:SourceResearch}.

\begin{theorem}[Polynomial identification of the Jacobian trace element]
\label{polynomial:thm:analyticcomparison}
Under the lower-total-degree hypothesis of
\eqref{polynomial:eq:polysystem} \emph{and} the positive-support hypothesis
just stated ---\,positivity being required of the errors $E_i$, not of the
ambient ring\,--- over the fixed-domain ring $\HHn(U)$,
\begin{equation}\label{polynomial:eq:residuecomparison}
 \Lambda_F(H)=\lambda_F(H)\qquad\text{for every polynomial }H .
\end{equation}
Consequently the trace element of the supplied residue pairing, restricted to
this polynomial family, is exactly the class of $J_F$, so that
$\Tr(m_H)=\Lambda_F(J_FH)$. All solutions in a specialized surcomplex fiber
are infinitesimal, so the global polynomial count $D$ agrees with the
finite-monad analytic count in this subclass.
\end{theorem}
\begin{proof}
Evaluate \eqref{polynomial:eq:sourceLambda} on a standard monomial
$H=x^\alpha$, of degree at most $\rho_\sharp$. For $k\ne0$ the extraction
degree is $\rho_\sharp+\sum_id_ik_i$, whereas the numerator has degree at most
$|\alpha|+\sum_i(d_i-1)k_i\le\rho_\sharp+\sum_id_ik_i-|k|$, strictly smaller.
So every $k\ne0$ term vanishes, and the $k=0$ term is one precisely for the
top monomial and zero for the other basis monomials. Since both functionals
annihilate the defining ideal, normal forms extend the equality to every
polynomial; their Gram matrices in the shared rectangular basis coincide, so
the element representing the trace under the perfect pairing is $J_F$ by
\cref{polynomial:thm:jacobian}.

For the location of the roots, specialize the parameters in an algebraically
closed Hahn workspace. The multiplication matrices of the $x_i$ have integral
coefficients, and reducing their entries modulo $\mm_K$ gives the nilpotent
coordinate shifts in $K[x]/(x_1^{d_1},\ldots,x_n^{d_n})$ reduced over $\C$.
Each characteristic polynomial therefore has the form $T^D+\sum_{j<D}c_jT^j$
with every $c_j\in\mm_K$. A noninfinitesimal root would have finite inverse,
and dividing by its $D$th power would express zero as one plus a finite sum of
infinitesimals, which is impossible. All eigenvalues, hence all solution
coordinates, are infinitesimal. The finite polynomial quotient and the
analytic quotient have the same rectangular basis and the same defining
multiplication relations, so the inclusion of polynomials induces an
isomorphism between them.
\end{proof}

This settles a precisely delimited comparison left unproved in the supplied
research manuscript: the Jacobian identity for its residue functional in the
lower-degree polynomial family, over the fixed-domain ring. It does
\emph{not} settle the comparison for arbitrary holomorphic errors, nor for
polynomials with higher-degree positive errors that can create roots outside
the finite monad --- \cref{polynomial:ex:escape} is a warning against exactly
that extrapolation. The general classical theory of complete-intersection
residues already contains Jacobian identities; the contribution here is a
direct normalization and support-compatible identification with the particular
supplied functional, not a claim that the classical identity was previously
unknown.

\begin{example}[A six-dimensional coupled algebra with an explicit dual basis]
\label{polynomial:ex:sixpoint}
Let $a,b,c,d$ be arbitrary surcomplex coefficients, or indeterminates in a
commutative coefficient ring, and take
\begin{equation}\label{polynomial:eq:sixsystem}
 F_1=x^2-ay-b,\qquad F_2=y^3-cx-d,
\end{equation}
so $(d_1,d_2)=(2,3)$ and the degree hypotheses hold. The quotient has rank six
with basis $(1,y,y^2,x,xy,xy^2)$ and $\lambda_F(H)=[xy^2]\NF_F(H)$; no nonzero
condition on $a,b,c,d$ is required. For a second tuple $(u,w)$,
\[
 D(X,Y)=\begin{pmatrix}x+u&-a\\-c&y^2+yw+w^2\end{pmatrix},
 \qquad \Delta_F=(x+u)(y^2+yw+w^2)-ac,
\]
whose coefficients in the first basis give the residue-dual basis
\begin{equation}\label{polynomial:eq:sixdual}
 (xy^2-ac,\ xy,\ x,\ y^2,\ y,\ 1),
\end{equation}
with diagonal $J_F=6xy^2-ac$ and residue Gram determinant $-1$ for every
coefficient choice. Thus $\Tr(m_H)=\lambda_F\bigl((6xy^2-ac)H\bigr)$,
including at all collisions. When $c\ne0$ ordinary elimination gives
$y^6-2dy^3-ac^2y+d^2-bc^2=0$ and $x=(y^3-d)/c$; at $c=0$ it is the division by
$c$ that fails, not the six-element basis or the residue theory --- a global
basis avoids exceptional denominators introduced by a particular elimination.

For a concrete infinitesimal splitting take $a=c=t>0$ and $b=d=0$. Solving
$x=y^3/t$ gives $y(y^5-t^3)=0$, so the roots are $(0,0)$ and
$(t^{4/5}\zeta^3,t^{3/5}\zeta)$ with $\zeta^5=1$. The Jacobian is $-t^2$ at
the origin and $5t^2$ at each of the other five, all six roots being simple,
and
\begin{equation}\label{polynomial:eq:sixweights}
 \lambda_F(H)=-\frac{H(0,0)}{t^2}
 +\frac1{5t^2}\sum_{\zeta^5=1}H\bigl(t^{4/5}\zeta^3,t^{3/5}\zeta\bigr).
\end{equation}
Each individual weight is infinite, yet $\lambda_F(1)=0$ by cancellation and
$\lambda_F(xy^2)=1$. In the underlying parameter family the value $a=c=b=d=0$
gives the length-six origin with algebra $K[x,y]/(x^2,y^3)$ and the same
nondegenerate residue determinant.
\end{example}

\section{Polynomials versus Hahn-coherent entire data}
\label{polynomial:subsec:uniformdegree}
There is a simple exact description of the intersection of the two languages.
For fixed $\Gamma$,
\begin{equation}\label{polynomial:eq:uniformdegree}
 K_\Gamma[z]=\bigcup_{d\ge0}\bigl(\C[z]_{\le d}((t^\Gamma))\bigr),
\end{equation}
equality of polynomial coefficient data, or of the corresponding
Hahn-coherent evaluations on the finite thickened plane. A finite-degree
polynomial has finitely many Hahn coefficients whose finite union of supports
is well ordered, and grouping by exponent gives ordinary polynomial
coefficient functions with the same uniform degree bound; conversely a Hahn
family of ordinary polynomials of degree at most $d$ regroups into $d+1$ Hahn
coefficients whose supports are subsets of the given common well-ordered
support.

The bound must be uniform in the exponent. The coherent datum
\eqref{polynomial:eq:geometricwarning} has ordinary entire coefficient
functions and well-ordered support, yet no uniform degree bound; on finite
surcomplex points it equals $1/(1-tz)$ by strong geometric summation, and it
is not a polynomial, since its ordinary coefficient functions have unbounded
degrees and coefficient uniqueness excludes a bounded-degree representation.
Its rational continuation has a pole at the infinite point $t^{-1}$. Thus
``entire on the thickened ordinary plane'' is not ``entire at all surreal
scales''; this is consistent with the all-scale polynomial rigidity theorem
stated in \cite{polynomial:SourceAnalysis}, which we take as given and do not
reprove.

For a finite root multiset, the supports of its finitely many symmetric
coefficients have well-ordered union, so the finite CRT, cluster-factor and
residue constructions above never encounter a failure caused merely by an
infinite number of spatial centers. They can still produce negative exponents
through small denominators, and root extraction can enlarge the exponent
group. The supplied global manuscript \cite{polynomial:SourceGlobal} proves a
different phenomenon for infinitely many ordinary centers: its cluster family
$P_m(u)=u-t^{1/m}$ cannot be the exact zero divisor of one Hahn-coherent
entire datum, because the union of the symmetric coefficient supports is
decreasing, whereas its family $P_m(u)=u^m-t$ is realizable, all symmetric
coefficients using a single exponent even though the individual root supports
contain the descending sequence $1/m$. These are results of that supplied
manuscript, not newly established global theorems here. A claim that every
separately specified infinite family can be assembled from the finite results
above would be unjustified.

Finally, three mechanisms have been used and are \emph{not} interchangeable.
\Cref{polynomial:thm:orderrouche} quantifies over an order circle and uses a
fixed-degree first-order statement. \Cref{polynomial:cor:valrouche} compares
finite coefficient valuations and preserves an initial polynomial. The
supplied contour theory integrates ordinary holomorphic coefficients and then
assembles a strongly summable Hahn series. Agreement of some consequences does
not make the mechanisms interchangeable.

\section{Conflicts recorded from the wider set}\label{polynomial:sec:conflicts}
The merge contract governing this report records two disagreements in the
wider manuscript set that touch this cluster. Their mathematics belongs to the
companion report on residues and contours, and is not restated here; what is
recorded is their \emph{status}, because the present report's residue
functionals are the polynomial restrictions of the objects at issue and it
would be easy to read more into them than is proved.

\subsection{Whether an arbitrary isolated complete intersection has a contour
representation of its perturbation residue}
Three manuscripts of the wider set treat this differently, and the
disagreement is about status, not arithmetic: the same statement is proved,
assumed, and declined.

One proves it unconditionally. With $A$ constructed from finite colength alone
so that $AF=(z_1^{e_1},\ldots,z_n^{e_n})$ and $G=AF$, it establishes
$\Res_F(H)=\Res_G(H\det A)$ and then a Hahn integral representation of
$\Res_F(H)$ over the coordinate torus $T_r$ as
$(2\pi i)^{-n}\int_{T_r}H\det A/(G_1\cdots G_n)$, needing only the strong
residue expansion, ordinary finite-family residue transformation and positive
rescaling. Its own two non-claims survive as well: the \emph{untransformed}
form has no such torus integral, and $\det A$ is not an optional Jacobian
factor, since $AF$ can introduce extra zeros.

A second reaches a determinant torus formula too, but its theorem opens by
assuming the finite-algebra and normal-form results of a companion geometry
manuscript and routes through that companion's univariate separating
denominators. Its general-case result is therefore \emph{conditional} on an
unrefereed companion, not independent.

A third explicitly \emph{declines} the same statement, proving the realization
for one variable and for a coordinate-power class only, and arguing that for
arbitrary isolated $f$ a single diagonal rescaling need not make the leading
equations coordinate powers or even keep the leading homogeneous system
isolated; that the usual cycle $\{|f_i|=\epsilon_i\}$ is not a coordinate
product torus; and that only a germwise, coefficient-dependent cycle
interpretation is available. It lists a general realization theorem as an open
problem.

\emph{Resolution as applied here.} We do not average these. The proved
theorem and its proof belong in the companion contours report, printed there
with the third manuscript's scope discussion reproduced immediately
afterwards, stating plainly that one of the three sources judged the general
case unproved and why, and that the second's route is conditional rather than
independent. For the present report the consequence is a scope statement:
\emph{nothing in this article asserts a contour representation of any residue.}
The functional $\lambda_P$ of \eqref{polynomial:eq:lambdadef} and the
functional $\lambda_F$ of \eqref{polynomial:eq:multilambda} are defined by
finite polynomial division, and
\cref{polynomial:prop:localresidues,polynomial:thm:jacobian} are proved from
that definition alone. The one place where a contour functional appears is
\cref{polynomial:thm:analyticcomparison}, which identifies $\Lambda_F$ with
$\lambda_F$ only under both the lower-total-degree hypothesis and the
positive-support hypothesis, over the fixed-domain ring, and which says
explicitly that the general analytic comparison needs another argument.

\subsection{Whether roots-of-unity quadrature recovers the Hahn contour
integral}
Two manuscripts of the wider set leave opposite impressions here, and the
hypothesis is doing all the work.

One treats it as a negative result. Its example takes $f(z)=1/(z-b)$ with
ordinary real $b=1/2$ on the unit circle, computes the $N$-point
roots-of-unity quadrature $Q_N(f)=1/(1-b^N)$, and observes that although the
ordinary complex limit is $1$, one has $v(Q_N(f)-1)=0$ for every $N$: there is
no convergence in the Hahn valuation topology, and none in the full fine
topology either, since the family is not eventually constant. With $b=t^\varrho$
a rank-one valuation can see convergence, but a higher-rank group with an
element exceeding every $N\varrho$ still cannot.

The other devotes a section to the positive side, with a rank-independent
aliasing criterion and a coefficientwise recovery theorem, under the decay
condition that for every $\gamma$ the set of $k$ with $v(A_k)\le\gamma$ is
finite.

\emph{Resolution as applied here.} The two are compatible; the first even
concedes in prose that taking ordinary coefficient limits gives the Hahn
functional, and that this is a legitimate alternative description but not a
limit in the internal valuation topology. The positive theorem is to be stated
\emph{with} its decay hypothesis, and the negative example kept immediately
alongside as the demonstration that the hypothesis cannot be dropped: the
decay condition demands eventual domination of \emph{every} value-group
threshold, and, as its own source notes, cannot be met by a nontrivial
sequence of finite valuations when no countable subset is cofinal in the
positive value group. For the present report the consequence is again a scope
statement, and one this article has enforced from
\cref{polynomial:rem:notalimit} onward: no construction here is justified by a
quadrature, a limit of partial sums, or an argument that a valuation improves
by a fixed positive amount at each iteration.

\section{Scope, provenance, and what is not proved}\label{polynomial:sec:scope}

\subsection{A consolidated status statement}
All three merged manuscripts carry a status disclaimer of the same shape. They
are consolidated here into one statement, which governs the whole article.

\begin{quote}\small
This is a theorem-driven research exposition. It supplies proofs of the stated
analogues and support estimates. It does not claim that these are first
occurrences in the literature, does not certify historical priority, and does
not claim to resolve a named published conjecture. Several central statements
are classical over real-closed fields, valued fields, or arbitrary commutative
rings, and are identified as such. The literature search behind these
attributions was conducted on 21 September 2026 and was targeted rather than
exhaustive. The article has not been independently refereed. The original
manuscript packages were not proof-assistant formalizations; the repository's
separate coverage ledger now records partial Lean results, without asserting
a complete formalization of this article. The symbolic checks recorded in the source archives use exact
rational, polynomial and finitely truncated formal-series inputs to validate
\emph{displayed finite examples and identities}; they are not machine
verification of the general proofs, and in particular they certify nothing
about arbitrary Hahn supports, arbitrary-rank exponent groups, transfinite
summability, class-set foundations, or real-closed-field quantifier
elimination. The supplied predecessor manuscripts
\cite{polynomial:SourceAnalysis,polynomial:SourceResearch,polynomial:SourceGlobal}
are treated as unpublished source material, credited for their Hahn-coherent
framework, their analytic antecedents and their reused examples, and are not
treated as independently verified publications. No theorem is declared new
merely because the word ``surcomplex'' is placed in its title.
\end{quote}

\subsection{Per-theorem non-claims}
The per-theorem non-claims of the three sources are specific, are \emph{not}
pooled into the statement above, and are preserved where they were made. They
are listed here for reference.

\begin{itemize}
\item \textbf{Classical inputs.} Euclidean division, finite Chinese
remainders, resultants, the Nullstellensatz and residue duality are classical
algebra. Gauss--Lucas is classical complex polynomial geometry, and the proof
of \cref{polynomial:thm:gausslucas} works over any real closed field.
Real-closed-field transfer and non-Archimedean root geometry have extensive
precedents
\cite{polynomial:TarskiSurvey,polynomial:Eisermann,polynomial:Faber,polynomial:CohenMahboubi}.
It would be misleading to label their surcomplex instances wholly new.
\Cref{polynomial:thm:fta} is an application of established algebraic
closedness, not a new proof of the surreal normal-form theorem.
\item \textbf{Support calculus.} \Cref{polynomial:lem:support} does
\emph{not} say that the relevant partial sums converge in the fine surreal
topology, and does \emph{not} say that $k\rho$ eventually exceeds every
positive value-group element; the latter is false in higher rank. Neither
false cofinality claim is ever substituted for the lemma. See
\cref{polynomial:rem:notalimit,polynomial:rem:belowbound}.
\item \textbf{No proper-class objects.} No assertion treats a
proper-class-indexed sum, an infinite-degree object as a polynomial, or a
strong Hahn sum as the limit of its partial sums in the full fine topology.
The formal expression \eqref{polynomial:eq:geometricwarning} is not a global
surcomplex polynomial; all global polynomial evaluations here are finite sums.
See \cref{polynomial:rem:degree,polynomial:rem:classsafety}.
\item \textbf{Gauss--Lucas.} The proof of
\cref{polynomial:thm:gausslucas} is finite and algebraic; it does not call on
compactness of a convex hull or a topological separation theorem. Neither its
conclusion nor that of \cref{polynomial:prop:weighted} describes how critical
points are distributed among infinitesimal clusters.
\item \textbf{Jensen disks.} \Cref{polynomial:thm:jensen} concerns
$F$-valued disks; it is not a statement about ordinary pictures obtained by
discarding infinitesimals.
\item \textbf{Schoenberg.} \Cref{polynomial:thm:schoenberg} is the
surcomplex version of a classical inequality, not a newly discovered general
principle; its matrix proof does not assert completeness of an
infinite-dimensional surcomplex Hilbert space.
\item \textbf{Transfer.} \Cref{polynomial:subsec:transferscope} states a
theorem \emph{scheme}, one sentence for each fixed finite degree bound. It
does not transfer unrestricted quantification over functions, sequences,
arbitrary subsets, or the standard natural numbers as a distinguished
definable set, and it does not make a proper class into a first-order
set-sized model. \Cref{polynomial:thm:orderrouche} is a transfer theorem, not
a new analytic proof; it does not define an integral along a fine-continuous
real-parameter surreal circle, and it is not the same statement as the
leading-coefficient Rouch\'e theorem on ordinary thickened domains in
\cite{polynomial:SourceAnalysis}, whose hypothesis is different. No
real-parameter fine-continuous circle or contour integral has been introduced
anywhere in this article.
\item \textbf{Sampling.} \Cref{polynomial:ex:circlewarning}: the exact
quantifier in \eqref{polynomial:eq:orderrouchehyp} cannot be replaced by
sampling ordinary points.
\item \textbf{Two Rouch\'e theorems.} The modulus and valuation versions have
different hypotheses and different balls. Neither is identified with the
contour theorem on an ordinary thickening in
\cite{polynomial:SourceAnalysis}.
\item \textbf{Hermite signatures.} \Cref{polynomial:thm:hermite} is a
classical signature method, related to the residue and B\'ezoutian forms of
\cite{polynomial:BCRS}; it is \emph{not} the nondegenerate residue pairing of
\cref{polynomial:sec:residue}.
\item \textbf{Newton polygons.} In
\cref{polynomial:thm:newton} and \cref{polynomial:cor:valrouche}, sufficiency
of the tropical double-minimum condition is a consequence of algebraic
closedness together with the exact count, \emph{not} an additional convergence
assertion about a Newton algorithm.
\item \textbf{Asymptotic notation.} In
\cref{polynomial:ex:newtoncubic,polynomial:ex:hensel}, $O(t^k)$ means an
integer-power formal Hahn series supported in exponents at least $k$, not a
fine-topological limiting estimate and not a numerical asymptotic assertion
obtained by varying $t$.
\item \textbf{Nondiscrete supports.} In
\cref{polynomial:ex:nondiscrete}, the binomial identity does not require that
the coefficient series converge after replacing $t$ by a positive real number
$q\in(0,1)$; in fact $q^{\gamma_k}\to q$, not $0$, so the summands would fail
to tend to zero. No such substitution is used.
\item \textbf{Differentiation.} The hypothesis of
\cref{polynomial:lem:initialderivative} cannot simply be removed: a constant
initial polynomial has zero derivative and does not determine the next nonzero
scale of $P'$. The argument can fail in positive residue characteristic, where
a nonconstant reduction may have zero derivative, and
\cref{polynomial:rem:rollehyp} records that a root-free ball is a genuinely
different case.
\item \textbf{Critical allocation.} \Cref{polynomial:cor:tree} is an
allocation theorem, not a claim that pairwise root valuations alone determine
every critical point; a residual critical polynomial with a multiple root may
require its own rescaling analysis. \Cref{polynomial:cor:nearest} is not a
proof of a Euclidean root--critical-point distance conjecture.
\item \textbf{Stability.} \Cref{polynomial:thm:stability} does not assert
equality of the roots, of their ordinary moduli, or of every Hahn coefficient
of the pairwise differences; \cref{polynomial:cor:treestability} preserves the
allocated critical directions, not the individual multiplicities of all
critical points. \Cref{polynomial:ex:sharpdisc} proves sharpness of the strict
boundary \emph{in that family}; it does not claim that $v(\Disc P)$ is the
optimal threshold for every polynomial, and
\cref{polynomial:ex:quartic} shows it is not.
\item \textbf{Residues.} \Cref{polynomial:thm:residuepairing} is the
polynomial core of the one-variable residue duality in
\cite{polynomial:SourceResearch}; it is not claimed as a new general duality
theorem. \Cref{polynomial:thm:trace} is a statement about the monogenic
algebra $R[z]/(P)$ and is not a proof of the general multivariable Jacobian
comparison. For coherent analytic numerators of unbounded Taylor degree one
must first use the fixed-domain division theorems of the supplied manuscripts;
the polynomial argument of
\cref{polynomial:subsec:coherentfamilies} does not replace that step.
\item \textbf{Several variables.} \Cref{polynomial:thm:analyticcomparison}
does not settle the comparison for arbitrary holomorphic errors, nor for
polynomials with higher-degree positive errors that can create roots outside
the finite monad; \cref{polynomial:ex:escape} is the warning against that
extrapolation. The degree hypothesis in
\eqref{polynomial:eq:polysystem} is a genuine restriction, not a disguised
assertion about all isolated complete intersections.
\item \textbf{Effectivity.} A finite algebraic operation on arbitrary Hahn
inputs is well defined but not automatically computationally effective. A
support may have no supplied computable enumeration, and equality of arbitrary
surreal coefficients may have no finite decision procedure in the chosen
representation. Noetherianity of the analytic coefficient rings of the
supplied manuscripts is not asserted; see
\cref{polynomial:rem:nonnoetherian}.
\item \textbf{Contours.} Residues here are either finite formal Laurent
coefficients or, in \cref{polynomial:thm:analyticcomparison} alone, the
specifically identified coherent contour functional; those interpretations are
never conflated with integration along fine-continuous real-parameter loops.
See also \cref{polynomial:sec:conflicts}.
\end{itemize}

\subsection{Hypotheses that cannot be discarded}
Degree is always finite; without that assumption neither the finite
factorization nor the finite root partitions and determinant proofs are
statements of the same kind. Divisibility of the exponent group is used
wherever roots, midpoints of scales or slopes are needed, and no claim is made
that $K_\Z$ is algebraically closed. Positive support in
\cref{polynomial:thm:hensel} is a \emph{common well-ordered support}
condition, not merely a pointwise assertion about an unrelated family of
errors; coprimality there is essential for the support conclusion. The exact
critical count requires a ball occupied by at least one root of the polynomial
being counted, and the mapping-ball form uses $P-P(a)$ to ensure that
occupation. All derivative arguments use residue characteristic zero, through
\eqref{polynomial:eq:integersunit}. Integral coefficients in the perturbation
estimates ensure that translating by a finite root preserves the coefficient
error bound, and
\cref{polynomial:cor:normalization,polynomial:subsec:rescaled} record the
change of coordinates that makes this harmless. The ordered Gauss--Lucas
theorem uses a real closed ordered real part and its conjugation; algebraic
closedness alone does not define its convex hull. The global
complete-intersection theorems use strict lower total degree. No general
closed-form solution by radicals is asserted anywhere, and the presence of
surreal coefficients does not make ordinary higher-degree Galois obstructions
disappear in a prescribed smaller coefficient field.

\subsection{A consolidated theorem map}
\begin{longtable}{@{}P{0.24\textwidth}P{0.31\textwidth}P{0.36\textwidth}@{}}
\toprule
Classical principle & Surcomplex conclusion & Mechanism, ring, and status\\
\midrule
\endhead
Fundamental theorem of algebra & Finite factorization in a set-sized Hahn
workspace & Established Hahn algebraic closedness, then finite induction, over
$K$; \cref{polynomial:thm:fta}.\\[0.4em]
Gauss--Lucas, Jensen, weighted incomplete polynomials & Convex-hull,
conjugate-disk and exact-hull statements with surreal weights & Finite
reciprocal identities over $F$ and $K$;
\cref{polynomial:thm:gausslucas,polynomial:thm:jensen,polynomial:prop:weighted}.\\[0.4em]
Schoenberg inequality & Sharp second-moment bound with collinearity equality
case & Finite matrix compression over a real closed field;
\cref{polynomial:thm:schoenberg}. Classical inequality.\\[0.4em]
Polynomial Rouch\'e; Hermite signatures & Order-circle counts at arbitrary
surreal radii; real-root counts as signatures & Precisely encoded
real-closed-field transfer, and finite trace forms;
\cref{polynomial:thm:orderrouche,polynomial:thm:hermite}.\\[0.4em]
Pseudozero sets & \emph{Exact} root loci under bounded coefficient error, in
both geometries & Explicit phase construction and a single-coefficient
witness, over $\SC$ and $K$; \cref{polynomial:thm:pseudozeros}.\\[0.4em]
Newton polygon & Counts in every ball, shell and residue direction, at
arbitrary group rank & One finite initial-polynomial product over $K$;
\cref{polynomial:thm:initialroots,polynomial:thm:newton}.\\[0.4em]
Non-Archimedean disk mapping & Exact image balls and two mapping degrees &
The same product theorem applied to $P-y$, over $K$;
\cref{polynomial:thm:imageballs}.\\[0.4em]
Finite flat maps and Riemann--Hurwitz & Constant fiber length, branch-value
polynomial, ramification sum $2n-2$ & Monic division and resultants over $K$;
\cref{polynomial:thm:finitemap,polynomial:thm:branchvalues,polynomial:prop:ramification}.\\[0.4em]
Rolle at a scale & Exactly $k-r$ roots of $P^{(r)}$ in an occupied $k$-root
ball ($1\le r\le k$); branch directions; nearest-neighbour identity &
Characteristic-zero differentiation of initial polynomials, over $K$;
\cref{polynomial:thm:criticalballs,polynomial:thm:branch,polynomial:cor:nearest}.\\[0.4em]
Hensel factorization & Unique coprime factors with support in $S^*$; parameter
version with no domain shrinking & Finite CRT linearization and well-founded
recursion, over $\OO_K$ and over $\HHn(U)$ for one fixed $U$;
\cref{polynomial:thm:hensel,polynomial:thm:parameterhensel}.\\[0.4em]
Continuity of root multisets & Optimal $\varepsilon/n$ matching, with
multiplicity, no separation hypothesis & Finite balanced valuation-ball
partition over $\OO_K$; \cref{polynomial:thm:holder}.\\[0.4em]
Simple-root perturbation & Rootwise threshold $T(P)$; full distance matrix and
leading coefficients preserved; exact displacement; second-order Newton error
& Product expansion plus valuation Rouch\'e over $\OO_K$;
\cref{polynomial:thm:stability,polynomial:cor:discprecision,polynomial:cor:treestability}.\\[0.4em]
Residue and trace duality & Perfect residue pairing through collisions over an
arbitrary commutative ring; trace determinant is the discriminant & Monic
division, explicit dual basis, finite CRT, formal Laurent coefficients;
\cref{polynomial:thm:residuepairing,polynomial:thm:trace,polynomial:thm:discriminant}.
Overlap with the supplied results is credited.\\[0.4em]
Nullstellensatz and local multiplicity & Finite-system solvability with a
class-safe certificate; finite local-algebra lengths & Classical algebra in
each set-sized workspace $K[x]$;
\cref{polynomial:thm:nullstellensatz,polynomial:prop:finitealgebra}.\\[0.4em]
B\'ezout, B\'ezoutian, Euler--Jacobi & Finite free rank $\prod_id_i$ over an
arbitrary commutative ring; coefficient-independent perfect pairing; Jacobian
trace formula through collisions & Degree-decreasing rewriting with a
critical-pair argument, divided-difference determinant, kernel map;
\cref{polynomial:thm:globalfree,polynomial:thm:multiperfect,polynomial:thm:jacobian,polynomial:cor:eulerjacobi}.\\[0.4em]
Comparison with a coherent contour residue & $\Lambda_F=\lambda_F$ in a
delimited polynomial subclass & Degree bookkeeping plus errors of positive
support, over $\HHn(U)$ for one fixed $U$; \cref{polynomial:thm:analyticcomparison}.
Conditional on both hypotheses; see \cref{polynomial:sec:conflicts}.\\
\bottomrule
\end{longtable}

\subsection{Directions justified by the proved results}
An effective representation of root-cluster refinements for restricted
computable classes of surreal coefficients, with certified precision
propagation for resultants and multiple critical clusters, would extend
\cref{polynomial:sec:stability}. Comparing support-controlled parameter
factors with ramification and monodromy across the ordinary discriminant would
extend \cref{polynomial:thm:parameterhensel}. Replacing the coordinate-power
highest terms of \eqref{polynomial:eq:polysystem} by a homogeneous regular
sequence with a known finite basis, using a graded division and duality
normalization in place of the rectangular argument, would extend
\cref{polynomial:sec:globalci}; the existing general theory of complete
intersections and B\'ezoutians is relevant, but a fixed-basis arbitrary-ring
statement must specify its hypotheses rather than assume freeness. A
perturbation theory for repeated-root trees would have to preserve a nested
collection of cluster algebras with explicit precision thresholds, since a
theorem preserving individual labels is impossible through genuine collisions.
The general analytic residue comparison for the full class of positive-support
holomorphic perturbations needs another argument. These are directions for
additional work, not conclusions silently assumed in the present proofs.

\appendix
\section{Finite-dependency and support audit}\label{polynomial:app:support}
The only genuinely infinite operations in this article are coprime factor
lifting and the coefficientwise expansions attached to it. All root-product
identities, derivatives, matrices, residue calculations, normal forms and CRT
decompositions involve only finitely many factors or coordinates. For lifting,
an output exponent receives finitely many support-word contributions by
\cref{polynomial:lem:support}; this remains true for a group of arbitrary rank
and for well-ordered supports with infinitely many terms below some upper
bound, which is exactly \cref{polynomial:rem:belowbound}.

\begin{center}
\begin{tabularx}{\textwidth}{@{}LL@{}}
\toprule
Operation & Size or support control\\
\midrule
Finite polynomial addition and multiplication & Finite unions and Minkowski
sums of coefficient supports.\\[0.4em]
Monic division of a finite polynomial & Finitely many ring operations;
integral coefficients remain integral.\\[0.4em]
Initial polynomial at $(a,\rho)$ & One finite translation and dilation, one
minimum, finitely many standard parts.\\[0.4em]
Newton profile and root-cluster tree & Finitely many active indices or root
pair separations; no cofinal sequence of scales.\\[0.4em]
Coprime factor correction & All outputs lie in $S^*$; at each exponent only
earlier correction coefficients occur.\\[0.4em]
Parameter factor correction & The same $S^*$, with one holomorphic finite
matrix inverse on the original fixed domain $U$.\\[0.4em]
H\"older matching & A finite partition of the combined root multiset; no
choice of an infinite limiting matching.\\[0.4em]
Threshold stability estimate & Exact finite product expansion with explicit
valuation remainders.\\[0.4em]
Multivariate normal form & At most $\deg H$ uses of an error on a branch;
support in $S+\bigcup_{k\le L}kT$ by
\eqref{polynomial:eq:finitepolysupport}, with no positivity needed.\\[0.4em]
Residue, trace, B\'ezoutian & Finite quotient matrices or finitely many formal
local Laurent coefficients; no contour limit.\\
\bottomrule
\end{tabularx}
\end{center}

\section{Notation index}\label{polynomial:app:notation}
\begin{longtable}{@{}P{0.29\textwidth}P{0.63\textwidth}@{}}
\toprule
Notation & Meaning\\
\midrule
\endhead
$\SC=\No[i]$ & Surcomplex class field.\\[0.3em]
$t^\gamma=\omega^{-\gamma}$ & Conway normal-form monomial.\\[0.3em]
$\Gamma$, $F=F_\Gamma$, $K=K_\Gamma$ & Set-sized divisible exponent group; the
real and complex Hahn workspaces $\R((t^\Gamma))$, $\C((t^\Gamma))$.\\[0.3em]
$|z|$, $v(z)$, $\lc(z)$ & Surreal-valued modulus; natural valuation; leading
complex coefficient.\\[0.3em]
$\OO_K$, $\mm_K$, $\st$ & Finite elements, infinitesimals, and the standard
part $\OO_K\to\C$.\\[0.3em]
$D(a,r)$, $\overline D(a,r)$ & Open and closed order disks defined by the
modulus.\\[0.3em]
$\Bge a\rho$, $\Bgt a\rho$ & Closed and open valuation balls,
$v(z-a)\ge\rho$ and $v(z-a)>\rho$.\\[0.3em]
$a+t^\rho(c+\mm_K)$ & Residue-direction ball at center $a$, scale $\rho$,
direction $c\in\C$.\\[0.3em]
$w_{a,\rho}(P)$, $I_{a,\rho}(P)$ & Weighted Gauss valuation and the ordinary
complex initial polynomial in the variable $X$.\\[0.3em]
$j_-$, $j_+$ & Smallest and largest indices attaining the weighted coefficient
minimum.\\[0.3em]
$\phi_P$ & Newton valuation profile at center zero.\\[0.3em]
$S^*$ & Additive monoid generated by a positive well-ordered support,
including zero.\\[0.3em]
$d_{ij}$, $d_i$, $\delta_i$, $T(P)$ & $v(\alpha_i-\alpha_j)$;
$v(P'(\alpha_i))=\sum_{j\ne i}d_{ij}$; $\max_{j\ne i}d_{ij}$; the threshold
$\max_i(d_i+\delta_i)$.\\[0.3em]
$E$, $\varepsilon$ & Perturbation and its coefficient precision
$\min_jv([z^j]E)$.\\[0.3em]
$A_P=K[z]/(P)$, $\lambda_P$ & Finite quotient and its top-remainder
coefficient functional $[z^{n-1}]\rem_P$.\\[0.3em]
$m_h$, $\Res(P,Q)$, $\Disc(P)$ & Multiplication endomorphism; resultant;
discriminant.\\[0.3em]
$\Delta_P$, $b^i$ & Univariate B\'ezoutian and the residue-dual basis
\eqref{polynomial:eq:dualbasis}.\\[0.3em]
$\mathcal B$, $D$, $\rho_\sharp$, $x^{\mathrm{top}}$ & Rectangular exponent
set, its size $\prod_id_i$, the top degree $\sum_i(d_i-1)$, and the top
monomial.\\[0.3em]
$\NF_F$, $\lambda_F$, $\Delta_F$, $J_F$ & Multivariate normal form, residue
functional, B\'ezoutian, Jacobian determinant.\\[0.3em]
$\Lambda_F$ & The supplied coherent contour residue
\eqref{polynomial:eq:sourceLambda}, used only in
\cref{polynomial:thm:analyticcomparison}.\\[0.3em]
$\HH(U)$, $\HHn(U)$, $\HHp(U)$ & $\OO(U)((t^\Gamma))$ on \emph{one fixed}
ordinary domain $U$, its nonnegative-support subring, and the positive-support
ideal of that subring.\\
\bottomrule
\end{longtable}

\begin{thebibliography}{99}
\addcontentsline{toc}{section}{References}
\small

\bibitem{polynomial:MsFactorization}
\emph{Surcomplex Polynomial Algebra: Factorization, Root Geometry, Multiscale
Stability, and Residue Duality}. Supplied manuscript dated September 21, 2026.
Merged here; source of
\cref{polynomial:thm:pseudozeros,polynomial:prop:clusterdisc,polynomial:thm:stability,polynomial:cor:discprecision,polynomial:thm:residuepairing,polynomial:thm:trace,polynomial:thm:finitemap,polynomial:thm:branchvalues,polynomial:prop:ramification,polynomial:thm:nullstellensatz}
and of \cref{polynomial:ex:circlewarning,polynomial:ex:quartic}.

\bibitem{polynomial:MsNewton}
\emph{Surcomplex Polynomial Algebra: Order Geometry, Newton Profiles, and
Scale-Resolved Stability}. Supplied manuscript dated September 21, 2026. Merged
here as the base; source of
\cref{polynomial:thm:initialroots,polynomial:thm:newton,polynomial:thm:imageballs,polynomial:thm:criticalballs,polynomial:thm:branch,polynomial:cor:tree,polynomial:thm:holder,polynomial:thm:jensen,polynomial:thm:schoenberg}
and of the second-order estimate
\eqref{polynomial:eq:newtonerror}.

\bibitem{polynomial:MsTrees}
\emph{Surcomplex Polynomial Algebra: Ordered Geometry, Valuative Root Trees,
and Residue Duality}. Supplied manuscript dated September 21, 2026. Merged
here; source of
\cref{polynomial:thm:hermite,polynomial:prop:weighted,polynomial:cor:nearest,polynomial:prop:disctree,polynomial:thm:globalfree,polynomial:thm:multiperfect,polynomial:thm:jacobian,polynomial:cor:eulerjacobi,polynomial:thm:polyparameters,polynomial:thm:analyticcomparison}
and of \cref{polynomial:ex:escape,polynomial:ex:sixpoint}.

\bibitem{polynomial:SourceAnalysis}
\emph{Surcomplex Analysis: Infinitesimal Calculus, Hahn-Coherent Holomorphy,
and Contour Theory}. Supplied predecessor manuscript dated September 21, 2026;
file \texttt{surcomplex\_analysis(3).tex}. In particular its workspaces,
preparation, Rouch\'e and all-scale rigidity sections. No publication or
authorship attribution is inferred from the filename.

\bibitem{polynomial:SourceResearch}
\emph{Hartogs Coherence, Finite Maps, and Residue Duality over the Surcomplex
Numbers}. Supplied predecessor manuscript dated September 21, 2026; file
\texttt{surcomplex\_research.tex}. In particular its positive operators,
parameter division, finite fibers, and one-variable residue and trace
formulas.

\bibitem{polynomial:SourceGlobal}
\emph{Global Divisors and Cohomological Obstructions in Hahn-Coherent
Surcomplex Analysis}. Supplied predecessor manuscript dated September 21, 2026;
file \texttt{surcomplex\_global\_theorems.tex}. In particular its
symmetric-support divisor criterion and support-restricted Chinese remainder
theorem.

\bibitem{polynomial:Conway}
John H. Conway, \emph{On Numbers and Games}, second edition, A K Peters, 2001;
first edition, Academic Press, 1976.

\bibitem{polynomial:Gonshor}
Harry Gonshor, \emph{An Introduction to the Theory of Surreal Numbers}, London
Mathematical Society Lecture Note Series \textbf{110}, Cambridge University
Press, 1986.

\bibitem{polynomial:Neumann}
B. H. Neumann, ``On ordered division rings,'' \emph{Transactions of the
American Mathematical Society} \textbf{66} (1949), 202--252.

\bibitem{polynomial:Poonen}
Bjorn Poonen, ``Maximally complete fields,'' \emph{L'Enseignement
Math\'ematique} (2) \textbf{39} (1993), 87--106. Corollary 4 supplies the
divisible-group algebraic-closure input.

\bibitem{polynomial:TarskiSurvey}
Lou van den Dries, ``Alfred Tarski's elimination theory for real closed
fields,'' \emph{The Journal of Symbolic Logic} \textbf{53} (1988), no.~1,
7--19.

\bibitem{polynomial:Swan}
Richard G. Swan, ``Tarski's principle and the elimination of quantifiers,''
expository notes on real closed fields and transfer.

\bibitem{polynomial:CohenMahboubi}
Cyril Cohen and Assia Mahboubi, ``Formal proofs in real algebraic geometry:
from ordered fields to quantifier elimination,'' \emph{Logical Methods in
Computer Science} \textbf{8} (2012), no.~1:02.

\bibitem{polynomial:Eisermann}
Michael Eisermann, ``The fundamental theorem of algebra made effective: an
elementary real-algebraic proof via Sturm chains,'' \emph{The American
Mathematical Monthly} \textbf{119} (2012), no.~9, 715--752.

\bibitem{polynomial:Eberl}
Manuel Eberl, ``Two theorems about the geometry of the critical points of a
complex polynomial,'' \emph{Archive of Formal Proofs}, November 14, 2023. Used
as a reference for ordinary Gauss--Lucas and Jensen geometry, not as a formal
verification of this article.

\bibitem{polynomial:Pereira}
Rajesh Pereira, ``Differentiators and the geometry of polynomials,''
\emph{Journal of Mathematical Analysis and Applications} \textbf{285} (2003),
no.~1, 336--348.

\bibitem{polynomial:KushelTyaglov}
Olga Kushel and Mikhail Tyaglov, ``Circulants and critical points of
polynomials,'' \emph{Journal of Mathematical Analysis and Applications}
\textbf{439} (2016), no.~2, 634--650. The Schoenberg inequality and the
matrix-differentiator precedents are classical inputs to the literature
comparison.

\bibitem{polynomial:Zhang}
Extensions of the Gauss--Lucas theorem to weighted incomplete polynomials in
the ordinary complex setting; used as the classical antecedent of
\cref{polynomial:prop:weighted}.

\bibitem{polynomial:Benedetto}
Robert L. Benedetto, \emph{Non-Archimedean Dynamics in Dimension One: Lecture
Notes}, Arizona Winter School, 2010. See \S3, especially Corollary 3.11 and
the discussion of valuation polygons, for the complete rank-one analytic
setting.

\bibitem{polynomial:ChoiLee}
Non-Archimedean root location and Newton-polygon estimates over complete
rank-one valued fields; classical antecedent for
\cref{polynomial:thm:newton,polynomial:cor:nearest}.

\bibitem{polynomial:Faber}
Xander Faber, ``Topology and geometry of the Berkovich ramification locus for
rational functions,'' with the associated non-Archimedean Rolle and
ramification results. Our proofs are direct arbitrary-rank polynomial
specializations and use no Berkovich space.

\bibitem{polynomial:FaroukiHan}
Rida T. Farouki and Chang Yong Han,
``Robust plotting of generalized lemniscates,''
\emph{Applied Numerical Mathematics} \textbf{51} (2004), nos.~2--3, 257--272.
\href{https://doi.org/10.1016/j.apnum.2004.05.007}{doi:10.1016/j.apnum.2004.05.007}.
The ordinary pseudozero-set antecedent of \cref{polynomial:thm:pseudozeros}.

\bibitem{polynomial:BCRS}
Eberhard Becker, Jean-Paul Cardinal, Marie-Fran\c{c}oise Roy,
and Zbigniew Szafraniec,
``Multivariate Bezoutians, Kronecker symbol and Eisenbud--Levine formula,''
in \emph{Algorithms in Algebraic Geometry and Applications},
Progress in Mathematics \textbf{143}, Birkh\"auser, 1996, 79--104.
\href{https://doi.org/10.1007/978-3-0348-9104-2_5}{doi:10.1007/978-3-0348-9104-2\_5}.
Standard reference for the multivariate duality of
\cref{polynomial:sec:bezoutian} and the Hermite form of
\cref{polynomial:thm:hermite}.

\bibitem{polynomial:BMP}
Thomas Brazelton, Stephen McKean, and Sabrina Pauli,
``B\'ezoutians and the $\mathbb A^1$-degree,''
\emph{Algebra \& Number Theory} \textbf{17} (2023), 1985--2012.
\href{https://doi.org/10.2140/ant.2023.17.1985}{doi:10.2140/ant.2023.17.1985};
\href{https://arxiv.org/abs/2103.16614}{arXiv:2103.16614}.
Broader context for
\cref{polynomial:thm:multiperfect,polynomial:thm:jacobian}.

\bibitem{polynomial:Stacks}
The Stacks Project Authors, \emph{The Stacks Project}, Tag 00FV: Hilbert
Nullstellensatz. Consulted September 21, 2026.

\bibitem{polynomial:UserReference}
\emph{Surreal number}, subsection ``Surcomplex numbers,'' Wikipedia, the
orientation reference specified in the original request. The proofs use the
primary and supplied sources identified above, not this overview as an
independent algebraic-closure proof.
\end{thebibliography}
\end{document}
